\documentclass[10pt,a4paper,twoside,bg=print,twocolumn,openany,nodeprecatedcode]{dndbook}
\usepackage[utf8]{inputenc}
\usepackage{tabulary}
\usepackage[normalem]{ulem}
\usepackage{color,soul}
\usepackage{float}
\usepackage{caption}
\usepackage{wrapfig}
\usepackage{tocloft}
\usepackage[toc]{multitoc}
\usepackage[hidelinks]{hyperref}
\raggedbottom
\extrafloats{2000}
\cftsetindents{section}{0em}{0em}
\cftsetindents{subsection}{0em}{0em}
\cftsetindents{subsubsection}{0.2em}{0em}
\setcounter{tocdepth}{1}
\renewcommand*{\multicolumntoc}{3}
\definecolor {mytable} {HTML} {f4edd8}
\definecolor {mycomment} {HTML} {ffffff}
\definecolor {mysidebar} {HTML} {d0cfc2}
\definecolor {myreadaloud} {HTML} {d9e8f1}
\definecolor {myreadaloudborder} {HTML} {5a6173}
\definecolor {mypagenum} {HTML} {2a2d2e}
\setlist[description]{nolistsep,listparindent=3pt,topsep=6pt}\title{ \Huge \scshape \nodesto\fontsize{50}{40}\selectfont
 Encyclopedia Exandria \\
 \medskip
 \nodesto
 \medskip\Huge
 A compendium of the Critical Role world}\date{}
\begin{document}
\frontmatter
\maketitle
\tableofcontents
\mainmatter
\chapter{Preface}
\DndDropCapLine{W}{}elcome to \textit{Tales from the Yawning Portal}. Within this book you will find seven of the deadliest dungeons from the history of D\&D, updated for the current edition of the game. Some are classics that have hosted an untold number of adventurers, while others are newer creations boldly staking their place in the pantheon of notable D\&D adventures.


Just as these dungeons have made an impression on D\&D players, so too have tales of their dangers spread across the D\&D multiverse. When the night grows long in Waterdeep, City of Splendors, and the fireplace in the taproom of the Yawning Portal dims to a deep crimson, adventurers from across the Sword Coast-and even some visiting from other D\&D worlds-spin tales and rumors of lost treasures.


\begin{itemize}
\item A wanderer from the distant Shou Empire speaks of strange, leering devil faces carved in dungeon walls that can devour an explorer in an instant, leaving behind not a single trace of the poor soul's passing.
\item A bald, stern wizard clad in blue robes and speaking with a strange accent tells of a wizard who claimed three powerful weapons from a city on the shores of a lake of unknown depths, who spirited them away to a slumbering volcano and dared adventurers to enter his lair and recover them.
\item A one-eyed dwarf spins tales of a castle that fell into the earth, and whose ruins stand above a subterranean grove dominated by a tree that spawns evil.
\end{itemize}

These are only a few of the tales that have spread across the Sword Coast from the furthest reaches of Faerûn and beyond. The minor details change with the telling. The dread tomb of Acererak shifts its location from a dismal swamp, to a searing desert, to some other forbidding clime in each telling. The key elements remain the same in each version of the tales, lending a thread of truth to the tale.

The seeds of those stories now rest in your hand. D\&D's deadliest dungeons are now part of your arsenal of adventures. Enjoy, and remember to keep a few spare character sheets handy.

\section{Using This Book}
Tales from the Yawning Portal contains seven adventures taken from across D\&D's history.

The introduction of each adventure provides ideas on adapting it to a variety of D\&D settings. Use that information to place it in your campaign or to give you an idea of how to adapt it.

These adventures provide the perfect side quest away from your current campaign. If you run published D\&D campaigns, such as \textit{Storm King's Thunder}, the higher level adventures presented here are an ideal way to extend the campaign beyond.

\section{About the Adventures}
\subsection{The Sunless Citadel}
\textit{The Sunless Citadel,} written by Bruce R. Cordell, was the first published adventure for the third edition of the D\&D game. It is designed for a party of four or five 1st level player characters.

Ever since its publication in 2000, The Sunless Citadel has been widely regarded as an excellent way to introduce new players to the game. It's also a great starting experience for someone looking to be a Dungeon Master for the first time.

\subsection{The Forge of Fury}
\textit{The Forge of Fury,} written by Richard Baker, was published in 2000 shortly after The Sunless Citadel. Characters who succeeded in that mission and advanced to 3rd level were now ready to take on the challenges of a ruined dwarven fortress.

Like its predecessor, \textit{The Forge of Fury} is tailored to provide increasingly tougher threats as the characters make their way through the fortress. Those who survive the experience can expect to advance to 5th level-seasoned adventurers ready to strive for greater glory and renown.

\subsection{The Hidden Shrine of Tamoachan}
\textit{The Hidden Shrine of Tamoachan,} written by Harold Johnson and Jeff R. Leason, made its debut under the title \textit{Lost Tamoachan} at the Origins game convention in 1979, where it was used in the official D\&D competition. The first published version of the adventure was produced in 1980.

The updated version of the adventure presented herein is designed for a group of four or five 5th-level player characters.

\subsection{White Plume Mountain}
Lawrence Schick, the author of \textit{White Plume Mountain}, related in the 2013 compilation Dungeons of Dread that he wrote the adventure as a way of persuading Gary Gygax to hire him as a game designer. Not only did he get the job, but \textit{White Plume} became an instant favorite when it was first published in 1979.

The version of the adventure in this book is tailored to a group of characters of 8th level.

\subsection{Dead in Thay}
\textit{Dead in Thay,} written by Scott Fitzgerald Gray, was created when the fifth edition D\&D game was in the testing stages. In its original form, it was used as the story of the D\&D Encounters season in the spring of 2014. Featuring an immense and lethal dungeon known as the Doomvault, the adventure serves as a tribute to \textit{Tomb of Horrors, Ruins of Undermountain}, and other "killer dungeons" throughout the history of the game.

The version of \textit{Dead in Thay} presented here is modified for use in home campaigns. It is designed for characters of 9th to 11th level.

\subsection{Against the Giants}
The three linked adventures that make up \textit{Against the Giants} were created and originally released in 1978, during the time when Gary Gygax was still writing the \textit{Player's Handbook} for the original AD\&D game. Despite being (in a sense) older than the game itself, these adventures continue to hold a special place in the hearts and memories of D\&D players of all ages.

The compilation of \textit{Steading of the Hill Giant Chief, Glacial Rift of the Frost Giant Jarl}, and \textit{Hall of the Fire Giant King} was published in 1981 as \textit{Against the Giants}. The version presented here is designed to be undertaken by characters of 11th level.

\subsection{Tomb of Horrors}
Before there was much of anything else in the world of the D\&D game, there was the \textit{Tomb of Horrors}.

The first version of the adventure was crafted for Gary Gygax's personal campaign in the early 1970s and went on to be featured as the official Dungeons \& Dragons event at the original Origins gaming convention in 1975. The first publication of \textit{Tomb of Horrors}, as a part of the Advanced D\&D game, came in 1978.

As a proving ground for characters and players alike, fabricated by the devious mind of the game's cocreator, Tomb of Horrors has no equal in the annals of D\&D's greatest adventures. Only high-level characters stand a chance of coming back alive, but every player who braves the Tomb will have the experience of a lifetime.

\section{Running the Adventures}
To run each of these adventures, you need the fifth edition \textit{Player's Handbook, Dungeon Master's Guide}, and \textit{Monster Manual}. Before you sit down with your players, read the text of the adventure all the way through and familiarize yourself with the maps as well, perhaps making notes about complex areas or places where the characters are certain to go, so you're well prepared before the action starts.

\begin{DndSidebar}[float=!b]{}
Text that appears in a box like this is meant to be read aloud or paraphrased for the players when their characters first arrive at a location or under a specific circumstance, as described in the text.



\end{DndSidebar}

The \textit{Monster Manual} contains stat blocks for most of the monsters and NPCs found in this book. When a monster's name appears in \textbf{bold} type, that's a visual cue pointing you to the creature's stat block in the Monster Manual. Descriptions and stat blocks for new monsters appear in appendix B. If a stat block is in that appendix, an adventure's text tells you so.

Spells and nonmagical objects or equipment mentioned in the book are described in the \textit{Player's Handbook}. Magic items are described in the \textit{Dungeon Master's Guide}, unless the adventure's text directs you to an item's description in appendix A.

\section{Creating a Campaign}
While these adventures were never meant to be combined into a full campaign-over 30 years separates the newest from the oldest-they have been selected to provide play across a broad range of levels. With a little work, you can run a complete campaign using only this book.

Starting with The Sunless Citadel, guide your players through the adventures in the order that they are presented in this book. Each one provides enough XP that, upon completing the adventure, the characters should be high enough level to advance to the next one.

The Yawning Portal, or some other tavern of your own invention or drawn from another D\&D setting, provides the perfect framing device for the campaign. The characters hear rumors of each dungeon, with just enough information available to lead them to the next adventure. Perhaps a friendly NPC drawn from the upcoming adventure visits the tavern in search of help, or some element of a character's background pushes the group down the proper road. In any case, these dungeons are designed to be easily portable to any campaign setting.

\section{The Yawning Portal}
Amid the bustle of Waterdeep, within the Castle Ward where barristers, nobles, and emissaries battle with word and contract, stands an inn not quite like any other. Before there was a Castle Ward or even what could be recognized as an ancestor of the City of Splendors, there was a dungeon, and in that dungeon begins the tale of the Yawning Portal.

In ages past, the mighty wizard Halaster built his tower at the foot of Mount Waterdeep and delved deep into tunnels first built by dwarves and drow in search of ever greater magical power. Halaster and his apprentices expanded the tunnels they found, worming out new lairs under the surface for reasons of their own. In time, their excavations grew into the vast labyrinth known today as Undermountain, the largest dungeon in all of the Forgotten Realms. Halaster eventually disappeared, as have all his apprentices, but the massive complex he built remains to this day.

For untold years, the secrets of Undermountain remained hidden from the surface world. Everyone who entered its halls failed to return. Its reputation as a death trap grew to the point that criminals in Waterdeep who were sentenced to die were forcibly escorted into the dungeon and left to fend for themselves.

All of that changed with the arrival of two men, a warrior named Durnan and a ne'er-do-well named Mirt. The duo were the first adventurers to return from Undermountain, laden with riches and magic treasures. While Mirt used his wealth to buy a mansion, Durnan had different plans. Durnan retired from adventuring and purchased the land on which sat the deep, broad well that was the only known entrance to the dungeon. Around this well he built a tavern and inn that caters to adventurers and those who seek their services, and he called it the Yawning Portal.

Some of the magic Durnan looted on his successful foray into Undermountain granted him a life span that exceeds even that of an elf. And for decades Durnan left delving into Undermountain to younger folk. Yet one day, something drew him back. Days of waiting for his triumphant return from the dungeon turned to months and then years. For nearly a century, citizens of Waterdeep thought him dead. But one night, a voice called up from the well. Few at first believed it could be Durnan, but folk as long-lived as he vouched it so. The Yawning Portal had passed into the hands of his ancestors, but Durnan returned with enough riches for them to quietly retire. Durnan took his customary place behind the bar, raised a toast to his own safe return, and then began serving customers as if he'd never left.

Adventurers from across Faerûn, and even from elsewhere in the great span of the multiverse, visit the Yawning Portal to exchange knowledge about Undermountain and other dungeons. Most visitors are content to swap stories by the hearth, but sometimes a group driven by greed, ambition, or desperation pays the toll for entry and descends the well. Most don't survive to make the return trip, but enough come back with riches and tales of adventure to tempt other groups into trying their luck.

\subsection{The Green Dragon Inn}
The Yawning Portal is not the only renowned tavern in D\&D lore. In the Free City of Greyhawk stands the Green Dragon Inn, which has been the starting point for some of the most successful expeditions to Castle Greyhawk and beyond. The place is crowded and smoke-filled. Patrons talk in low voices, and anyone attempting to strike up a conversation without making a clear intent to pay can expect a cold reception. Paranoia and suspicion run rampant here, as befits a free city that stands at the nexus between a devil-haunted empire, a vast domain locked in the iron-tight grip of a demigod of evil, and a splintered, bickering host of kingdoms nominally committed to justice and weal. In the battered, weary world of Greyhawk, profit and power take precedence over heroics.

\section{Features of the Yawning Portal}
The Yawning Portal's taproom fills the first floor of the building. The 40-foot-diameter well that provides access to Undermountain dominates the space. The "well" is all that remains of Halaster's tower, and now, devoid of the stairways and floors that formed subterranean levels, it drops as an open shaft for 140 feet. Stirges, spiders, and worse have been known to invade the Yawning Portal from below.

Balconies on the tavern's second and third floors overlook the well, with those floors accessed by way of wooden stairs that rise up from the taproom. Guests sitting at the tables on the balconies have an excellent view of the well and the action below.

\subsection{Entering the Well}
Those who wish to enter Undermountain for adventure (or the daring tourists who just want to "ride the rope") must pay a gold piece to be lowered down. The return trip also costs a piece of gold, sent up in a bucket in advance. Once the initial payment is made, a few stairs takes one to the top of the waisthigh lip of the well. The rope that hangs in the center of the well is levered over to the lip by a beam in the rafters, and when those who have paid are ready, they mount the rope and take the long ride down.

\subsection{Oddities on Display}
A staggering variety of curios and oddities adorn the taproom. Traditionally, adventurers who recover a strange relic from Undermountain present it to Durnan as a trophy of their success. Other adventurers leave such curios to mark their visits to the tavern, or relinquish them after losing a bet with Durnan, who likes to wager on the fate of adventuring bands that enter the dungeon. Occasionally, something that strikes Durnan's fancy can be used to pay a bar tab.

\begin{DndTable}[header=Yawning Portal Taproom Curios]{cX}

d20 & Item\\
1 & A key carved from bone\\
2 & A small box with no apparent way to open it\\
3 & A mummified troglodyte's hand\\
4 & Half of an iron symbol of Bane\\
5 & A small burlap pouch filled with various teeth\\
6 & Burnt fragments of a scroll\\
7 & A lute missing its strings\\
8 & A bloodstained map\\
9 & An iron gauntlet that is hot to the touch\\
10 & A gold coin stamped with a worn, hawk-wing helm crest\\
11 & A troll finger, still wriggling\\
12 & A silver coin that makes no noise when dropped\\
13 & An empty jar\\
14 & A clockwork owl\\
15 & A blue, glowing crystal shard\\
16 & A statuette of a panther, wooden and painted black\\
17 & A piece of parchment, listing fourteen magical pools and their effects when touched\\
18 & A vial filled with a dark, fizzy liquid that is sealed and cannot be opened\\
19 & A feeler taken from a slain rust monster\\
20 & A wooden pipe marked with Elminster's sigil\\
\end{DndTable}

(See also Template in the next Story)

\subsection{A Typical Evening}
On quiet nights, guests in the Yawning Portal gather around a large fireplace in the taproom and swap tales of distant places, strange monsters, and valuable treasures. On busier nights, the place is loud and crowded. The balconies overflow with merchants and nobles, while the tables on the ground floor are filled with adventurers and their associates. Invariably, the combination of a few drinks and the crowd's encouragement induces some folk to pay for a brief trip down into Undermountain. Most folk pay in advance for a ride down and immediately back up, though a few ambitious souls might launch impromptu expeditions into the dungeon. Few such ill-prepared parties ever return.

Groups seeking to enter Undermountain for a specific reason generally come to the tavern during its quiet hours. Even at such times, there are still a few prying eyes in the taproom, lurkers who carry news of the comings and goings from Undermountain to the Zhentarim, dark cults, criminal gangs, and other interested parties.

\section{Starting the Story}
Kicking off a dungeon adventure can be as simple as having a mysterious stranger offer the characters a quest while they are at the Yawning Portal (or some other tavern). This approach is a cliché, but it is an effective one. Use the following two tables to generate a couple of details, then tailor the particulars of the quest and the quest giver to suit the adventure you plan to run.

\begin{DndTable}[header=Mysterious Stranger Quest]{cX}

d8 & Objective\\
1 & Recover a particular item\\
2 & Find and return with an NPC or monster\\
3 & Slay a terrible monster or NPC\\
4 & Guard a person while they perform a ritual\\
5 & Create an accurate map of part of the dungeon\\
6 & Discover secret lore hidden in the dungeon\\
7 & Destroy an object\\
8 & Sanctify part of the dungeon to a god of good\\
\end{DndTable}

\begin{DndTable}[header=Mysterious Stranger Secret]{cX}

d8 & Secret\\
1 & Intends to betray the party\\
2 & Unwittingly provides false information\\
3 & Has a secret agenda (roll another quest)\\
4 & Is a devil in disguise\\
5 & Has led other parties to their doom\\
6 & Is the charmed thrall of a mind flayer\\
7 & Is possessed by a ghost\\
8 & Is a solar in disguise\\
\end{DndTable}

(See also the Template below)

\subsection{Durnan}
The proprietor of the Yawning Portal is something of an enigma. Blessed with a seemingly limitless life span by treasures he brought back from his expedition nearly two centuries ago, he is as much a fixture in the tap room as the well.

Durnan is a man of few words. He expects to be paid for his time, and will offer insight and rumors only in return for hard cash. "We know the odds and take our chances," he says, whether he is breaking up a card game that has turned violent or refusing the pleas of adventurers trapped at the bottom of the well who are unable to pay for a ride up. Despite his stony heart, he is an excellent source of information about Undermountain and other dungeons, provided one can pay his price.

\includegraphics[width=0.4\textwidth]{images/../5e.tools/img/adventure/TftYP/Durnan.png}
\subsubsection{Personality Trait: Isolation}
It's a cruel world. All people have to fend for themselves. Self-sufficiency is the only path to success.

\subsubsection{Ideal: Independence}
Someone who can stand alone can stand against anything.

\subsubsection{Bond: The Yawning Portal}
This place is my only home. My friends and family are long gone. I love this place, but I try not to get attached to the people here. I'll outlive them all. Lucky me.

\subsubsection{Flaw: Heartless}
If you want sympathy, the Temple of Ilmater is in the Sea Ward. No matter how bad things are, you'll be gone in a blink of an eye.

\subsection{Other Denizens}
The Yawning Portal is host to a variety of regular visitors, most of whom offer services to adventurers. Chapter 4 of the \textit{Dungeon Master's Guide} provides plenty of resources for generating nonplayer characters. The following table provides some possibilities for why an individual is visiting the Yawning Portal.

\begin{DndTable}[header=Denizens]{cX}

d10 & Denizen\\
1 & Devotee of Tymora, encourages adventures to seek out quests, can cast bless\\
2 & Bored, retired adventurer, claims to have explored dungeon of note and can describe first few areas (20 percent chance of an accurate description)\\
3 & Heckler, mocks cowards and makes bets that adventurers won't return from an expedition\\
4 & Con artist, selling fake treasure maps (but a 10 percent chance that a map is genuine)\\
5 & Wizard's apprentice, carefully making exact sketches of various curios at her master's command\\
6 & Spouse of a slain adventurer, who pays the toll for anyone wanting to exit Undermountain and plots against Durnan\\
7 & Zhentarim agent, seeks rumors of treasure, tails any folk who return from Undermountain and notes their home base for future robbery\\
8 & Agent of the Xanathar, ordered to "steal the hat worn by the eighth person to enter the taproom this night"\\
9 & Magically preserved corpse in a coffin leaning against the bar; if asked about it, Durnan says, "He's waiting for someone," and nothing more\\
10 & Elminster, incognito; 10 percent chance he is on an errand of cosmic importance; otherwise, he's pressing Durnan for gossip\\
\end{DndTable}

\chapter{White Plume Mountain}
\DndDropCapLine{W}{}hite Plume Mountain has always been a subject of superstitious awe to the neighboring villagers. People still travel many miles to gaze upon this natural wonder, though few will approach it closely, as it is reputed to be the haunt of various demons and devils. The occasional disappearance of those who stray too close to the Plume reinforces this belief.


Thirteen hundred years ago, the wizard Keraptis was searching for a suitable haven where he could indulge his eccentricities without fear of interference. He visited White Plume Mountain, going closer than most dared to, and discovered the system of old lava-tubes that riddle the cone and the underlying strata. With a little alteration, he thought, these would be perfect for his purposes. The area already had a bad reputation, and he could think of a few ways to make it worse. So he disappeared below White Plume Mountain and vanished from the knowledge of the surface world.

Today, the once-feared name of Keraptis is not widely known even among learned scholars. Or it \textit{was} not widely known, that is, until several weeks ago, when three highly valued magic weapons named \textit{Wave}, \textit{Whelm}, and \textit{Blackrazor} disappeared from the vaults of their owners. Rewards were posted, servants hanged, even the sanctuary of the thieves' guild was violated in the frantic search for the priceless arms, but not even a single clue was turned up until the weapons' former owners each received a copy of the following note:

\begin{DndSidebar}[float=!b]{}
\textit{Search ye far or search ye near}



\textit{You'll find no trace of the three}



\textit{Unless you follow instructions clear}



\textit{For the weapons abide with me.}



\textit{North past forest, farm and furrow}



\textit{You must go to the feathered mound}



\textit{Then down away from the sun you'll burrow}



\textit{Forget life, forget light, forget sound.}



\textit{To rescue Wave, you must do battle}



\textit{With the Beast in the Boiling Bubble}



\textit{Crost cavern vast, where chain-links rattle}



\textit{Lies Whelm, past water-spouts double.}



\textit{Blackrazor yet remains to be won}



\textit{Underneath inverted ziggurat}.



\textit{That garnered, think not that you're done}



\textit{For now you'll find you are caught}



\textit{I care not, former owners brave}



\textit{What heroes you seek to hire.}



\textit{Though mighty, I'll make each one my slave}



\textit{Or send him to the fire.}



\end{DndSidebar}

All the notes were signed with the symbol of Keraptis.

White Plume Mountain has tentatively been identified as the "feathered mound" of the poem. The former owners of \textit{Wave}, \textit{Whelm}, and \textit{Blackrazor} are outfitting a group of intrepid heroes to take up the challenge. If the adventurers can rescue the weapons from this false Keraptis (for who can believe it is really the magician of legend, after thirteen hundred years?), the wealthy collectors have promised to grant them whatever they desire, if it is within their power to do so.

\section{Running the Adventure}
This version of \textit{White Plume Mountain} is designed for a group of 8th-level player characters. Your players will need both brains and brawn to successfully complete their mission, as there are situations here which cannot be resolved by frontal assault. If your players are unused to hack-proof dilemmas, they may find this adventure frustrating or even boring. But if your players are used to using their wits, they should find this an intriguing balance of problems and action. Unless you are used to mastering lengthy adventures, it will probably take more than one session for a party to investigate all three branches of the dungeon. If this is the case, it would be best if the party were required to leave the dungeon and reenter upon resumption of the game. If they stay in the nearest village (several miles away) they will be relatively safe, but if they camp near White Plume Mountain it would be a good idea to roll for random encounters.

\begin{DndSidebar}[float=!b]{Placing the Adventure}
White Plume Mountain is located in the Greyhawk campaign setting, in the northeastern part of the Shield Lands, near the Bandit Kingdoms and the Great Rift. Here are suggestions for where you can place the mountain in another world. Wherever you place it, the party may be required to journey to the vicinity through the wilderness. How they get there is up to you.



\subsubsection{Forgotten Realms}
The mountain can be placed near Mount Hotenow in the region of Neverwinter.



\subsubsection{Dragonlance}
Found near Neraka in the Khalkist Mountains, the mountain might be a place of interest not only to adventurers, but also to the armies of Takhisis.



\subsubsection{Eberron}
On the continent of Xen'Drik, the mountain could stand in the range known as the Fangs of Argarak.



\end{DndSidebar}

\subsection{Adventure Start}
The party has arrived at White Plume Mountain, which stands alone in a vast area of dismal moors and tangled thickets. They will probably arrange to leave their horses and possessions either at the nearest village (about 5 miles from the mountain) or hidden in the Dead Gnoll's Eye Socket, a small natural cave in the side of a hill about 2 miles from the Plume. There is really no other shelter available. The villagers know about the cave and may have mentioned it. If the party leaves no guard, they will just have to trust the villagers not to steal their belongings. (Dishonest villagers will have to weigh their fear of White Plume Mountain against their certain belief that the party will never be seen again.) The cave is easily barricaded to keep out unintelligent wandering monsters.

White Plume Mountain is an almost perfectly conical volcanic hill formed from an ancient slow lava leakage. It is about 1,000 yards in diameter at the base, and rises about 800 feet above the surrounding land. The white plume that gives the mountain its name and fame is a continuous geyser that spouts from the very summit of the mountain another 300 feet into the air, trailing off to the east under the prevailing winds like a great white feather. The spray collects in depressions downslope and merges into a sizable stream. Steam vents are visible in various spots on the slopes of the mountain, but none of them are large enough to allow entry.

Map 4.1 depicts a cross-section of the mountain, showing the lava pool and the shaft of the geyser. The numbers refer to key areas inside the mountain, showing their orientation with respect to one another.

The only possible entrance into the cone is a cave on the south slope known as the Wizard's Mouth. This cave actually seems to breathe, exhaling a large cloud of steam and then slowly inhaling, like a person breathing on a cold day. Each cycle takes about 30 seconds. Approaching the cave, the party will hear a whistling noise coinciding with the wind cycle. If it were not for the continuous roaring of the Plume, this whistling could be heard for a great distance.

The cave is about 8 feet in diameter and 40 feet long. At the end of the cave, near the roof, is a long, horizontal crevice about a foot wide. The air is sucked into this crack at great speed, creating the loud whistling noise and snuffing out torches. Shortly the rush of air slows down, stops for a couple of seconds, and then comes back out in a great blast of steam. This steam is not hot enough to scald anyone who keeps low and avoids the crevice, but it does make the cave very uncomfortable, like a hot sauna bath interrupted by blasts of cold air.

The ceiling and walls of the cave are slick with the condensed steam that runs down them. The floor is covered with several inches of fine muck. Only careful probing of the muck near the back of the cave will reveal a small trapdoor with a rusted iron ring set in it. Once the muck has been cleared away, a successful DC 20 Strength check is required to pull up the encrusted door. Magic such as \textit{knock} or \textit{passwall} can also help open or bypass the door.

Directly beneath the door is a 20-foot-square vertical shaft and the beginning of a spiral staircase that leads down.

\begin{DndSidebar}[float=!b]{}
\textbf{ABOUT THE ORIGINAL}



\textit{White Plume Mountain}, by Lawrence Schick, was originally published in 1979 as an adventure for the first edition of the D\&D game.



Schick related in the 2013 compilation \textit{Dungeons of Dread} that he wrote the adventure as a way of persuading Gary Gygax to hire him as a game designer.



\includegraphics[width=0.4\textwidth]{images/../5e.tools/img/adventure/TftYP-WPM/000-totyp-04-book.png}


\end{DndSidebar}

\subsection{Dungeon: General Features}
All corridors in the dungeon are 10 feet in height, and have been carved out of and, in some places, seemingly melted through solid rock. Unless otherwise noted, doors are 8 feet by 8 feet, made of oak and bound in iron. Though the doors are swollen by the dampness, and thus difficult to open (requiring a successful DC 15 Strength check), the wood is not by any means rotten.

The water on the floor is about 1 foot deep, and the floor itself is covered with slippery mud. Except where flights of steps lead up out of it, this scummy water covers the floors of all rooms and corridors. The water and mud reduce average movement by one-third (speed 30 becomes 20, speed 20 becomes 15), and will necessitate continuous probing of the floor by the party as they advance. It will be very difficult to keep silent, run (without falling), or depend on invisibility (waves and foot-shaped holes in the water give one away).

\includegraphics[width=0.4\textwidth]{images/../5e.tools/img/adventure/TftYP-WPM/001-wpm01.png}
\subsection{Random Encounters}
Check once every 10 minutes for random encounters by rolling a d12; an encounter occurs on a roll of 1. If an encounter is indicated, roll a d6 and refer to the table below.

These are monsters that Keraptis has released into the dungeon specifically for the purpose of giving the intruders a hard time. All will attack immediately. The ogres and the bugbears are magically controlled and cannot be persuaded to betray Keraptis.


\begin{DndTable}{cX}

d6 & Encounter\\
1 & 1 \textbf{black pudding}\\
2 & 1d4 + 1 bugbears\\
3 & 2 gargoyles\\
4 & 1 \textbf{invisible stalker}\\
5 & 1d3 ogres\\
6 & 1d2 wights\\
\end{DndTable}

\section{Locations in the Dungeon}
The following locations are identified on map 4.2.

\includegraphics[width=0.4\textwidth]{images/../5e.tools/img/adventure/TftYP-WPM/002-wpm02.png}
\includegraphics[width=0.4\textwidth]{images/../5e.tools/img/adventure/TftYP-WPM/003-wpm01.png}
\begin{DndSidebar}[float=!b]{The Legend of Keraptis}
In the original publication of \textit{White Plume Mountain}, "The Legend of Keraptis" was presented on the inside back cover. Although these details of the wizard's former life don't play a direct part in the adventure, a DM who shares this is information with the players can deepen the characters' understanding of the situation and strengthen their motivation for delving beneath the mountain.



✫ ✫ ✫



Well over a millennium ago, the wizard Keraptis rose to power in the valleys of the northern mountains, bringing the local warlords under his thumb with gruesome threats-threats that were fulfilled just often enough to keep the leaders in line. Under Keraptis's overlordship, the influx of rapacious monsters and raids from the wild mountains decreased markedly, dwindled, and then almost stopped. Seeing this, the populace did not put up much resistance to paying the wizard's heavy taxes and tithes, especially when stories were circulated of what happened to those who balked. Any nobles who protested disappeared in the night and were replaced by the next in the line of succession, who was usually inclined to be more tractable than the previous lord.



Gradually, as dissension was stilled, the taxes and levies became even more burdensome, until eventually the wizard was taking a great piece of everything that was grown, made, or sold in the valleys, including the newborn young of livestock. Around this time, numerous reports arose in the land concerning sudden madness, demonic possessions, and sightings of apparitions and undead. Furthermore, monstrous incursions into the settled lands began to increase as raiding parties of humanoids assaulted villages, and evil and fantastic monsters appeared from nowhere to prey upon the harried peasants. At the height of this unrest, Keraptis's tax collectors came forth with word of a new levy: one-third of all newborn children were henceforth to be turned over to the wizard.



That edict turned out to be the tipping point. As one, the people rose up to overwhelm the wizard's lackeys and marched on his keep, where, led by a powerful and good cleric and his ranger acolytes, they destroyed Keraptis's final guardians. The great wizard barely managed to escape, accompanied only by his personal bodyguard company of deranged and fanatical evil gnomes. Keraptis fled to the cities of the south and west, but wherever he went, his reputation preceded him, and he was unable to stay anywhere for long. Once again moving north, he came to the shores of the Lake of Unknown Depths, where he heard tales about haunted White Plume Mountain. After investigating further, he at last found the refuge he was looking for in the tangled maze of volcanic tunnels beneath the cone. He and his gnomes vanished into the shadow of the Plume, and humankind heard no more of the evil wizard.



That was almost thirteen hundred years before the present day. Now, seemingly, the hand of Keraptis is once again interfering in human affairs. If it is in truth the ancient wizard at work here, can he be thwarted before his power grows once more? What is his purpose in presenting this bizarre challenge to the world's heroes? There is only one way to find out.



\end{DndSidebar}

\subsection{1. Spiral Staircase}
\begin{DndReadAloud}{}
The staircase is badly rusted but appears to be sturdy. The air inside the passageway is warm, humid, and rather foul. You reach the bottom of the stairs with a splash—the floor is submerged beneath a foot of water!

\end{DndReadAloud}

The spiral staircase descends about 100 feet before ending in area 1. Sensitive characters will feel it thrumming to a continuous low vibration (this vibration from the Plume geyser will be noticeable everywhere in the dungeon). In the humidity, lamps and torches will burn fitfully and give off a lot of smoke.

Floating on the water are splotches of green and white subterranean algae. This algae or algae-like fungus also clings in patches to walls and ceilings. It is harmless, and can be found almost everywhere in the dungeon where there is water.

\subsection{2. Riddling Guardian}
A rather mangy, bedraggled \textbf{gynosphinx} squats in the space where the three corridors converge. A \textit{wall of force} along the south side separates the gynosphinx from those who approach from that direction. This wall of force is weaker than most, and can be brought down by \textit{disintegrate}, \textit{dispel magic}, or \textit{passwall}.

The sphinx will let the characters pass (by removing the wall of force) if they can answer the following riddle:

\begin{DndSidebar}[float=!b]{}
\textit{Round she is, yet flat as a board}



\textit{Altar of the Lupine Lords}



\textit{Jewel on black velvet, pearl in the sea}



\textit{Unchanged but e'erchanging, eternally}



\end{DndSidebar}

The answer is "the Moon." If the wall of force is destroyed or circumvented, the sphinx will attack.

\subsection{3. Hidden Slime}
Midway along the corridor that runs northeast, the floor is covered by a large patch of green slime (see "green slime" in chapter 5 of the \textit{Dungeon Master's Guide}). Because the slime is covered by water, it is not easily detectable, and characters might walk through it and not even notice they have done so until it has eaten through their boots and started on their feet.

\subsection{4. Glass Globes}
The door to this room appears normal on the way in: a large iron-bound oak door, swollen by dampness and difficult to open. When the characters have entered the room (or as many of them as are going in), the door will slam shut behind them. No tools, weapons, or magic available to the adventurers will open the door or prevent it from closing. Only the proper key inserted in the keyhole on the inside of the door will unlock it.

In the room, suspended from the ceiling by unbreakable wires, are nine silvered glass spheres, each about 2 feet in diameter. Unless otherwise noted, magically looking inside a sphere, such as with a \textit{clairvoyance} spell or a \textit{ring of x-ray vision}, will show that it contains some apparent treasure and a key. (Each sphere holds a key, but only one of the keys opens the door.) A good, hard crack with a weapon will shatter any of the spheres (each has AC 13 and 3 hit points), dropping its contents (if not caught) into the muck below.

Number the globes 1 through 9 for your own reference. The globes contain the following items:


\begin{description}
\item[Sphere 1.]
Three folded-up shadows, 300 worthless lead pieces, and a false key. (The shadows aren't visible to magical inspection.)
\item[Sphere 2.]
A \textit{spell scroll} of \textit{fear} and a false key.
\item[Sphere 3.]
Jewelry worth 12,000 gp, a false key, and an angry \textbf{air elemental}.
\item[Sphere 4.]
A \textit{potion of flying} and a false key.
\item[Sphere 5.]
Eleven worthless glass gems and a false key.
\item[Sphere 6.]
Phony glass jewelry, a false key, and a \textbf{gray ooze}. (The ooze fills the entire globe and isn't visible to magical inspection.)
\item[Sphere 7.]
A \textit{spell scroll} of \textit{hold person} and a false key.
\item[Sphere 8.]
The actual key and a silver ring. When the ring is released from the sphere, it speaks to the characters telepathically:
\end{description}

\begin{DndReadAloud}{}
"Stop before you put me on. I confer the following powers upon my wearer: \textit{invisibility}, \textit{haste}, immunity to charms, \textit{fly} once per day, \textit{detect magic}, and one \textit{wish}. I also provide the benefits of protection and spell turning. The only drawback is that once a year I permanently eat a small part of your life. I must be worn before I can leave this room; merely carrying me away is not possible. If ever I am removed from my wearer's finger, however, all my powers are lost. So you must decide right now who will wear me forever."

\end{DndReadAloud}

This situation is a basic loyalty and intelligence test. Will the party members cut each others' throats over the ring? Will they be suspicious enough to simply leave it alone? If they take time to think about it, they'll likely realize that the ring must be a hoax.

Someone who puts it on can cause it to exhibit any of the powers mentioned above, except for the \textit{wish} spell, as long as it remains in the room. While it is located here, the ring enables the wearer to produce the indicated spell effects at will (except for \textit{fly}), and it also acts as a \textit{ring of spell turning} and a \textit{ring of protection}.

Once the ring leaves the area, however, it has no abilities and can't talk.


\begin{description}
\item[Sphere 9.]
Gems worth 600 gp and a false key.
\end{description}

\subsection{5. Numbered Golems}
\begin{DndReadAloud}{}
Five flesh golems are clustered against the north wall. Each has a number on its chest: 5, 7, 9, 11, and 13. Number 5 says, "One of us does not belong with the others. If you can pick it out, it will serve you, and the others will allow you passage. If you pick the wrong one, we will kill you. You have sixty seconds."

\end{DndReadAloud}

The correct choice is 9, because it is not a prime number. Give the players an actual 60 seconds to figure it out.

If the characters give the wrong answer, roll initiative. The golems lumber forward to confront the intruders, trying to prevent anyone from moving through the door to the north. The monsters will pursue enemies that flee to the south, but they won't climb the stairs that led to area 6.

If the characters answer the riddle correctly, events unfold as promised: the golem numbered 9 becomes an ally of the party, accompanying the characters if they so desire, and the other four golems become inert again.

\subsection{6. Turnstile}
\begin{DndReadAloud}{}
A short flight of stairs leads up to a dry corridor. Just around the corner is a turnstile.

\end{DndReadAloud}

The characters will discover that the turnstile rotates in only one direction (counterclockwise). They can pass through it easily when moving to the east, but it will probably have to be destroyed on the way back. A golem or a strong character could rip it out with a successful DC 24 Strength (Athletics) check.

\subsection{7. Geysers and Chains}
\includegraphics[width=0.4\textwidth]{images/../5e.tools/img/adventure/TftYP-WPM/003-totyp-04-05.png}
\begin{DndReadAloud}{}
The door opens onto a stone platform in a large natural cave. Opposite the entrance in the distance is another stone platform. Between them, a series of wooden disks is suspended from the ceiling by massive steel chains. The cave floor seems to be covered by a pool of boiling mud.

\end{DndReadAloud}

The ceiling is 50 feet above the level of the platforms. The cave floor is 50 feet below. Two spots in the mud are the locations of geysers. The northern one erupts once every 5 minutes, the southern one every 3 minutes. The stone platform opposite the entrance is approximately 90 feet away.

The disks are about 4 feet in diameter and 3 feet apart. Each disk is attached to its chain by a giant staple fixed in its center. The disks swing freely and will tilt when weight is placed upon them.

The disks and the chains, as well as the walls of the cavern, are covered with a wet, slippery algal scum. This coating gives off a feeble phosphorescent glow. Climbing the chains or the walls requires a successful DC 15 Strength (Athletics) check.

When the geysers erupt, they reach nearly to the roof of the cavern, and creatures holding onto the disks or the chains might be washed off to fall into the boiling mud. The damage a creature takes from a geyser depends on how close a creature is to the geyser when it erupts (see the table below). Creatures that succeed on a DC 15 Dexterity saving throw take half damage.

In addition, a creature that is on a disk or holding onto a chain when a geyser erupts must succeed on a Strength saving throw (see the table for DCs) or be knocked off and fall into the boiling mud.


\begin{DndTable}{XXX}

Location & Damage & DC\\
Adjacent to geyser & 27 (5d10) fire damage & 14\\
One disk away & 22 (4d10) fire damage & 13\\
Two disks away & 16 (3d10) fire damage & 12\\
Three disks away & 11 (2d10) fire damage & 11\\
Four disks away & 5 (1d10) fire damage & 10\\
Anywhere else in the area & 3 (1d6) fire damage & —\\
\end{DndTable}

Any creature that falls into the boiling mud takes 44 (8d10) fire damage at the start of each of its turns for as long as it remains in the mud.

\includegraphics[width=0.4\textwidth]{images/../5e.tools/img/adventure/TftYP-WPM/004-wpm03.png}
\subsection{8. Coffin}
\begin{DndReadAloud}{}
The door opens into an area of utter blackness.

\end{DndReadAloud}

\includegraphics[width=0.4\textwidth]{images/../5e.tools/img/adventure/TftYP-WPM/005-totyp-04-06.png}
The room beyond the door is the lair of a \textbf{vampire} named Ctenmiir. He is compelled by a curse to remain here in a trance except when roused to defend his treasure, which lies in a niche in the floor under his coffin. The vampire automatically awakes at the approach of intruders.

The room is affected by a \textit{darkness} spell, which the vampire is not hindered by because he has truesight. If the vampire needs to leave this area to pursue intruders, the door to the room is littered with tiny holes through which he can pass in mist form. If the magical darkness is dispelled, it renews again automatically at dawn.

\subparagraph{Treasure}
The space beneath the coffin contains \textit{Whelm}, a sentient warhammer (see "whelm" in chapter 7 of the \textit{Dungeon Master's Guide}), 10,000 sp and 9,000 gp in six leather sacks, a \textit{potion of mind reading}, and three Spell Scroll (\textit{conjure minor elementals}, \textit{dispel magic}, and \textit{magic mouth}).

\subsection{9. Pool and Drain}
As the characters move along the water-covered corridor, they can see that the water is deeper in a small circular area to the east. This space is a 10-foot-deep pit. At the bottom is a wheel connected to a valve. Turning the wheel requires a successful DC 20 Strength check.

When the wheel is turned, a channel will open, and all the water in the wet corridors will drain out in 1 hour. Also at the bottom of the pit is a secret door that is concealed by illusion magic. The door cannot be detected by sight, but can be discovered through the use of magic or by someone who examines the area by touch and succeeds on a DC 15 Intelligence (Investigation) check. The doorway leads to Keraptis's Indoctrination Center (see "Escaping the Dungeon" at the end of the adventure).

\subsection{10. Deceptively Deep Room}
\begin{DndReadAloud}{}
Ahead is what appears to be a water-covered room, with steps rising out of the muck on the far side.

\end{DndReadAloud}

Most of the area is actually a 15-foot-deep pool. The only shallow spaces are those that run mostly around the perimeter of the room, where the water is only 1 foot deep.

The deep area that makes up most of the room is inhabited by two kelpies (see appendix B). As the characters move into the room, the kelpies will rise to the surface, and each will attempt to charm a character.

In the eastern wing of the pool are two partly enclosed areas with entrances that are accessible only from beneath the surface of the water. These areas are partially covered by a roof along the east edge, where the water is only 1 foot deep. The southern chamber is the kelpies' lair, which contains their treasure. The northern chamber is empty.

\subparagraph{Treasure}
Scattered about the kelpies' lair are 600 gp, a piece of jewelry worth 2,000 gp, and a suit of +1 Chain Mail.

\subsection{11. Spinning Cylinder}
\begin{DndReadAloud}{}
The stone corridor changes abruptly to a spinning cylinder, apparently made of some light-colored metal. The inner surface rotates rapidly. It is painted in a dizzying black-and-white spiral pattern, and it glistens as if coated with some substance.

\end{DndReadAloud}

The 30-foot-long cylinder is 10 feet in diameter and spins counterclockwise at about 10 feet per second. The inner surface is covered with slippery oil. It is possible to slide through the cylinder by propelling oneself along the floor, but walking through this area without being knocked prone requires a successful DC 20 Dexterity (Acrobatics) check.

\subsection{12. Burket's Guardpost}
Watching through the arrow slit at the end of the passage is an alert guard named Burket (LE male human \textbf{veteran}). If he sees intruders approaching, he will wait until they are halfway through the spinning cylinder and then ignite the slippery oil with a flaming arrow.

The oil burns for 2d4 rounds if the fire is not extinguished. On the first round, the burning oil deals 9 (2d8) fire damage to any creature that enters the spinning cylinder or starts its turn inside the cylinder. On later rounds, it deals 2 (1d4) fire damage to such creatures.

\subparagraph{Development}
After Burket shoots, he then warns his lover, the wizard \textbf{Snarla}, to close and lock the shutter over the arrow slit and move to defend the door that adjoins the corridor. They stand ready to engage anyone that enters.

\textbf{Snarla} is a \textbf{werewolf}, her Intelligence is 16 (+3), and she has the following additional feature, which increases her challenge rating to 5 (1,800 XP):

\begin{DndSidebar}[float=!b]{}
\subparagraph{Spellcasting}
\textbf{Snarla} is a 6th-level spellcaster. Her spellcasting ability is Intelligence (spell save DC 14, +6 to hit with spell attacks). She has the following wizard spells prepared:

Cantrips (at will): \textit{fire bolt}, \textit{light}, \textit{mage hand}, \textit{shocking grasp}

1st level (4 slots): \textit{magic missile}, \textit{shield}, \textit{thunderwave}

2nd level (3 slots): \textit{mirror image}, \textit{web}

3rd level (3 slots): \textit{dispel magic}, \textit{fear}, \textit{haste}, \textit{stinking cloud}



\end{DndSidebar}

If Burket is killed or if \textbf{Snarla} finds herself in a bad situation, she changes into werewolf form and attacks with desperate savagery, giving her advantage on all her attack rolls.

If she is captured alive and made to talk, she will tell the party only that she is charged with keeping the kelpies and other denizens of the dungeon fed. There are strange gaps in her memory concerning her employer or any section of the dungeon other than her own. She has never been past the doors at area 14. Burket knows even less than she does.

\subparagraph{Treasure}
The room has a couple of benches and a table, upon which are a large candlestick (worth 10 gp) and \textbf{Snarla}'s spellbook. The book contains the spells that she has prepared and no others. It is protected by an explosive runes \textit{glyph of warding} that deals 5d6 fire damage.

\subsection{13. Snarla's Sanctum}
\begin{DndReadAloud}{}
Unlike the room you just left, this place is beautifully decorated. The floor is covered by fine rugs, the walls by erotic tapestries and shimmering curtains, the ceiling by an intricate mosaic depicting a summer sky dotted with fleecy clouds. In the northeast corner is a large and lavishly covered bed, strewn with cushions. Next to it on a low table is a buffet of sweetmeats, cakes, and other delicious-looking comestibles. In the northwest corner of the room is a brass-bound oak chest.

\end{DndReadAloud}

Anyone who tries out the bed will find that it feels quite uncomfortable, and anybody who samples the food will be disappointed in the extreme, finding it tough and not very tasty.

In fact, the room is covered in illusions; a \textit{true seeing} spell or similar illusion-piercing magic reveals that the opulent bed is only an old straw mattress, and the delicious treats are just ordinary rations. The walls, floor, and ceiling of the room are quite bare. A \textit{dispel magic} spell will automatically remove the illusion.

\subparagraph{Treasure}
Only the brass-bound chest appears as it actually is. The chest must be opened while uttering a command word known only to \textbf{Snarla}, or it will emit a \textit{stinking cloud} spell that lasts for 1 minute. Inside are 400 ep, 300 gp, and seven gems worth a total of 1,300 gp.

\subsection{14. Flood Doors}
The three doors along the corridor are made of thick metal, their edges flanged so that they overlap the door jamb on the north side and thus can be opened only by pivoting them to the north. The north side of each door has a handle so that it can be pulled open from that direction.

These barriers are emergency doors, whose purpose is to prevent the dungeon from being flooded by the boiling lake at area 15, in case of an "accident."

\subsection{15. Boiling Lake}
This boiling lake is several hundred feet deep, extending down to the red-hot rock below, and reaching nearly to the ceiling of the cavern it occupies, 50 feet above the level of the sunken ledge described in area 17. It is fed by an underground stream that enters from the northwest at a depth 100 feet below that of the ledge. Its run-off flows through a channel to the east, above the ledge, near the ceiling of the cavern.

Any creature that enters the boiling lake takes 44 (8d10) fire damage immediately and again at the start of each of its turns for as long as it remains in the lake.

\includegraphics[width=0.4\textwidth]{images/../5e.tools/img/adventure/TftYP-WPM/006-wpm04.png}
\subsection{16. Blow Hole}
The run-off from the boiling lake cascades down through a series of near-vertical lava-tubes to the base of the blow hole, 800 feet below the level of the dungeon. There the water strikes molten rock and is instantly converted to steam. It is ejected up the blow hole and out the top of the volcanic cone, forming the continuous geyser of White Plume Mountain. The boiling water here is just as dangerous as the water in area 15.

\subsection{17. The Boiling Bubble}
A sunken stone ledge projects out into the boiling lake. The corridor from the dungeon continues out into the lake under a magical force field that keeps out the water by forming a sort of elastic skin. The shape of the corridor is not square in cross-section, but rather semicircular, as if a series of hoops were supporting the ceiling.

The protective skin is soft, resilient, and uncomfortably warm to the touch. Under any pressure it immediately becomes taut, and any character unwise enough to puncture it with a piercing weapon will cause a stream of scalding water to rush into the corridor, hopefully burning the idiot who made the hole for 2 (1d4) fire damage). Thereafter, any creature that enters the space with this stream of scalding water or that starts its turn there takes that damage. The skin will not "heal" once it is compromised. Major damage to the skin, as from a slash with a sword or an axe, will collapse the field like a deflating balloon in 1d6 rounds.

After 30 feet the corridor gives way to an oval-shaped, domed area enclosed by the protective skin. Here lives the guardian of the treasure, a Huge \textbf{giant crab} that has the following changes, which increase its challenge rating to 8 (3,900 XP):


\begin{itemize}
\item It has 161 (14d12 + 70) hit points.
\item Its Strength and Constitution are 20 (+5).
\item On one of its claws it wears a rune-covered copper band that makes it immune to being charmed, frightened, and paralyzed. (The copper band is worthless as a treasure, for the magic is keyed to this crab.)
\item It has an improved claw attack:
\end{itemize}

\begin{DndSidebar}[float=!b]{}
\subparagraph{Claw}
\textsl{Melee Weapon Attack:} +9 to hit, reach 10 ft., one target. Hit: 27 (4d10 + 5) bludgeoning damage, and the target is grappled (escape DC 14). The crab has two claws, each of which can grapple only one target.



\end{DndSidebar}

The crab will intelligently attack intruders, being careful not to bump the protective skin walls. The crab is experienced in fighting in this manner, as is evidenced by the bones scattered about, but the characters are not. You will have to watch for characters whose actions might rip the water skin, especially any foolish enough to use two-handed weapons or a violent spell such as \textit{fireball} or \textit{lightning bolt}. Such people are likely to get the whole party boiled.

\includegraphics[width=0.4\textwidth]{images/../5e.tools/img/adventure/TftYP-WPM/007-totyp-04-08.png}
\subparagraph{Treasure}
At the north end of the domed area is a heavy chest firmly attached to the floor. In it is \textit{Wave}, a sentient trident (see "wave" in chapter 7 of the \textit{Dungeon Master's Guide}), 1,000 gp in small sacks, twenty gems (two big ones worth 1,000 gp each, one big one worth 5,000 gp, and seventeen others worth a total of 3,935 gp), \textit{goggles of night}, and a Stone of Good Luck.

\subparagraph{Development}
A character who grabs \textit{Wave} while the protective skin is collapsing can save the lives of those nearby by using the trident as a \textit{cube of force}. \textit{Wave} will instantly make its bearer aware of this property and allow the bearer to instantly become attuned to it if that person worships a god of the sea or is willing to convert on the spot.

Characters protected by the cube will probably end up being blown out the geyser at the top of the mountain. The air-filled cube will float, drain down the cascade, and be ejected from the Plume—a rocky ride.

Characters could also survive the boiling lake with a combination of immunity to fire damage and the ability to breathe water.

\subsection{18. Hall Pit}
Halfway down the corridor that leads to the northwest is a 10-foot-long, 10-foot-deep open pit, filled with and hidden by water. If the characters aren't testing the floor ahead of them, those in the first rank must succeed on a DC 15 Dexterity saving throw to avoid falling into the pit.

\subsection{19. Metal-Heating Corridor}
\begin{DndReadAloud}{}
A series of copper-colored metal plates lines the walls of the path before you.

\end{DndReadAloud}

The plates are 6 feet high and 6 feet wide, and cannot be damaged or removed. They produce an invisible electrical field that extends from floor to ceiling throughout the 70-foot-long corridor. The field isn't directly harmful, but metal objects that pass between the plates become heated. Metal will become uncomfortably warm after moving 20 feet into the field, painfully hot after 30 feet, and hot enough to deal fire damage at 40 feet and beyond. The field affects armor, weapons, equipment, and even treasure made of metal.

A character in metal armor who tries to move through this corridor would take damage as follows: 4 (1d8) fire damage at 40 feet, another 9 (2d8) fire damage at 50 feet, an additional 16 (3d10) fire damage at 60 feet, and another 22 (4d10) fire damage at 70 feet.

Characters not wearing metal armor but carrying metal weapons or equipment will feel only slight discomfort when passing between the plates. Metal carried in wrappings of cloth will burn through by the 50-foot mark, and it will similarly burn through leather by the 60-foot mark. Armored characters might have no recourse other than to remove their armor, then drag, push, or use magic to get their metal armor and weapons through the corridor, and then suit up again. Armor pulled through the corridor by ropes will heat up enough to burn through the ropes at 60 feet, leaving a pile of hot metal lying in the water.

The only sovereign remedies for this dilemma are area effects that deal cold damage, such as \textit{ice storm} or \textit{cone of cold}, which will nullify the effect long enough for the party to dash through.

\subsection{20. Ghoul Ambushers}
Behind the secret door, eight ghouls wait in ambush for intruders to come through the heat-induction corridor. These ghouls wear amulets that make them immune to the Turn Undead ability.

\subsection{21. Stairs Up}
These stairs lead up out of the water to the dry corridor beyond.

\subsection{22. Frictionless Trap}
When the characters pass through the door from area 21, describe the features of the room beyond depending on how much of the chamber they can see.

\begin{DndReadAloud}{}
The path to the west is broken by a sizable gap, and you can see the glint of metal at the bottom of the opening. The floor beyond this area has a silvery sheen.

In the distance you can see another hole, beyond which is a patch of floor that adjoins the western wall.

\end{DndReadAloud}

The openings in the floor are pits, each 5 feet wide and 10 feet deep. Their bottoms are lined with rusty, razor-like blades. Anyone who falls in takes 6 (1d12) damage and must succeed on a DC 15 Constitution saving throw or contract a disease called super-tetanus (see below).

The walls, ceiling, and floor of the area between the pits are covered with a substance that is totally frictionless. This silvery stuff is inert and utterly unaffected by any force, magical or otherwise. Anything that alights on the surface will move in the direction of its last horizontal impetus, bouncing off the walls (if it strikes them) like a billiard ball, until it slides into a razor-pit. It is impossible for a creature even to stand still on the surface, for the slightest movement in any direction would unbalance it enough to send it moving in that direction.

Magic that augments movement (such as \textit{fly}, \textit{jump}, and \textit{levitate}) and teleportation effects do not work in this room.

The line near the west edge of the room is the illusion of a wall—the actual wall is 10 feet farther away. Objects that make contact with the illusory wall will pass through it and seemingly disappear. This fact tends to frustrate (at least initially) any scheme for attaching a rope to the west wall from afar.

In fact, however, one way to get safely across is to get a rope strung through the room and fastened securely at both ends. Once this is done, characters can pull themselves across the floor and are able to control their speed. A clever party might be able to come up with other methods. Ingenuity is required.

\subparagraph{Super-Tetanus}
A creature that contracts the disease of super-tetanus is wracked with pain as its heart races and its muscles spasm hard enough to break its bones. The creature takes 11 (2d10) damage at the start of each of its turns. If a victim is not cured by other means, it can repeat the saving throw at the end of every minute after becoming exposed, ending the effect on itself with a successful save.

\subsection{23. Floating Stream}
\begin{DndReadAloud}{}
Water not only flows through this room, it floats.

Entering a hole in the western wall, two feet off the floor, is a stream seemingly suspended in mid-air. It flows out of another hole near the northeast corner. The water is about three feet deep. You can see a few blind cave fish being carried along in the brisk current. On either side of the door are a total of six kayaks, each able to carry two riders, but there are no paddles to be found.

\end{DndReadAloud}

The stream is hemispherical in cross-section where it flows through the room. The entrance and exit holes are 6 feet in diameter; they are the ends of a tunnel that connects this room with area 24.

The water is lukewarm. Objects can be thrust through the sides of the stream, but no water other than a few drops will escape. It would even be possible to walk right through the stream, but only a very strong creature could do so without being swept off its feet (requiring a DC 15 Strength saving throw). The stream flows quite quickly from west to east.

If characters choose to go boating into the unknown tunnel, they will have to figure out how to get in the kayaks once the kayaks are placed in the stream, the surface of which is 5 feet off the floor. These vessels tip over easily. If they successfully board the kayaks, the characters will bump along through a twisting tunnel 6 feet across. They can regulate their speed by pushing against the walls. Eventually they will emerge into area 24.

\subsection{24. Sir Bluto's Guardpost}
A fallen knight named Sir Bluto Sans Pite (a NE \textbf{champion}; see appendix B) and his eight minions (NE human knights) wait here to ambush any who come through the tunnel. They will be alerted by disturbances in the flow of water as the party tampers with it upstream. The stream that flows through this room is suspended in the same manner as in area 23. It continues from north to south out of area 24 and back to area 23, completing the circuit.

Sir Bluto's knights work in two teams of four, with two team members on each side of the stream. When a kayak comes out, a team will throw a large net over it and attempt to drag it and its occupants out of the stream and down onto the floor. When it falls, they move in to finish off the occupants with swords. To escape from a net, a creature can use its action to make a DC 10 Strength check, freeing itself or another creature within its reach on a successful check. Dealing 5 slashing damage to the net (AC 10) also frees a creature without harming it; any other creature in the same net can then use its action to free itself through that opening.

Sir Bluto was a respected Knight of the Realm before his indictment in the River of Blood mass-murder case. His mysterious disappearance from prison left even the Royal Magician-Detectives baffled, and a reward of 10,000 gp was posted for his capture. Someone in the party is sure to recognize his one-of-a-kind face.

\subparagraph{Treasure}
In addition to his weapons and armor, Sir Bluto wears \textit{boots of striding and springing}. In addition, he carries the key that opens the secret doors at area 25.

\subsection{25. Magical Secret Doors}
The secret doors at either end of the narrow corridor will reveal themselves and open only to the bearer of Sir Bluto's magic key.

\subsection{26. Terraced Aquarium}
\begin{DndReadAloud}{}
You are looking out and down into an enormous chamber defined by terraced steps that ring the entire area and descend toward a central enclosure.

A few feet south of the door, water laps at the edge of the stone floor. Looking to the side, you can see that the ring of the terrace that lies beyond is a ten-foot-deep ring of water held in by a nearly transparent wall. Another of these watery steps in the terrace is somewhat lower than the first; at the level in between these aquariums, the terrace steps are dry, but the area is still enclosed by a glassy wall like the others.

\end{DndReadAloud}

Each tier in the room is 10 feet in depth and width. The three middle levels are enclosed by 10-foot-tall glass walls that keep the inhabitants of those tiers confined. The two on either side are filled with water. All the creatures in this room have been charmed and ordered to stay in their areas as long as their glass walls are intact.

\includegraphics[width=0.4\textwidth]{images/../5e.tools/img/adventure/TftYP-WPM/008-wpm05.png}
A 10-foot-square area of glass wall has AC 15, and 20 hit points. The glass can also be broken with a successful DC 15 Strength check.

The only exit from the room is the door that leads south out of the bottom tier. It opens into a corridor that passes under the higher tiers of the room.

\subparagraph{Breaking Glass}
If any of the glass walls are broken, a weak \textit{wall of force} will activate immediately in front of the door on the bottom tier (preventing water from forcing the door open and escaping into the corridor beyond). If the glass is broken on both of the aquariums, the volume of water is sufficient to fill the two lower tiers and cover the floor of the middle tier to a depth of about 2 feet. The water will slowly drain out of four small drains in the corners of the bottom tier, but it will take a good three hours to do so. Of course, characters who choose to wait will be subject to the possibility of random encounters. Once the water is finally gone, the wall of force will disappear.

If the characters destroy the wall of force, the pressure of the water will push open the door and the water will rush into the corridor beyond, pulling along any swimming characters and miscellaneous debris nearby. The water will collide with another wall of force that covers the door to area 27, and then begin draining out through a grating in the floor in the last 10 feet of the passage. The water will take only 20 minutes to drain out through this grating. (The water drains straight down through an old lava tube to a large cave with no other exits.) When the water is gone, the wall of force barring entry to area 27 dissipates.

\subparagraph{Creatures}
The occupants of each level will be randomly distributed when the party enters, but as the intruders come near, the monsters will move to follow, expecting to be fed. They are accustomed to live food, and will ignore dead meat or other nonliving sustenance. They consider any living creature that enters their domain as food, and will attempt to eat it.

No creatures inhabit the top tier of the terrace. The aquarium tier inside that area is occupied by six \textbf{giant crayfish} (see appendix B). The next lower tier is a dry level where the glass walls enclose six giant scorpions. Beneath and inside that is an aquarium that holds four sea lions (see appendix B).

On the bottom tier are three wing-clipped manticores that are unable to fly. The manticores will not hesitate to fire their spikes at any they recognize as intruders.

\subparagraph{Treasure}
A safe is set in the north wall opposite the door on the bottom tier. If it is opened incorrectly (that is, the trap not removed), a vibration device in the wall is triggered that will shatter the glass walls in this area in 1d6 rounds. The safe contains 6,000 sp and one piece of jewelry worth 3,000 gp.

\includegraphics[width=0.4\textwidth]{images/../5e.tools/img/adventure/TftYP-WPM/009-totyp-04-10.png}
\subsection{27. Luxurious Prison}
\begin{DndReadAloud}{}
Lavish furnishings and decorations are everywhere in this large room. The floor is strewn with rugs and cushions, and tapestries cover the walls. A hookah as tall as an adult human stands in one corner. The largest piece of furniture is a sumptuous divan.

\end{DndReadAloud}

\includegraphics[width=0.4\textwidth]{images/../5e.tools/img/adventure/TftYP-WPM/010-totyp-04-11.png}
The room is the residence of Qesnef, an \textbf{oni} who lost a bet with Keraptis and as a result must confine himself in these luxurious surroundings, guarding his treasure, for 1,001 years. A \textit{magic mouth} spell warns him of the approach of trespassers, so he will disguise himself by using his Change Shape ability to take the form of a doughty halfling warrior, claiming to be someone who has been trapped by the evil wizard.

\subparagraph{Treasure}
Qesnef's valuables have been casually shoved beneath the divan. The hoard includes \textit{Blackrazor}, a sentient greatsword (see "blackrazor" in chapter 7 of the \textit{Dungeon Master's Guide}), 1,000 ep, 200 pp, four pieces of jewelry worth a total of 11,000 gp, a Potion of Greater Healing, a \textit{scroll of protection} (fiends), and \textit{armor of vulnerability (slashing)}.

In addition, Qesnef wears a \textit{ring of protection} on his left hand and a \textit{ring of spell storing} (with two \textit{mirror image} spells in it) on his right.

\section{Escaping the Dungeon}
If the characters obtain two or even all three of the magic weapons and are finally leaving for good, they might be stopped at area 2 by the return of the \textit{wall of force}. A voice will speak to them out of the air:

\begin{DndReadAloud}{}
"Not thinking of leaving, are you? You've been so very entertaining. I just couldn't think of letting you go, especially with those little items of mine. And since you've eliminated all of their guardians, why, you'll simply have to stay... to take their places. I'll have to ask you to leave all of your ridiculous weapons behind and let Nix and Nox escort you to the Indoctrination Center. I'll be most disappointed if you cause me any trouble and Nix and Nox have to eliminate you. Don't worry—you'll like it here."

\end{DndReadAloud}

\subparagraph{Creatures}
The wall of force disappears, but coming up the south passage are Nix and Nox, two efreet. If the party can get past them, they're home free.

Of course, this whole episode can be omitted if the party has already taken too much damage. Conversely, if the characters have had too easy a time of it, this final challenge can be made tougher by the addition of one or two more efreet (called Box and Cox).

If, for some foolish reason, the party decides to comply with Keraptis's request and go with Nix and Nox to the Indoctrination Center, you will have to play it by ear. It's not too difficult—use your imagination and make it up as you go. Just make sure that the characters are extremely sorry they ever decided to submit to Keraptis's demands. They probably will end up as the brainwashed new guards in Keraptis's renewed version of White Plume Mountain's dungeon.

\appendix
\chapter{Magic Items}
The magic items that are introduced in this book are detailed here in alphabetical order. The adventure in which an item appears is given at the end of its description.


\begin{itemize}
\begin{multicols}{3}
\item \textit{Amulet of Protection from Turning}
\item \textit{Balance of Harmony}
\item \textit{Bracelet of Rock Magic}
\item \textit{Eagle Whistle}
\item \textit{Hell Hound Cloak}
\item \textit{Loadstone}
\item \textit{Mirror of the Past}
\item \textit{Night Caller}
\item \textit{Potion of Mind Control}
\item \textit{Robe of Summer}
\item \textit{Shatterspike}
\item \textit{Spear of Backbiting}
\item \textit{Stone of Ill Luck}
\item \textit{Wand of Entangle}
\item \textit{Waythe}
\end{multicols}

\end{itemize}

\chapter{Creatures}
This appendix details creatures and nonplayer characters that are mentioned in this book and that don't appear in the \textit{Monster Manual}. That book's introduction explains how to interpret a stat block.

Some of these creatures are available in \textit{Volo's Guide to Monsters} but are reproduced here for your convenience.

The creatures are presented in alphabetical order.


\begin{itemize}
\begin{multicols}{4}
\item \textbf{Animated Table}
\item \textbf{Barghest}
\item \textbf{Centaur Mummy}
\item \textbf{Champion}
\item \textbf{Choker}
\item \textbf{Conjurer}
\item \textbf{Deathlock Wight}
\item \textbf{Dread Warrior}
\item \textbf{Duergar Spy}
\item \textbf{Enchanter}
\item \textbf{Evoker}
\item \textbf{Giant Crayfish}
\item \textbf{Giant Ice toad}
\item \textbf{Giant Lightning eel}
\item \textbf{Giant Skeleton}
\item \textbf{Giant Subterranean lizard}
\item \textbf{Greater Zombie}
\item \textbf{Illusionist}
\item \textbf{Kalka-Kylla}
\item \textbf{Kelpie}
\item \textbf{Leucrotta}
\item \textbf{Malformed Kraken}
\item \textbf{Martial Arts Adept}
\item \textbf{Nereid}
\item \textbf{Necromancer}
\item \textbf{Ooze master}
\item \textbf{Sea Lion}
\item \textbf{Sharwyn Hucrele}
\item \textbf{Sir Braford}
\item \textbf{Siren}
\item \textbf{Tarul var}
\item \textbf{Tecuziztecatl}
\item \textbf{Thayan Apprentice}
\item \textbf{Thayan Warrior}
\item \textbf{Thorn Slinger}
\item \textbf{Transmuter}
\item \textbf{Vampiric Mist}
\item \textbf{White Maw}
\item \textbf{Yusdrayl}
\end{multicols}

\end{itemize}

\chapter{Credits}

\begin{description}
\begin{multicols}{2}
\item[Compilers.]
Kim Mohan, Mike Mearls

\item[Lead Rules Developer.]
Jeremy Crawford

\item[Fifth Edition Conversion.]
Chris Sims, Sean K Reynolds, Jennifer Clarke Wilkes

\item[Managing Editor.]
Jeremy Crawford

\item[Editors.]
Kim Mohan, Michele Carter

\item[Editorial Assistance.]
Chris Dupuis, Ben Petrisor, Matt Sernett

\item[Art Director.]
Kate Irwin

\item[Additional Art Direction.]
Shauna Narciso, Richard Whitters

\item[Graphic Designer.]
Emi Tanji

\item[Cover Illustrator.]
Tyler Jacobson

\item[Interior Illustrators.]
Mark Behm, Eric Belisle, Zoltan Boros, Noah Bradley, Sam Carr, Jedd Chevrier, Bud Cook, Olga Drebas, Wayne England, Lake Hurwitz, Izzy, Tyler Jacobson, Titus Lunter, Brynn Metheney, Scott Murphy, Claudio Pozas, Ned Rogers, Chris Seaman, Cory Trego-Erdner, Franz Vohwinkel, Mark Winters, Sam Wood, Ben Wootten

\item[Cartographers.]
Jason A. Engle, Rob Lazzaretti, Mike Schley, Ben Wootten

\item[Producer.]
Stan!

\item[Project Manager.]
Heather Fleming

\item[Product Engineer.]
Cynda Callaway

\item[Imaging Technicians.]
Sven Bolen, Carmen Cheung, Kevin Yee

\item[Art Administration.]
David Gershman

\item[Prepress Specialist.]
Jefferson Dunlap

\item[Other D\&D Team Members.]
Bart Carroll, John Feil, Trevor Kidd, Adam Lee, Christopher Lindsay, Shelly Mazzanoble, Christopher Perkins, Hilary Ross, Liz Schuh, Nathan Stewart, Greg Tito, Shawn Wood

\end{multicols}

\end{description}

\section{Credits from the Original Adventures}

\begin{description}
\begin{multicols}{2}
\section{Tomb of Horrors (1978)}

\begin{description}
\item[Design.]
Gary Gygax

\end{description}

\section{White Plume Mountain (1979)}

\begin{description}
\item[Design.]
Lawrence Schick

\item[Editing and Suggestions.]
Mike Carr, Allen Hammack, Harold Johnson, Tim Jones, Jeff Leason, Dave Sutherland, Jean Wells

\item[Art.]
Dave Sutherland, Erol Otus, Darlene Pekul, Jeff Dee, David S. LaForce, Jim Roslof, Bill Willingham

\end{description}

\section{The Hidden Shrine of Tamoachan (1980)}

\begin{description}
\item[Design.]
Harold Johnson, Jeff R. Leason

\item[Able Assistance.]
Dave Cook, Lawrence Schick

\item[Editing.]
Harold Johnson

\item[Editing and Production.]
Dave Cook, Jeff R. Leason, Lawrence Schick

\item[Illustrations.]
Erol Otus, Jeff Dee, Gregory K. Fleming, David S. LaForce, David C. Sutherland III

\end{description}

\section{Against the Giants (1981)}

\begin{description}
\item[Design.]
Gary Gygax

\item[Editing.]
Mike Carr, Timothy Jones, Jon Pickens, Lawrence Schick

\item[Art.]
David C. Sutherland III, David A. Trampier, Jeff Dee, David S. LaForce, Erol Otis, Bill Willingham

\end{description}

\section{The Sunless Citadel (2000)}

\begin{description}
\item[Design.]
Bruce R. Cordell

\item[Editing.]
Miranda Horner

\item[Cartography.]
Todd Gamble

\item[Illustrations.]
Dennis Cramer, Todd Lockwood

\end{description}

\section{The Forge of Fury (2000)}

\begin{description}
\item[Design.]
Richard Baker

\item[Editing.]
Miranda Horner

\item[Cartography.]
Todd Gamble

\item[Illustrations.]
Dennis Cramer, Todd Lockwood

\end{description}

\section{Dead in Thay (2014)}

\begin{description}
\item[Design.]
Scott Fitzgerald Gray

\item[Editing.]
Ray Vallese

\item[Cartography.]
Mike Schley

\item[Illustrations.]
Eric Belisle, Sam Carr, Tyler Jacobson, Miles Johnstone, Mark Winters

\end{description}

\end{multicols}

\end{description}

\section{On the Cover}
\includegraphics[width=0.4\textwidth]{images/../5e.tools/img/adventure/TftYP/credits.png}
\includegraphics[width=0.4\textwidth]{images/../5e.tools/img/adventure/TftYP/credits1.png}
(1. Gargoyle (Tomb of Horrors); 2. Tarul Var (Dead in Thay); 3. Mialee (Sunless Citadel); 4. Sir Bluto Sans Pite (White Plume Mountain); 5. Tordek (Forge of Fury); 6. Xipe, the Oni (Hidden Shrine of Tamoachan); 7. Manticore (White Plume Mountain); 8. Kieren, Chosen of Ilmater (Dead in Thay))

\chapter{Magic Items}
\DndItemHeader{Amulet of Protection from Turning}{r, a, r, e None}
While you wear this amulet of silver and turquoise, you have advantage on saving throws against effects that turn undead.

If you fail a saving throw against such an effect, you can choose to succeed instead. You can do so three times, and expended uses recharge daily at dawn. Each time an effect that turns undead is used against you, the amulet glows with silvery blue light for a few seconds.

\DndItemHeader{Balance of Harmony}{u, n, c, o, m, m, o, n}
This scale bears celestial symbols on one pan and fiendish symbols on the other. You can use the scale to cast \textit{detect evil and good} as a ritual. Doing so requires you to place the scale on a solid surface, then sprinkle the pans with holy water or place a transparent gem worth 100 gp in each pan. The scale remains motionless if it detects nothing, tips to one side or the other for good (consecrated) or evil (desecrated), and fluctuates slightly if it detects a creature appropriate to the spell but neither good nor evil. By touching the scales after casting the ritual, you instantly learn any information the spell can normally convey, and then the effect ends.

\DndItemHeader{Blackrazor}{l, e, g, e, n, d, a, r, y None}
Hidden in the dungeon of White Plume Mountain, Blackrazor shines like a piece of night sky filled with stars. Its black scabbard is decorated with pieces of cut obsidian.

You gain a +3 bonus to attack and damage rolls made with this magic weapon. It has the following additional properties.

\subparagraph{Devour Soul}
Whenever you use it to reduce a creature to 0 hit points, the sword slays the creature and devours its soul, unless it is a construct or an undead. A creature whose soul has been devoured by Blackrazor can be restored to life only by a \textit{wish} spell.

When it devours a soul, Blackrazor grants you temporary hit points equal to the slain creature's hit point maximum. These hit points fade after 24 hours. As long as these temporary hit points last and you keep Blackrazor in hand, you have advantage on attack rolls, saving throws, and ability checks.

If you hit an undead with this weapon, you take 1d10 necrotic damage and the target regains 1d10 hit points. If this necrotic damage reduces you to 0 hit points, Blackrazor devours your soul.

\subparagraph{Soul Hunter}
While you hold the weapon. you are aware of the presence of Tiny or larger creatures within 60 feet of you that aren't constructs or undead. You also can't be charmed or frightened.

Blackrazor can cast the \textit{haste} spell on you once per day. It decides when to cast the spell and maintains concentration on it so that you don't have to.

\subparagraph{Sentience}
Blackrazor is a sentient chaotic neutral weapon with an Intelligence of 17, a Wisdom of 10, and a Charisma of 19. It has hearing and darkvision out to a range of 120 feet.

The weapon can speak, read, and understand Common, and can communicate with its wielder telepathically. Its voice is deep and echoing. While you are attuned to it, Blackrazor also understands every language you know.

\subparagraph{Personality}
Blackrazor speaks with an imperious tone, as though accustomed to being obeyed.

The sword's purpose is to consume souls. It doesn't care whose souls it eats, including the wielder's. The sword believes that all matter and energy sprang from a void of negative energy and will one day return to it. Blackrazor is meant to hurry that process along.

Despite its nihilism, Blackrazor feels a strange kinship to \textit{Wave} and \textit{Whelm}, two other weapons locked away under White Plume Mountain. It wants the three weapons to be united again and wielded together in combat, even though it violently disagrees with \textit{Whelm} and finds \textit{Wave} tedious.

Blackrazor's hunger for souls must be regularly fed. If the sword goes three days or more without consuming a soul, a conflict between it and its wielder occurs at the next sunset.

\DndItemHeader{Bracelet of Rock Magic}{v, e, r, y,  , r, a, r, e None}
While you wear this gold bracelet, it grants you immunity to being petrified, and it allows you to cast \textit{flesh to stone} (save DC 15) as an action. Once the spell has been cast three times, the bracelet can no longer cast it. Thereafter, you can cast \textit{stone shape} as an action. After you have done this thirteen times, the bracelet loses its magic and turns from gold to lead.

\subparagraph{Curse}
The bracelet's affinity with earth manifests as an unusual curse. Creatures of flesh that are strongly related to earth and stone, such as stone giants and dwarves, have advantage on the saving throw against \textit{flesh to stone} cast from the bracelet. If such a creature's save is successful, the bracelet breaks your attunement to it and casts the spell on you. You make your saving throw with disadvantage, and on a failed save you are petrified instantly.

\DndItemHeader{Eagle Whistle}{r, a, r, e}
While you blow an eagle whistle continuously, you can fly twice as fast as your walking speed. You can blow the whistle continuously for a number of rounds equal to 5 + five times your Constitution modifier (minimum of 1 round) or until you talk, hold your breath, or start suffocating. A use of the whistle also ends if you land. If you are aloft when you stop blowing the whistle, you fall. The whistle has three uses. It regains expended uses daily at dawn.

\DndItemHeader{Hell Hound Cloak}{r, a, r, e None}
This dark cloak is made of cured \textbf{hell hound} hide. As an action, you can command the cloak to transform you into a \textbf{hell hound} for up to 1 hour. The transformation otherwise functions as the \textit{polymorph} spell, but you can use a bonus action to revert to your normal form.

\subparagraph{Curse}
This cloak is cursed with the essence of a \textbf{hell hound}, and becoming attuned to it extends the curse to you. Until the curse is broken with \textit{remove curse} or similar magic, you are unwilling to part with the cloak, keeping it within reach at all times.

The sixth time you use the cloak, and each time thereafter, you must make a DC 15 Charisma saving throw. On a failed save, the transformation lasts until dispelled or until you drop to 0 hit points, and you can't willingly return to normal form. If you ever remain in \textbf{hell hound} form for 6 hours, the transformation becomes permanent and you lose your sense of self. All your statistics are then replaced by those of a \textbf{hell hound}. Thereafter, only \textit{remove curse} or similar magic allows you to regain your identity and return to normal. If you remain in this permanent form for 6 days, only a \textit{wish} spell can reverse the transformation.

\DndItemHeader{Loadstone}{r, a, r, e}
This stone is a large gem worth 150 gp.

\subparagraph{Curse}
The stone is cursed, but its magical nature is hidden; \textit{detect magic} doesn't detect it. An \textit{identify} spell reveals the stone's true nature. If you use the Dash or Disengage action while the stone is on your person, its curse activates. Until the curse is broken with \textit{remove curse} or similar magic, your speed is reduced by 5 feet, and your maximum load and maximum lift capacities are halved. You also become unwilling to part with the stone.

\DndItemHeader{Mirror of the Past}{r, a, r, e}
The holder of this platinum hand mirror can learn something about the history of a specific object or creature by taking an action to gaze into the mirror and think of the target. Instead of the holder's reflection, the mirror presents scenes from the target's past. Information conveyed is accurate, but it is random and cryptic, and presented in no particular order. Once it is activated, the mirror gives its information for 1 minute or less, then returns to normal. It can't be used again until the next dawn.

\DndItemHeader{Night Caller}{u, n, c, o, m, m, o, n}
This whistle is carved from transparent crystal, and it resembles a tiny dragon curled up like a snail. The name Night Caller is etched on the whistle in Dwarvish runes. If a character succeeds on a DC 20 Intelligence (Arcana or History) check, the character recalls lore that says the \textbf{duergar} made several such whistles for various groups in an age past.

If you blow the whistle in darkness or under the night sky, it allows you to cast the \textit{animate dead} spell. The target can be affected through up to 10 feet of soft earth or similar material, and if it is, it takes 1 minute to claw its way to the surface to serve you. Once the whistle has animated an undead creature, it can't do so again until 7 days have passed.

Once every 24 hours, you can blow the whistle to reassert control over one or two creatures you animated with it.

\DndItemHeader{Potion of Greater Healing}{u, n, c, o, m, m, o, n}
You regain 4d4 + 4 hit points when you drink this potion. The potion's red liquid glimmers when agitated.

\DndItemHeader{Robe of Summer}{r, a, r, e None}
This elegant garment is made from fine cloth in hues of red, orange, and gold. While you wear the robe, you have resistance to cold damage. In addition, you are comfortable as if the temperature were that of a balmy day, so you suffer no ill effects from the weather's temperature extremes.

\DndItemHeader{Shatterspike}{u, n, c, o, m, m, o, n None}
You have a +1 bonus to attack and damage rolls made with this magic weapon. If it hits an object, the hit is automatically a critical hit, and it can deal bludgeoning or slashing damage to the object (your choice). Further, damage from nonmagical sources can't harm the weapon.

\DndItemHeader{Spear of Backbiting}{v, e, r, y,  , r, a, r, e None}
You gain a +2 bonus to attack and damage rolls made with this magic weapon. When you throw it, its normal and long ranges both increase by 30 feet, and it deals one extra die of damage on a hit. After you throw it and it hits or misses, it flies back to your hand immediately.

\subparagraph{Curse}
This weapon is cursed, and becoming attuned to it extends the curse to you. Until the curse is broken with \textit{remove curse} or similar magic, you are unwilling to part with the weapon, keeping it within reach at all times. In addition, you have disadvantage on attack rolls made with weapons other than this one.

Whenever you roll a 1 on an attack roll using this weapon, the weapon bends or flies to hit you in the back. Make a new attack roll with advantage against your own AC. If the result is a hit, you take damage as if you had attacked yourself with the spear.

\DndItemHeader{Stone of Good Luck}{u, n, c, o, m, m, o, n None}
While this polished agate is on your person, you gain a +1 bonus to ability checks and saving throws.

\DndItemHeader{Stone of Ill Luck}{u, n, c, o, m, m, o, n None}
This polished agate appears to be a \textit{stone of good luck} to anyone who tries to identify it, and it confers that item's property while on your person.

\subparagraph{Curse}
This item is cursed. While it is on your person, you take a -2 penalty to ability checks and saving throws. Until the curse is discovered, the DM secretly applies this penalty, assuming you are adding the item's bonus. You are unwilling to part with the stone until the curse is broken with \textit{remove curse} or similar magic.

\DndItemHeader{Wand of Entangle}{u, n, c, o, m, m, o, n None}
This wand has 7 charges. While holding it, you can use an action to expend 1 of its charges to cast the \textit{entangle} spell (save DC 13) from it.

The wand regains 1d6 + 1 expended charges daily at dawn. If you expend the wand's last charge, roll a d20. On a 1, the wand crumbles into ashes and is destroyed.

\DndItemHeader{Wave}{l, e, g, e, n, d, a, r, y None}
Held in the dungeon of White Plume Mountain, this trident is an exquisite weapon engraved with images of waves, shells, and sea creatures. Although you must worship a god of the sea to attune to this weapon, Wave happily accepts new converts.

You gain a +3 bonus to attack and damage rolls made with this magic weapon. If you score a critical hit with it, the target takes extra necrotic damage equal to half its hit point maximum.

The weapon also functions as a \textit{trident of fish command} and a \textit{weapon of warning}. It can confer the benefit of a \textit{cap of water breathing} while you hold it, and you can use it as a \textit{cube of force} by choosing the effect, instead of pressing cube sides to select it.

\subparagraph{Sentience}
Wave is a sentient weapon of neutral alignment, with an Intelligence of 14, a Wisdom of 10, and a Charisma of 18. It has hearing and darkvision out to a range of 120 feet.

The weapon communicates telepathically with its wielder and can speak, read, and understand Aquan. It can also speak with aquatic animals as if using a \textit{speak with animals} spell, using telepathy to involve its wielder in the conversation.

\subparagraph{Personality}
When it grows restless, Wave has a habit of humming tunes that vary from sea chanteys to sacred hymns of the sea gods.

Wave zealously desires to convert mortals to the worship of one or more sea gods, or else to consign the faithless to death. Conflict arises if the wielder fails to further the weapon's objectives in the world. The trident has a nostalgic attachment to the place where it was forged, a desolate island called Thunderforge. A sea god imprisoned a family of storm giants there, and the giants forged Wave in an act of devotion to—or rebellion against—that god.

Wave harbors a secret doubt about its own nature and purpose. For all its devotion to the sea gods, Wave fears that it was intended to bring about a particular sea god's demise. This destiny is something Wave might not be able to avert.

\DndItemHeader{Waythe}{l, e, g, e, n, d, a, r, y None}
Waythe is a unique greatsword most recently in the possession of a high-ranking \textbf{cloud giant} ambassador.

You gain a +1 bonus to attack and damage rolls made with this magic weapon. When you hit a creature of the giant type with it, the giant takes an extra 2d6 slashing damage, and it must succeed on a DC 15 Strength saving throw or fall prone.

The sword also functions as a \textit{wand of enemy detection}. It regains all of its expended charges at dawn and isn't at risk of crumbling if its last charge is used.

\subparagraph{Sentience}
Waythe is a sentient weapon of neutral good alignment, with an Intelligence of 12, a Wisdom of 2, and a Charisma of 14. It has hearing and darkvision out to a range of 120 feet.

The weapon can speak and understand Giant and Common, and it can communicate telepathically with its wielder.

\subparagraph{Personality}
This sword believes in freedom and allowing others to live as they see fit. It is protective of its friends, and wants to be friends with a like-minded wielder. (It takes only 1 minute for a good-aligned character to gain attunement with the sword.) Waythe is courageous to the point of foolhardiness, however, and vocally urges bold action. It is likely to come into conflict with an evil or a timid wielder.

\DndItemHeader{Whelm}{l, e, g, e, n, d, a, r, y None}
Whelm is a powerful warhammer forged by dwarves and lost in the dungeon of White Plume Mountain.

You gain a +3 bonus to attack and damage rolls made with this magic weapon. At dawn the day after you first make an attack roll with Whelm, you develop a fear of being outdoors that persists as long as you remain attuned to the weapon. This causes you to have disadvantage on attack rolls, saving throws, and ability checks while you can see the daytime sky.

\subparagraph{Thrown Weapon}
Whelm has the thrown property, with a normal range of 20 feet and a long range of 60 feet. When you hit with a ranged weapon attack using it, the target takes an extra 1d8 bludgeoning damage, or an extra 2d8 bludgeoning damage if the target is a giant. Each time you throw the weapon, it flies back to your hand after the attack. If you don't have a hand free, the weapon lands at your feet.

\subparagraph{Shock Wave}
You can use an action to strike the ground with Whelm and send a shock wave out from the point of impact. Each creature of your choice on the ground within 60 feet of that point must succeed on a DC 15 Constitution saving throw or become stunned for 1 minute. A creature can repeat the saving throw at the end of each of its turns, ending the effect on itself on a success. Once used, this property can't be used again until the next dawn.

\subparagraph{Supernatural Awareness}
While you are holding the weapon, it alerts you to the location of any secret or concealed doors within 30 feet of you. In addition, you can use an action to cast \textit{detect evil and good} or \textit{locate object} from the weapon. Once you cast either spell, you can't cast it from the weapon again until the next dawn.

\subparagraph{Sentience}
Whelm is a sentient lawful neutral weapon with an Intelligence of 15, a Wisdom of 12, and a Charisma of 15. It has hearing and darkvision out to a range of 120 feet.

The weapon communicates telepathically with its wielder and can speak, read, and understand Dwarvish, Giant, and Goblin. It shouts battle cries in Dwarvish when used in combat.

\subparagraph{Personality}
Whelm's purpose is to slaughter giants and goblinoids. It also seeks to protect dwarves against all enemies. Conflict arises if the wielder fails to destroy goblins and giants or to protect dwarves. Whelm has ties to the dwarf clan that created it, variously called the Dankil or the Mightyhammer clan. It longs to be returned to that clan. It would do anything to protect those dwarves from harm. The hammer also carries a secret shame. Centuries ago, a dwarf named Ctenmiir wielded it valiantly for a time. But Ctenmiir was turned into a \textbf{vampire}. His will was strong enough that he bent Whelm to his evil purposes, even killing members of his own clan.

\DndItemHeader{Armor of Vulnerability (Slashing)}{r, a, r, e None}
While wearing this armor, you have resistance to slashing damage.

\subparagraph{Curse}
This armor is cursed, a fact that is revealed only when an \textit{identify} spell is cast on the armor or you attune to it. Attuning to the armor curses you until you are targeted by the \textit{remove curse} spell or similar magic; removing the armor fails to end the curse. While cursed you have vulnerability to bludgeoning and piercing damage.

\DndItemHeader{Blackrazor}{l, e, g, e, n, d, a, r, y None}
Hidden in the dungeon of White Plume Mountain, Blackrazor shines like a piece of night sky filled with stars. Its black scabbard is decorated with pieces of cut obsidian.

You gain a +3 bonus to attack and damage rolls made with this magic weapon. It has the following additional properties.

\subparagraph{Devour Soul}
Whenever you use it to reduce a creature to 0 hit points, the sword slays the creature and devours its soul, unless it is a construct or an undead. A creature whose soul has been devoured by Blackrazor can be restored to life only by a \textit{wish} spell.

When it devours a soul, Blackrazor grants you temporary hit points equal to the slain creature's hit point maximum. These hit points fade after 24 hours. As long as these temporary hit points last and you keep Blackrazor in hand, you have advantage on attack rolls, saving throws, and ability checks.

If you hit an undead with this weapon, you take 1d10 necrotic damage and the target regains 1d10 hit points. If this necrotic damage reduces you to 0 hit points, Blackrazor devours your soul.

\subparagraph{Soul Hunter}
While you hold the weapon. you are aware of the presence of Tiny or larger creatures within 60 feet of you that aren't constructs or undead. You also can't be charmed or frightened.

Blackrazor can cast the \textit{haste} spell on you once per day. It decides when to cast the spell and maintains concentration on it so that you don't have to.

\subparagraph{Sentience}
Blackrazor is a sentient chaotic neutral weapon with an Intelligence of 17, a Wisdom of 10, and a Charisma of 19. It has hearing and darkvision out to a range of 120 feet.

The weapon can speak, read, and understand Common, and can communicate with its wielder telepathically. Its voice is deep and echoing. While you are attuned to it, Blackrazor also understands every language you know.

\subparagraph{Personality}
Blackrazor speaks with an imperious tone, as though accustomed to being obeyed.

The sword's purpose is to consume souls. It doesn't care whose souls it eats, including the wielder's. The sword believes that all matter and energy sprang from a void of negative energy and will one day return to it. Blackrazor is meant to hurry that process along.

Despite its nihilism, Blackrazor feels a strange kinship to \textit{Wave} and \textit{Whelm}, two other weapons locked away under White Plume Mountain. It wants the three weapons to be united again and wielded together in combat, even though it violently disagrees with \textit{Whelm} and finds \textit{Wave} tedious.

Blackrazor's hunger for souls must be regularly fed. If the sword goes three days or more without consuming a soul, a conflict between it and its wielder occurs at the next sunset.

\DndItemHeader{Boots of Striding and Springing}{u, n, c, o, m, m, o, n None}
While you wear these boots, your walking speed becomes 30 feet, unless your walking speed is higher, and your speed isn't reduced if you are encumbered or wearing heavy armor. In addition, you can jump three times the normal distance, though you can't jump farther than your remaining movement would allow.

\DndItemHeader{Cube of Force}{r, a, r, e None}
This cube is about an inch across. Each face has a distinct marking on it that can be pressed. The cube starts with 36 charges, and it regains 1d20 expended charges daily at dawn.

You can use an action to press one of the cube's faces, expending a number of charges based on the chosen face, as shown in the Cube of Force Faces table. Each face has a different effect. If the cube has insufficient charges remaining, nothing happens. Otherwise, a barrier of invisible force springs into existence, forming a cube 15 feet on a side. The barrier is centered on you, moves with you, and lasts for 1 minute, until you use an action to press the cube's sixth face, or the cube runs out of charges. You can change the barrier's effect by pressing a different face of the cube and expending the requisite number of charges, resetting the duration. If your movement causes the barrier to come into contact with a solid object that can't pass through the cube, you can't move any closer to that object as long as the barrier remains.

\begin{DndTable}[header=Cube of Force Faces]{ccX}

Face & Charges & Effect\\
1 & 1 & Gases, wind, and fog can't pass through the barrier.\\
2 & 2 & Nonliving matter can't pass through the barrier. Walls, floors, and ceilings can pass through at your discretion.\\
3 & 3 & Living matter can't pass through the barrier.\\
4 & 4 & Spell effects can't pass through the barrier.\\
5 & 5 & Nothing can pass through the barrier. Walls, floors, and ceilings can pass through at your discretion.\\
6 & 0 & The barrier deactivates.\\
\end{DndTable}

The cube loses charges when the barrier is targeted by certain spells or comes into contact with certain spell or magic item effects, as shown in the table below.


\begin{DndTable}{XX}

Spell or item & Charges Lost\\
\textit{Disintegrate} & 1d12\\
\textit{Horn of blasting} & 1d10\\
\textit{Passwall} & 1d6\\
\textit{Prismatic spray} & 1d20\\
\textit{Wall of fire} & 1d4\\
\end{DndTable}

\DndItemHeader{Goggles of Night}{u, n, c, o, m, m, o, n}
While wearing these dark lenses, you have darkvision out to a range of 60 feet. If you already have darkvision, wearing the goggles increases its range by 60 feet.

\DndItemHeader{Potion of Flying}{v, e, r, y,  , r, a, r, e}
When you drink this potion, you gain a flying speed equal to your walking speed for 1 hour and can hover. If you're in the air when the potion wears off, you fall unless you have some other means of staying aloft. This potion's clear liquid floats at the top of its container and has cloudy white impurities drifting in it.

\DndItemHeader{Potion of Mind Reading}{r, a, r, e}
When you drink this potion, you gain the effect of the \textit{detect thoughts} spell (save DC 13). The potion's dense, purple liquid has an ovoid cloud of pink floating in it.

\DndItemHeader{Ring of Protection}{r, a, r, e None}
You gain a +1 bonus to AC and saving throws while wearing this ring.

\DndItemHeader{Ring of Spell Storing}{r, a, r, e None}
This ring stores spells cast into it, holding them until the attuned wearer uses them. The ring can store up to 5 levels worth of spells at a time. When found, it contains 1d6 - 1 levels of stored spells chosen by the DM.

Any creature can cast a spell of 1st through 5th level into the ring by touching the ring as the spell is cast. The spell has no effect, other than to be stored in the ring. If the ring can't hold the spell, the spell is expended without effect. The level of the slot used to cast the spell determines how much space it uses.

While wearing this ring, you can cast any spell stored in it. The spell uses the slot level, spell save DC, spell attack bonus, and spellcasting ability of the original caster, but is otherwise treated as if you cast the spell. The spell cast from the ring is no longer stored in it, freeing up space.

\DndItemHeader{Ring of Spell Turning}{l, e, g, e, n, d, a, r, y None}
While wearing this ring, you have advantage on saving throws against any spell that targets only you (not in an area of effect). In addition, if you roll a 20 for the save and the spell is 7th level or lower, the spell has no effect on you and instead targets the caster, using the slot level, spell save DC, attack bonus, and spellcasting ability of the caster.

\DndItemHeader{Ring of X-ray Vision}{r, a, r, e None}
While wearing this ring, you can use an action to speak its command word. When you do so, you can see into and through solid matter for 1 minute. This vision has a radius of 30 feet. To you, solid objects within that radius appear transparent and don't prevent light from passing through them. The vision can penetrate 1 foot of stone, 1 inch of common metal, or up to 3 feet of wood or dirt. Thicker substances block the vision, as does a thin sheet of lead.

Whenever you use the ring again before taking a long rest, you must succeed on a DC 15 Constitution saving throw or gain one level of exhaustion.

\DndItemHeader{Wave}{l, e, g, e, n, d, a, r, y None}
Held in the dungeon of White Plume Mountain, this trident is an exquisite weapon engraved with images of waves, shells, and sea creatures. Although you must worship a god of the sea to attune to this weapon, Wave happily accepts new converts.

You gain a +3 bonus to attack and damage rolls made with this magic weapon. If you score a critical hit with it, the target takes extra necrotic damage equal to half its hit point maximum.

The weapon also functions as a \textit{trident of fish command} and a \textit{weapon of warning}. It can confer the benefit of a \textit{cap of water breathing} while you hold it, and you can use it as a \textit{cube of force} by choosing the effect, instead of pressing cube sides to select it.

\subparagraph{Sentience}
Wave is a sentient weapon of neutral alignment, with an Intelligence of 14, a Wisdom of 10, and a Charisma of 18. It has hearing and darkvision out to a range of 120 feet.

The weapon communicates telepathically with its wielder and can speak, read, and understand Aquan. It can also speak with aquatic animals as if using a \textit{speak with animals} spell, using telepathy to involve its wielder in the conversation.

\subparagraph{Personality}
When it grows restless, Wave has a habit of humming tunes that vary from sea chanteys to sacred hymns of the sea gods.

Wave zealously desires to convert mortals to the worship of one or more sea gods, or else to consign the faithless to death. Conflict arises if the wielder fails to further the weapon's objectives in the world. The trident has a nostalgic attachment to the place where it was forged, a desolate island called Thunderforge. A sea god imprisoned a family of storm giants there, and the giants forged Wave in an act of devotion to—or rebellion against—that god.

Wave harbors a secret doubt about its own nature and purpose. For all its devotion to the sea gods, Wave fears that it was intended to bring about a particular sea god's demise. This destiny is something Wave might not be able to avert.

\DndItemHeader{Whelm}{l, e, g, e, n, d, a, r, y None}
Whelm is a powerful warhammer forged by dwarves and lost in the dungeon of White Plume Mountain.

You gain a +3 bonus to attack and damage rolls made with this magic weapon. At dawn the day after you first make an attack roll with Whelm, you develop a fear of being outdoors that persists as long as you remain attuned to the weapon. This causes you to have disadvantage on attack rolls, saving throws, and ability checks while you can see the daytime sky.

\subparagraph{Thrown Weapon}
Whelm has the thrown property, with a normal range of 20 feet and a long range of 60 feet. When you hit with a ranged weapon attack using it, the target takes an extra 1d8 bludgeoning damage, or an extra 2d8 bludgeoning damage if the target is a giant. Each time you throw the weapon, it flies back to your hand after the attack. If you don't have a hand free, the weapon lands at your feet.

\subparagraph{Shock Wave}
You can use an action to strike the ground with Whelm and send a shock wave out from the point of impact. Each creature of your choice on the ground within 60 feet of that point must succeed on a DC 15 Constitution saving throw or become stunned for 1 minute. A creature can repeat the saving throw at the end of each of its turns, ending the effect on itself on a success. Once used, this property can't be used again until the next dawn.

\subparagraph{Supernatural Awareness}
While you are holding the weapon, it alerts you to the location of any secret or concealed doors within 30 feet of you. In addition, you can use an action to cast \textit{detect evil and good} or \textit{locate object} from the weapon. Once you cast either spell, you can't cast it from the weapon again until the next dawn.

\subparagraph{Sentience}
Whelm is a sentient lawful neutral weapon with an Intelligence of 15, a Wisdom of 12, and a Charisma of 15. It has hearing and darkvision out to a range of 120 feet.

The weapon communicates telepathically with its wielder and can speak, read, and understand Dwarvish, Giant, and Goblin. It shouts battle cries in Dwarvish when used in combat.

\subparagraph{Personality}
Whelm's purpose is to slaughter giants and goblinoids. It also seeks to protect dwarves against all enemies. Conflict arises if the wielder fails to destroy goblins and giants or to protect dwarves. Whelm has ties to the dwarf clan that created it, variously called the Dankil or the Mightyhammer clan. It longs to be returned to that clan. It would do anything to protect those dwarves from harm. The hammer also carries a secret shame. Centuries ago, a dwarf named Ctenmiir wielded it valiantly for a time. But Ctenmiir was turned into a \textbf{vampire}. His will was strong enough that he bent Whelm to his evil purposes, even killing members of his own clan.

\chapter{Creatures}
\setlist[description]{nolistsep,listparindent=3pt,topsep=6pt,font={\bfseries\sffamily\small}}
\pagestyle{empty}

\begin{DndMonster}[float=t]{Air Elemental}
\DndMonsterType{Large elemental, neutral}
\DndMonsterBasics[
armor-class = {15},
hit-points = {\DndDice{12d10 + 24}},
speed = {0 ft., fly 90 ft. \textit{(hover)}}
]
\DndMonsterAbilityScores[
 str = 14,
 dex = 20,
 con = 14,
 int = 6,
 wis = 10,
 cha = 6
]
\DndMonsterDetails[
damage-resistances = {lightning; thunder; bludgeoning, piercing, slashing from nonmagical attacks},
damage-immunities = {poison},
condition-immunities = {exhaustion, grappled, paralyzed, petrified, poisoned, prone, restrained, unconscious},
senses = {darkvision 60 ft., passive perception 10},
languages = {Auran},
challenge = 5,
]
\DndMonsterAction{Air Form}
The elemental can enter a hostile creature's space and stop there. It can move through a space as narrow as 1 inch wide without squeezing.

\DndMonsterSection{Actions}
\DndMonsterAction{Multiattack}
The elemental makes two slam attacks.

\DndMonsterAction{Slam}
\textsl{Melee Weapon Attack:} 8 to hit, reach 5 ft., one target. \textsl{Hit:} 14 (2d8 + 5) bludgeoning damage.

\DndMonsterAction{Whirlwind (Recharge 4\textendash 6)}
Each creature in the elemental's space must make a DC 13 Strength saving throw. On a failure, a target takes 15 (3d8 + 2) bludgeoning damage and is flung up 20 feet away from the elemental in a random direction and knocked prone. If a thrown target strikes an object, such as a wall or floor, the target takes 3 (1d6) bludgeoning damage for every 10 feet it was thrown. If the target is thrown at another creature, that creature must succeed on a DC 13 Dexterity saving throw or take the same damage and be knocked prone.

If the saving throw is successful, the target takes half the bludgeoning damage and isn't flung away or knocked prone.

\end{DndMonster}



\begin{DndMonster}[float=t]{Animated Table}
\DndMonsterType{Large construct, unaligned}
\DndMonsterBasics[
armor-class = {15 (natural armor)},
hit-points = {\DndDice{6d10 + 6}},
speed = {40 ft.}
]
\DndMonsterAbilityScores[
 str = 18,
 dex = 8,
 con = 13,
 int = 1,
 wis = 3,
 cha = 1
]
\DndMonsterDetails[
damage-immunities = {poison, psychic},
condition-immunities = {blinded, charmed, deafened, exhaustion, frightened, paralyzed, petrified, poisoned},
senses = {blindsight 60 ft. (blind beyond this radius), passive perception 6},
challenge = 2,
]
\DndMonsterAction{Antimagic Susceptibility}
The table is incapacitated while in the area of an \textit{antimagic field}. If targeted by \textit{dispel magic}, the table must succeed on a Constitution saving throw against the caster's spell save DC or fall unconscious for 1 minute.

\DndMonsterAction{False Appearance}
While the table remains motionless, it is indistinguishable from a normal table.

\DndMonsterAction{Charge}
If the table moves at least 20 feet straight toward a target and then hits it with a ram attack on the same turn, the target takes an extra 9 (2d8) bludgeoning damage. If the target is a creature, it must succeed on a DC 15 Strength saving throw or be knocked prone.

\DndMonsterSection{Actions}
\DndMonsterAction{Ram}
\textsl{Melee Weapon Attack:} 6, reach 5 ft., one target. \textsl{Hit:} 13 (2d8 + 4) bludgeoning damage.

\end{DndMonster}



\begin{DndMonster}[float=t]{Barghest}
\DndMonsterType{Large fiend (shapechanger), neutral evil}
\DndMonsterBasics[
armor-class = {17 (natural armor)},
hit-points = {\DndDice{12d10 + 24}},
speed = {60 ft. (30 ft. in goblin form)}
]
\DndMonsterAbilityScores[
 str = 19,
 dex = 15,
 con = 14,
 int = 13,
 wis = 12,
 cha = 14
]
\DndMonsterDetails[
skills = {Deception +4, Intimidation +4, Perception +5, Stealth +4},
damage-resistances = {cold; fire; lightning; bludgeoning, piercing, slashing from nonmagical attacks},
damage-immunities = {acid, poison},
condition-immunities = {poisoned},
senses = {blindsight 60 ft., darkvision 60 ft., passive perception 15},
languages = {Abyssal, Common, Goblin, Infernal, telepathy 60 ft.},
challenge = 4,
]
\DndMonsterAction{Shapechanger}
The barghest can use its action to polymorph into a Small goblin or back into its true form. Other than its size and speed, its statistics are the same in each form. Any equipment it is wearing or carrying isn't transformed. The barghest reverts to its true form if it dies.

\DndMonsterAction{Fire Banishment}
When the barghest starts its turn engulfed in flames that are at least 10 feet high or wide, it must succeed on a DC 15 Charisma saving throw or be instantly banished to Gehenna. Instantaneous bursts of flame (such as a red dragon's breath or a \textit{fireball} spell) don't have this effect on the barghest.

\DndMonsterAction{Keen Smell}
The barghest has advantage on Wisdom (Perception) checks that rely on smell.

\DndMonsterAction{Soul Feeding}
A barghest can feed on the corpse of a humanoid that it killed that has been dead for less than 10 minutes, devouring both flesh and soul in doing so. This feeding takes at least 1 minute, and it destroys the victim's body. The victim's soul is trapped in the barghest for 24 hours, after which time it is digested. If the barghest dies before the soul is digested, the soul is released.

While a humanoid's soul is trapped in a barghest, any form of revival that could work has only a 50 percent chance of doing so, freeing the soul from the barghest if it is successful. Once a creature's soul is digested, however, no mortal magic can return that humanoid to life.

\DndMonsterAction{Innate Spellcasting}
The barghest's innate spellcasting ability is Charisma (spell save DC 12). It can innately cast the following spells, requiring no material components:
\begin{DndMonsterSpells}
  \DndInnateSpellLevel{\textit{levitate}, \textit{minor illusion}, \textit{pass without trace}}
  \DndInnateSpellLevel[1e]{\textit{charm person}, \textit{dimension door}, \textit{suggestion}}
\end{DndMonsterSpells}
\DndMonsterSection{Actions}
\DndMonsterAction{Bite}
\textsl{Melee Weapon Attack:} 6 to hit, reach 5 ft., one target. \textsl{Hit:} 13 (2d8 + 4) piercing damage.

\DndMonsterAction{Claws}
\textsl{Melee Weapon Attack:} 6 to hit, reach 5 ft., one target. \textsl{Hit:} 8 (1d8 + 4) slashing damage.

\end{DndMonster}



\begin{DndMonster}[float=t]{Black Pudding}
\DndMonsterType{Large ooze, unaligned}
\DndMonsterBasics[
armor-class = {7},
hit-points = {\DndDice{10d10 + 30}},
speed = {20 ft., climb 20 ft.}
]
\DndMonsterAbilityScores[
 str = 16,
 dex = 5,
 con = 16,
 int = 1,
 wis = 6,
 cha = 1
]
\DndMonsterDetails[
damage-immunities = {acid, cold, lightning, slashing},
condition-immunities = {blinded, charmed, deafened, exhaustion, frightened, prone},
senses = {blindsight 60 ft. (blind beyond this radius), passive perception 8},
challenge = 4,
]
\DndMonsterAction{Amorphous}
The pudding can move through a space as narrow as 1 inch wide without squeezing.

\DndMonsterAction{Corrosive Form}
A creature that touches the pudding or hits it with a melee attack while within 5 feet of it takes 4 (1d8) acid damage. Any nonmagical weapon made of metal or wood that hits the pudding corrodes. After dealing damage, the weapon takes a permanent and cumulative −1 penalty to damage rolls. If its penalty drops to −5, the weapon is destroyed. Nonmagical ammunition made of metal or wood that hits the pudding is destroyed after dealing damage. The pudding can eat through 2-inch-thick, nonmagical wood or metal in 1 round.

\DndMonsterAction{Spider Climb}
The pudding can climb difficult surfaces, including upside down on ceilings, without needing to make an ability check.

\DndMonsterSection{Actions}
\DndMonsterAction{Pseudopod}
\textsl{Melee Weapon Attack:} 5 to hit, reach 5 ft., one target. \textsl{Hit:} 6 (1d6 + 3) bludgeoning damage plus 18 (4d8) acid damage. In addition, nonmagical armor worn by the target is partly dissolved and takes a permanent and cumulative −1 penalty to the AC it offers. The armor is destroyed if the penalty reduces its AC to 10.

\DndMonsterSection{Reactions}
\DndMonsterAction{Split}
When a pudding that is Medium or larger is subjected to lightning or slashing damage, it splits into two new puddings if it has at least 10 hit points. Each new pudding has hit points equal to half the original pudding's, rounded down. New puddings are one size smaller than the original pudding.

\end{DndMonster}



\begin{DndMonster}[float=t]{Bugbear}
\DndMonsterType{Medium humanoid (goblinoid), chaotic evil}
\DndMonsterBasics[
armor-class = {16 (\textit{hide armor})},
hit-points = {\DndDice{5d8 + 5}},
speed = {30 ft.}
]
\DndMonsterAbilityScores[
 str = 15,
 dex = 14,
 con = 13,
 int = 8,
 wis = 11,
 cha = 9
]
\DndMonsterDetails[
skills = {Stealth +6, Survival +2},
senses = {darkvision 60 ft., passive perception 10},
languages = {Common, Goblin},
challenge = 1,
]
\DndMonsterAction{Brute}
A melee weapon deals one extra die of its damage when the bugbear hits with it (included in the attack).

\DndMonsterAction{Surprise Attack}
If the bugbear surprises a creature and hits it with an attack during the first round of combat, the target takes an extra 7 (2d6) damage from the attack.

\DndMonsterSection{Actions}
\DndMonsterAction{Morningstar}
\textsl{Melee Weapon Attack:} 4 to hit, reach 5 ft., one target. \textsl{Hit:} 11 (2d8 + 2) piercing damage.

\DndMonsterAction{Javelin}
\textsl{Melee or Ranged Weapon Attack:} 4 to hit, reach 5 ft. or range 30/120 ft., one target. \textsl{Hit:} 9 (2d6 + 2) piercing damage in melee or 5 (1d6 + 2) piercing damage at range.

\end{DndMonster}



\begin{DndMonster}[float=t]{Centaur Mummy}
\DndMonsterType{Large undead, lawful evil}
\DndMonsterBasics[
armor-class = {13 (natural armor)},
hit-points = {\DndDice{10d10 + 30}},
speed = {30 ft.}
]
\DndMonsterAbilityScores[
 str = 20,
 dex = 12,
 con = 16,
 int = 5,
 wis = 14,
 cha = 12
]
\DndMonsterDetails[
saving-throws = {Wis +5},
damage-vulnerabilities = {fire},
damage-resistances = {bludgeoning, piercing, slashing from nonmagical attacks},
damage-immunities = {necrotic, poison},
condition-immunities = {charmed, exhaustion, frightened, paralyzed, poisoned},
senses = {darkvision 60 ft., passive perception 12},
languages = {Common, Sylvan},
challenge = 6,
]
\DndMonsterAction{Charge}
If the centaur mummy moves at least 20 feet straight toward a target and then hits it with a pike attack on the same turn, the target takes an extra 10 (3d6) piercing damage.

\DndMonsterSection{Actions}
\DndMonsterAction{Multiattack}
The centaur mummy makes two melee attacks, one with its pike and one with its hooves, or it attacks with its pike and uses Dreadful Glare.

\DndMonsterAction{Pike}
\textsl{Melee Weapon Attack:} 8 to hit, reach 10 ft., one target. \textsl{Hit:} 10 (1d10 + 5) piercing damage.

\DndMonsterAction{Hooves}
\textsl{Melee Weapon Attack:} 8 to hit, reach 5 ft., one target. \textsl{Hit:} 12 (2d6 + 5) piercing damage plus 10 (3d6) necrotic damage. If the target is a creature, it must succeed on a DC 14 Constitution saving throw or be cursed with mummy rot. The cursed target can't regain hit points, and its hit point maximum decreases by 10 (3d6) for every 24 hours that elapse. If the curse reduces the target's hit point maximum to 0, the target dies, and its body turns to dust. The curse lasts until removed by the \textit{remove curse} spell or similar magic.

\DndMonsterAction{Dreadful Glare}
The centaur mummy targets one creature it can see within 60 feet of it. If the target can see the mummy, the target must succeed on a DC 12 Wisdom saving throw against this magic or become frightened until the end of the mummy's next turn. If the target fails the saving throw by 5 or more, it is also paralyzed for the same duration. A target that succeeds on the saving throw is immune to the Dreadful Glare of all mummies (but not mummy lords) for the next 24 hours.

\end{DndMonster}



\begin{DndMonster}[float=t]{Champion}
\DndMonsterType{Medium humanoid (any race), any alignment}
\DndMonsterBasics[
armor-class = {18 (\textit{plate armor})},
hit-points = {\DndDice{22d8 + 44}},
speed = {30 ft.}
]
\DndMonsterAbilityScores[
 str = 20,
 dex = 15,
 con = 14,
 int = 10,
 wis = 14,
 cha = 12
]
\DndMonsterDetails[
saving-throws = {Str +9, Con +6},
skills = {Athletics +9, Intimidation +5, Perception +6},
languages = {any one language (usually Common)},
challenge = 9,
]
\DndMonsterAction{Indomitable (2/Day)}
The champion rerolls a failed saving throw.

\DndMonsterAction{Second Wind (Recharges after a Short or Long Rest)}
As a bonus action, the champion can regain 20 hit points.

\DndMonsterSection{Actions}
\DndMonsterAction{Multiattack}
The champion makes three attacks with its greatsword or its shortbow.

\DndMonsterAction{Greatsword}
\textsl{Melee Weapon Attack:} 9 to hit, reach 5 ft., one target. \textsl{Hit:} 12 (2d6 + 5) slashing damage, plus 7 (2d6) slashing damage if the champion has more than half of its total hit points remaining.

\DndMonsterAction{Shortbow}
\textsl{Ranged Weapon Attack:} 6 to hit, range 80/320 ft., one target. \textsl{Hit:} 5 (1d6 + 2) piercing damage, plus 7 (2d6) piercing damage if the champion has more than half of its total hit points remaining.

\end{DndMonster}



\begin{DndMonster}[float=t]{Choker}
\DndMonsterType{Small aberration, chaotic evil}
\DndMonsterBasics[
armor-class = {16 (natural armor)},
hit-points = {\DndDice{3d6 + 3}},
speed = {30 ft.}
]
\DndMonsterAbilityScores[
 str = 16,
 dex = 14,
 con = 13,
 int = 4,
 wis = 12,
 cha = 7
]
\DndMonsterDetails[
skills = {Stealth +6},
senses = {darkvision 60 ft., passive perception 11},
languages = {Deep Speech},
challenge = 1,
]
\DndMonsterAction{Aberrant Quickness (Recharges after a Short or Long Rest)}
The choker can take an extra action on its turn.

\DndMonsterAction{Boneless}
The choker can move through and occupy a space as narrow as 4 inches wide without squeezing.

\DndMonsterAction{Spider Climb}
The choker can climb difficult surfaces, including upside down on ceilings, without needing to make an ability check.

\DndMonsterSection{Actions}
\DndMonsterAction{Multiattack}
The choker makes two tentacle attacks.

\DndMonsterAction{Tentacle}
\textsl{Melee Weapon Attack:} 5 to hit, reach 10 ft., one target. \textsl{Hit:} 5 (1d4 + 3) bludgeoning damage plus 3 (1d6) piercing damage. If the target is a Large or smaller creature, it is grappled (escape DC 15). Until this grapple ends, the target is restrained, and the choker can't use this tentacle on another target. The choker has two tentacles. If this attack is a critical hit, the target also can't breathe or speak until the grapple ends.

\end{DndMonster}



\begin{DndMonster}[float=t]{Cloud Giant}
\DndMonsterType{Huge giant, neutral good (50\%) neutral evil (50\%)}
\DndMonsterBasics[
armor-class = {14 (natural armor)},
hit-points = {\DndDice{16d12 + 96}},
speed = {40 ft.}
]
\DndMonsterAbilityScores[
 str = 27,
 dex = 10,
 con = 22,
 int = 12,
 wis = 16,
 cha = 16
]
\DndMonsterDetails[
saving-throws = {Con +10, Wis +7, Cha +7},
skills = {Insight +7, Perception +7},
languages = {Common, Giant},
challenge = 9,
]
\DndMonsterAction{Keen Smell}
The giant has advantage on Wisdom (Perception) checks that rely on smell.

\DndMonsterAction{Innate Spellcasting}
The giant's innate spellcasting ability is Charisma. It can innately cast the following spells, requiring no material components:
\begin{DndMonsterSpells}
  \DndInnateSpellLevel{\textit{detect magic}, \textit{fog cloud}, \textit{light}}
  \DndInnateSpellLevel[3e]{\textit{feather fall}, \textit{fly}, \textit{misty step}, \textit{telekinesis}}
  \DndInnateSpellLevel[1e]{\textit{control weather}, \textit{gaseous form}}
\end{DndMonsterSpells}
\DndMonsterSection{Actions}
\DndMonsterAction{Multiattack}
The giant makes two morningstar attacks.

\DndMonsterAction{Morningstar}
\textsl{Melee Weapon Attack:} 12 to hit, reach 10 ft., one target. \textsl{Hit:} 21 (3d8 + 8) piercing damage.

\DndMonsterAction{Rock}
\textsl{Ranged Weapon Attack:} 12 to hit, range 60/240 ft., one target. \textsl{Hit:} 30 (4d10 + 8) bludgeoning damage.

\end{DndMonster}



\begin{DndMonster}[float=t]{Conjurer}
\DndMonsterType{Medium humanoid (any race), any alignment}
\DndMonsterBasics[
armor-class = {12 (15 with \textit{mage armor})},
hit-points = {\DndDice{9d8}},
speed = {30 ft.}
]
\DndMonsterAbilityScores[
 str = 9,
 dex = 14,
 con = 11,
 int = 17,
 wis = 12,
 cha = 11
]
\DndMonsterDetails[
saving-throws = {Int +6, Wis +4},
skills = {Arcana +6, History +6},
languages = {any four languages},
challenge = 6,
]
\DndMonsterAction{Benign Transportation (Recharges after the Conjurer Casts a Conjuration Spell of 1st Level or Higher)}
As a bonus action, the conjurer teleports up to 30 feet to an unoccupied space that it can see. If it instead chooses a space within range that is occupied by a willing Small or Medium creature, they both teleport, swapping places.

\DndMonsterAction{Spellcasting}
The conjurer is a 9th-level spellcaster. Its spellcasting ability is intelligence (spell save DC 14, 6 to hit with spell attacks). The conjurer has the following wizard spells prepared:
\begin{DndMonsterSpells}
  \DndMonsterSpellLevel{\textit{acid splash}, \textit{mage hand}, \textit{poison spray}, \textit{prestidigitation}}
   \DndMonsterSpellLevel[1][4]{\textit{mage armor}, \textit{magic missile}, \textit{unseen servant}*}
   \DndMonsterSpellLevel[2][3]{\textit{cloud of daggers}*, \textit{misty step}*, \textit{web}*}
   \DndMonsterSpellLevel[3][3]{\textit{fireball}, \textit{stinking cloud}*}
   \DndMonsterSpellLevel[4][3]{\textit{Evard's black tentacles}*, \textit{stoneskin}}
   \DndMonsterSpellLevel[5][2]{\textit{cloudkill}*, \textit{conjure elemental}*}
\end{DndMonsterSpells}
\DndMonsterSection{Actions}
\DndMonsterAction{Dagger}
\textsl{Melee or Ranged Weapon Attack:} 5 to hit, reach 5 ft. or range 20/60 ft., one target. \textsl{Hit:} 4 (1d4 + 2) piercing damage.

\end{DndMonster}



\begin{DndMonster}[float=t]{Deathlock Wight}
\DndMonsterType{Medium undead, lawful evil}
\DndMonsterBasics[
armor-class = {12 (15 with \textit{mage armor})},
hit-points = {\DndDice{5d8 + 15}},
speed = {30 ft.}
]
\DndMonsterAbilityScores[
 str = 11,
 dex = 14,
 con = 16,
 int = 12,
 wis = 14,
 cha = 16
]
\DndMonsterDetails[
saving-throws = {Wis +4},
skills = {Arcana +3, Perception +4},
damage-resistances = {necrotic; bludgeoning, piercing, slashing from nonmagical attacks},
damage-immunities = {poison},
condition-immunities = {exhaustion, poisoned},
senses = {darkvision 60 ft., passive perception 14},
languages = {the languages it knew in life},
challenge = 3,
]
\DndMonsterAction{Sunlight Sensitivity}
While in sunlight, the wight has disadvantage on attack rolls, as well as on Wisdom (Perception) checks that rely on sight.

\DndMonsterAction{Innate Spellcasting}
The wight's innate spellcasting ability is Charisma (spell save DC 13). It can innately cast the following spells, requiring no verbal or material components:
\begin{DndMonsterSpells}
  \DndInnateSpellLevel{\textit{detect magic}, \textit{disguise self}, \textit{mage armor}}
  \DndInnateSpellLevel[1e]{\textit{fear}, \textit{hold person}, \textit{misty step}}
\end{DndMonsterSpells}
\DndMonsterSection{Actions}
\DndMonsterAction{Multiattack}
The wight attacks twice with Grave Bolt.

\DndMonsterAction{Grave Bolt}
\textsl{Ranged Spell Attack:} 5 to hit, range 120 ft., one target. \textsl{Hit:} 7 (1d8 + 3) necrotic damage.

\DndMonsterAction{Life Drain}
\textsl{Melee Weapon Attack:} 4 to hit, reach 5 ft., one creature. \textsl{Hit:} 9 (2d6 + 2) necrotic damage. The target must succeed on a DC 13 Constitution saving throw or its hit point maximum is reduced by an amount equal to the damage taken. This reduction lasts until the target finishes a long rest. The target dies if this effect reduces its hit point maximum to 0.

A humanoid slain by this attack rises 24 hours later as a \textbf{zombie} under the wight's control, unless the humanoid is restored to life or its body is destroyed. The wight can have no more than twelve zombies under its control at one time.

\end{DndMonster}



\begin{DndMonster}[float=t]{Dread Warrior}
\DndMonsterType{Medium undead, neutral evil}
\DndMonsterBasics[
armor-class = {18 (\textit{chain mail})},
hit-points = {\DndDice{5d8 + 15}},
speed = {30 ft.}
]
\DndMonsterAbilityScores[
 str = 15,
 dex = 11,
 con = 16,
 int = 10,
 wis = 12,
 cha = 10
]
\DndMonsterDetails[
saving-throws = {Wis +3},
skills = {Athletics +4, Perception +3},
damage-immunities = {poison},
condition-immunities = {exhaustion, poisoned},
senses = {darkvision 60 ft., passive perception 13},
languages = {Common},
challenge = 1,
]
\DndMonsterAction{Undead Fortitude}
If damage reduces the dread warrior to 0 hit points, it must make a Constitution saving throw with a DC of 5+the damage taken, unless the damage is radiant or from a critical hit. On a success, the dread warrior drops to 1 hit point instead.

\DndMonsterSection{Actions}
\DndMonsterAction{Multiattack}
The dread warrior makes two melee attacks.

\DndMonsterAction{Battleaxe}
\textsl{Melee Weapon Attack:} 4 to hit, reach 5 ft., one target. \textsl{Hit:} 6 (1d8 + 2) slashing damage, or 7 (1d10 + 2) slashing damage if wielded with two hands.

\DndMonsterAction{Javelin}
\textsl{Melee or Ranged Weapon Attack:} 4 to hit, reach 5 ft. or range 30/120 ft., one target. \textsl{Hit:} 5 (1d6 + 2) piercing damage.

\end{DndMonster}



\begin{DndMonster}[float=t]{Duergar}
\DndMonsterType{Medium humanoid (dwarf), lawful evil}
\DndMonsterBasics[
armor-class = {16 (\textit{scale mail})},
hit-points = {\DndDice{4d8 + 4}},
speed = {25 ft.}
]
\DndMonsterAbilityScores[
 str = 14,
 dex = 11,
 con = 14,
 int = 11,
 wis = 10,
 cha = 9
]
\DndMonsterDetails[
damage-resistances = {poison},
senses = {darkvision 120 ft., passive perception 10},
languages = {Dwarvish, Undercommon},
challenge = 1,
]
\DndMonsterAction{Duergar Resilience}
The duergar has advantage on saving throws against poison, spells, and illusions, as well as to resist being charmed or paralyzed.

\DndMonsterAction{Sunlight Sensitivity}
While in sunlight, the duergar has disadvantage on attack rolls, as well as on Wisdom (Perception) checks that rely on sight.

\DndMonsterSection{Actions}
\DndMonsterAction{Enlarge (Recharges after a Short or Long Rest)}
For 1 minute, the duergar magically increases in size, along with anything it is wearing or carrying. While enlarged, the duergar is Large, doubles its damage dice on Strength-based weapon attacks (included in the attacks), and makes Strength checks and Strength saving throws with advantage. If the duergar lacks the room to become Large, it attains the maximum size possible in the space available.

\DndMonsterAction{War Pick}
\textsl{Melee Weapon Attack:} 4 to hit, reach 5 ft., one target. \textsl{Hit:} 6 (1d8 + 2) piercing damage, or 11 (2d8 + 2) piercing damage while enlarged.

\DndMonsterAction{Javelin}
\textsl{Melee or Ranged Weapon Attack:} 4 to hit, reach 5 ft. or range 30/120 ft., one target. \textsl{Hit:} 5 (1d6 + 2) piercing damage, or 9 (2d6 + 2) piercing damage while enlarged.

\DndMonsterAction{Invisibility (Recharges after a Short or Long Rest)}
The duergar magically turns invisible until it attacks, casts a spell, or uses its Enlarge, or until its concentration is broken, up to 1 hour (as if concentration on a spell). Any equipment the duergar wears or carries is invisible with it.

\end{DndMonster}



\begin{DndMonster}[float=t]{Duergar Spy}
\DndMonsterType{Medium humanoid (dwarf), lawful evil}
\DndMonsterBasics[
armor-class = {15 (studded leather armor)},
hit-points = {\DndDice{6d8 + 6}},
speed = {25 ft.}
]
\DndMonsterAbilityScores[
 str = 10,
 dex = 16,
 con = 12,
 int = 12,
 wis = 10,
 cha = 13
]
\DndMonsterDetails[
skills = {Deception +5, Insight +2, Investigation +5, Perception +4, Persuasion +3, Sleight Of Hand +5, Stealth +7},
damage-resistances = {poison},
senses = {darkvision 120 ft., passive perception 14},
languages = {Dwarvish, Undercommon},
challenge = 2,
]
\DndMonsterAction{Cunning Action}
On each of its turns, the spy can use a bonus action to take the Dash, Disengage, or Hide action.

\DndMonsterAction{Duergar Resilience}
The spy has advantage on saving throws against poison, spells, and illusions, as well as to resist being charmed or paralyzed.

\DndMonsterAction{Sneak Attack}
Once per turn, the spy can deal an extra 7 (2d6) damage when it hits a target with a weapon attack and has advantage on the attack roll, or when the target is within 5 feet of an ally of the spy that isn't incapacitated and the spy doesn't have disadvantage on the attack roll.

\DndMonsterAction{Sunlight Sensitivity}
While in sunlight, the spy has disadvantage on attack rolls, as well as on Wisdom (Perception) checks that rely on sight.

\DndMonsterSection{Actions}
\DndMonsterAction{Multiattack}
The spy makes two shortsword attacks.

\DndMonsterAction{Enlarge (Recharges after a Short or Long Rest)}
For 1 minute, the spy magically increases in size, along with anything it is wearing or carrying. While enlarged, the spy is Large, doubles her damage dice on Strength-based weapon attacks (included in the attacks), and makes Strength checks and Strength saving throws with advantage. If the spy lacks the room to become Large, it attains the maximum size possible in the space available.

\DndMonsterAction{Shortsword}
\textsl{Melee Weapon Attack:} 5 to hit, reach 5 ft., one target. \textsl{Hit:} 6 (1d6 + 3) piercing damage, or 10 (2d6 + 3) piercing damage while enlarged.

\DndMonsterAction{Hand Crossbow}
\textsl{Ranged Weapon Attack:} 5 to hit, range 30/120 ft., one target. \textsl{Hit:} 6 (1d6 + 3) piercing damage.

\DndMonsterAction{Invisibility (Recharges after a Short or Long Rest)}
The spy magically turns invisible until it attacks, deals damage, casts a spell, or uses its Enlarge, or until its concentration is broken, up to 1 hour (as if concentration on a spell). Any equipment the spy wears or carries is invisible with it.

\end{DndMonster}



\begin{DndMonster}[float=t]{Efreeti}
\DndMonsterType{Large elemental, lawful evil}
\DndMonsterBasics[
armor-class = {17 (natural armor)},
hit-points = {\DndDice{16d10 + 112}},
speed = {40 ft., fly 60 ft.}
]
\DndMonsterAbilityScores[
 str = 22,
 dex = 12,
 con = 24,
 int = 16,
 wis = 15,
 cha = 16
]
\DndMonsterDetails[
saving-throws = {Int +7, Wis +6, Cha +7},
damage-immunities = {fire},
senses = {darkvision 120 ft., passive perception 12},
languages = {Ignan},
challenge = 11,
]
\DndMonsterAction{Elemental Demise}
If the efreeti dies, its body disintegrates in a flash of fire and puff of smoke, leaving behind only equipment the efreeti was wearing or carrying.

\DndMonsterAction{Innate Spellcasting}
The efreeti's innate spellcasting ability is Charisma (spell save DC 15, 7 to hit with spell attacks). It can innately cast the following spells, requiring no material components:
\begin{DndMonsterSpells}
  \DndInnateSpellLevel{\textit{detect magic}}
  \DndInnateSpellLevel[3e]{\textit{enlarge/reduce}, \textit{tongues}}
  \DndInnateSpellLevel[1e]{\textit{conjure elemental} (\textbf{fire elemental} only), \textit{gaseous form}, \textit{invisibility}, \textit{major image}, \textit{plane shift}, \textit{wall of fire}}
\end{DndMonsterSpells}
\DndMonsterSection{Actions}
\DndMonsterAction{Multiattack}
The efreeti makes two scimitar attacks or uses its Hurl Flame twice.

\DndMonsterAction{Scimitar}
\textsl{Melee Weapon Attack:} 10 to hit, reach 5 ft., one target. \textsl{Hit:} 13 (2d6 + 6) slashing damage plus 7 (2d6) fire damage.

\DndMonsterAction{Hurl Flame}
\textsl{Ranged Spell Attack:} 7 to hit, range 120 ft., one target. \textsl{Hit:} 17 (5d6) fire damage.

\end{DndMonster}



\begin{DndMonster}[float=t]{Enchanter}
\DndMonsterType{Medium humanoid (any race), any alignment}
\DndMonsterBasics[
armor-class = {12 (15 with \textit{mage armor})},
hit-points = {\DndDice{9d8}},
speed = {30 ft.}
]
\DndMonsterAbilityScores[
 str = 9,
 dex = 14,
 con = 11,
 int = 17,
 wis = 12,
 cha = 11
]
\DndMonsterDetails[
saving-throws = {Int +6, Wis +4},
skills = {Arcana +6, History +6},
languages = {any four languages},
challenge = 5,
]
\DndMonsterAction{Spellcasting}
The enchanter is a 9th-level spellcaster. Its spellcasting ability is Intelligence (spell save DC 14, 6 to hit with spell attacks). The enchanter has the following wizard spells prepared:
\begin{DndMonsterSpells}
  \DndMonsterSpellLevel{\textit{friends}, \textit{mage hand}, \textit{mending}, \textit{message}}
   \DndMonsterSpellLevel[1][4]{\textit{charm person}*, \textit{mage armor}, \textit{magic missile}}
   \DndMonsterSpellLevel[2][3]{\textit{hold person}*, \textit{invisibility}, \textit{suggestion}*}
   \DndMonsterSpellLevel[3][3]{\textit{fireball}, \textit{haste}, \textit{tongues}}
   \DndMonsterSpellLevel[4][3]{\textit{dominate beast}*, \textit{stoneskin}}
   \DndMonsterSpellLevel[5][2]{\textit{hold monster}*}
\end{DndMonsterSpells}
\DndMonsterSection{Actions}
\DndMonsterAction{Quarterstaff}
\textsl{Melee Weapon Attack:} 2 to hit, reach 5 ft., one target. \textsl{Hit:} 2 (1d6 - 1) bludgeoning damage, or 3 (1d8 - 1) bludgeoning damage if used with two hands.

\DndMonsterSection{Reactions}
\DndMonsterAction{Instinctive Charm (Recharges after the Enchanter Casts an Enchantment Spell of 1st level or Higher)}
The enchanter tries to magically divert an attack made against it, provided that the attacker is within 30 feet of it and visible to it. The enchanter must decide to do so before the attack hits or misses.

The attacker must make a DC 14 Wisdom saving throw. On a failed save, the attacker targets the creature closest to it, other than the enchanter or itself. If multiple creatures are closest, the attacker chooses which one to target.

\end{DndMonster}



\begin{DndMonster}[float=t]{Evoker}
\DndMonsterType{Medium humanoid (any race), any alignment}
\DndMonsterBasics[
armor-class = {12 (15 with \textit{mage armor})},
hit-points = {\DndDice{12d8 + 12}},
speed = {30 ft.}
]
\DndMonsterAbilityScores[
 str = 9,
 dex = 14,
 con = 12,
 int = 17,
 wis = 12,
 cha = 11
]
\DndMonsterDetails[
saving-throws = {Int +7, Wis +5},
skills = {Arcana +7, History +7},
languages = {any four languages},
challenge = 9,
]
\DndMonsterAction{Sculpt Spells}
When the evoker casts an evocation spell that forces other creatures it can see to make a saving throw, it can choose a number of them equal to 1 + the spell's level. These creatures automatically succeed on their saving throws against the spell. If a successful save means a chosen creature would take half damage from the spell, it instead takes no damage from it.

\DndMonsterAction{Spellcasting}
The evoker is a 12th-level spellcaster. Its spellcasting ability is Intelligence (spell save DC 15, 7 to hit with spell attacks). The evoker has the following wizard spells prepared:
\begin{DndMonsterSpells}
  \DndMonsterSpellLevel{\textit{fire bolt}*, \textit{light}*, \textit{prestidigitation}, \textit{ray of frost}*}
   \DndMonsterSpellLevel[1][4]{\textit{burning hands}*, \textit{mage armor}, \textit{magic missile}*}
   \DndMonsterSpellLevel[2][3]{\textit{mirror image}, \textit{misty step}, \textit{shatter}*}
   \DndMonsterSpellLevel[3][3]{\textit{counterspell}, \textit{fireball}*, \textit{lightning bolt}*}
   \DndMonsterSpellLevel[4][3]{\textit{ice storm}*, \textit{stoneskin}}
   \DndMonsterSpellLevel[5][2]{\textit{Bigby's hand}*, \textit{cone of cold}*}
   \DndMonsterSpellLevel[6][1]{\textit{chain lightning}*, \textit{wall of ice}*}
\end{DndMonsterSpells}
\DndMonsterSection{Actions}
\DndMonsterAction{Quarterstaff}
\textsl{Melee Weapon Attack:} 3 to hit, reach 5 ft., one target. \textsl{Hit:} 2 (1d6 - 1) bludgeoning damage, or 3 (1d8 - 1) bludgeoning damage if used with two hands.

\end{DndMonster}



\begin{DndMonster}[float=t]{Flesh Golem}
\DndMonsterType{Medium construct, neutral}
\DndMonsterBasics[
armor-class = {9},
hit-points = {\DndDice{11d8 + 44}},
speed = {30 ft.}
]
\DndMonsterAbilityScores[
 str = 19,
 dex = 9,
 con = 18,
 int = 6,
 wis = 10,
 cha = 5
]
\DndMonsterDetails[
damage-immunities = {lightning; poison; bludgeoning, piercing, slashing from nonmagical attacks that aren't adamantine},
condition-immunities = {charmed, exhaustion, frightened, paralyzed, petrified, poisoned},
senses = {darkvision 60 ft., passive perception 10},
languages = {understands the languages of its creator but can't speak},
challenge = 5,
]
\DndMonsterAction{Berserk}
Whenever the golem starts its turn with 40 hit points or fewer, roll a d6. On a 6, the golem goes berserk. On each of its turns while berserk, the golem attacks the nearest creature it can see. If no creature is near enough to move to and attack, the golem attacks an object, with preference for an object smaller than itself. Once the golem goes berserk, it continues to do so until it is destroyed or regains all its hit points.

The golem's creator, if within 60 feet of the berserk golem, can try to calm it by speaking firmly and persuasively. The golem must be able to hear its creator, who must take an action to make a DC 15 Charisma (Persuasion) check. If the check succeeds, the golem ceases being berserk. If it takes damage while still at 40 hit points or fewer, the golem might go berserk again.

\DndMonsterAction{Aversion of Fire}
If the golem takes fire damage, it has disadvantage on attack rolls and ability checks until the end of its next turn.

\DndMonsterAction{Immutable Form}
The golem is immune to any spell or effect that would alter its form.

\DndMonsterAction{Lightning Absorption}
Whenever the golem is subjected to lightning damage, it takes no damage and instead regains a number of hit points equal to the lightning damage dealt.

\DndMonsterAction{Magic Resistance}
The golem has advantage on saving throws against spells and other magical effects.

\DndMonsterAction{Magic Weapons}
The golem's weapon attacks are magical.

\DndMonsterSection{Actions}
\DndMonsterAction{Multiattack}
The golem makes two slam attacks.

\DndMonsterAction{Slam}
\textsl{Melee Weapon Attack:} 7 to hit, reach 5 ft., one target. \textsl{Hit:} 13 (2d8 + 4) bludgeoning damage.

\end{DndMonster}



\begin{DndMonster}[float=t]{Gargoyle}
\DndMonsterType{Medium elemental, chaotic evil}
\DndMonsterBasics[
armor-class = {15 (natural armor)},
hit-points = {\DndDice{7d8 + 21}},
speed = {30 ft., fly 60 ft.}
]
\DndMonsterAbilityScores[
 str = 15,
 dex = 11,
 con = 16,
 int = 6,
 wis = 11,
 cha = 7
]
\DndMonsterDetails[
damage-resistances = {bludgeoning, piercing, slashing from nonmagical attacks that aren't adamantine},
damage-immunities = {poison},
condition-immunities = {exhaustion, petrified, poisoned},
senses = {darkvision 60 ft., passive perception 10},
languages = {Terran},
challenge = 2,
]
\DndMonsterAction{False Appearance}
While the gargoyle remains motionless, it is indistinguishable from an inanimate statue.

\DndMonsterSection{Actions}
\DndMonsterAction{Multiattack}
The gargoyle makes two attacks: one with its bite and one with its claws.

\DndMonsterAction{Bite}
\textsl{Melee Weapon Attack:} 4 to hit, reach 5 ft., one target. \textsl{Hit:} 5 (1d6 + 2) piercing damage.

\DndMonsterAction{Claws}
\textsl{Melee Weapon Attack:} 4 to hit, reach 5 ft., one target. \textsl{Hit:} 5 (1d6 + 2) slashing damage.

\end{DndMonster}



\begin{DndMonster}[float=t]{Ghoul}
\DndMonsterType{Medium undead, chaotic evil}
\DndMonsterBasics[
armor-class = {12},
hit-points = {\DndDice{5d8}},
speed = {30 ft.}
]
\DndMonsterAbilityScores[
 str = 13,
 dex = 15,
 con = 10,
 int = 7,
 wis = 10,
 cha = 6
]
\DndMonsterDetails[
damage-immunities = {poison},
condition-immunities = {charmed, exhaustion, poisoned},
senses = {darkvision 60 ft., passive perception 10},
languages = {Common},
challenge = 1,
]
\DndMonsterSection{Actions}
\DndMonsterAction{Bite}
\textsl{Melee Weapon Attack:} 2 to hit, reach 5 ft., one creature. \textsl{Hit:} 9 (2d6 + 2) piercing damage.

\DndMonsterAction{Claws}
\textsl{Melee Weapon Attack:} 4 to hit, reach 5 ft., one target. \textsl{Hit:} 7 (2d4 + 2) slashing damage. If the target is a creature other than an elf or undead, it must succeed on a DC 10 Constitution saving throw or be paralyzed for 1 minute. The target can repeat the saving throw at the end of each of its turns, ending the effect on itself on a success.

\end{DndMonster}


\clearpage

\begin{DndMonster}[float=t]{Giant Crab}
\DndMonsterType{Medium beast, unaligned}
\DndMonsterBasics[
armor-class = {15 (natural armor)},
hit-points = {\DndDice{3d8}},
speed = {30 ft., swim 30 ft.}
]
\DndMonsterAbilityScores[
 str = 13,
 dex = 15,
 con = 11,
 int = 1,
 wis = 9,
 cha = 3
]
\DndMonsterDetails[
skills = {Stealth +4},
senses = {blindsight 30 ft., passive perception 9},
challenge = 1/8,
]
\DndMonsterAction{Amphibious}
The crab can breathe air and water.

\DndMonsterSection{Actions}
\DndMonsterAction{Claw}
\textsl{Melee Weapon Attack:} 3 to hit, reach 5 ft., one target. \textsl{Hit:} 4 (1d6 + 1) bludgeoning damage, and the target is grappled (escape DC 11). The crab has two claws, each of which can grapple only one target.

\end{DndMonster}



\begin{DndMonster}[float=t]{Giant Crayfish}
\DndMonsterType{Large beast, unaligned}
\DndMonsterBasics[
armor-class = {15 (natural armor)},
hit-points = {\DndDice{7d10 + 7}},
speed = {30 ft., swim 30 ft.}
]
\DndMonsterAbilityScores[
 str = 15,
 dex = 13,
 con = 13,
 int = 1,
 wis = 9,
 cha = 3
]
\DndMonsterDetails[
skills = {Stealth +3},
senses = {blindsight 30 ft., passive perception 9},
challenge = 2,
]
\DndMonsterAction{Amphibious}
The giant crayfish can breathe air and water.

\DndMonsterSection{Actions}
\DndMonsterAction{Multiattack}
The giant crayfish makes two claw attacks.

\DndMonsterAction{Claw}
\textsl{Melee Weapon Attack:} 4 to hit, reach 5 ft., one target. \textsl{Hit:} 7 (1d10 + 2) bludgeoning damage, and the target is grappled (escape DC 12). The crayfish has two claws, each of which can grapple only one target.

\end{DndMonster}



\begin{DndMonster}[float=t]{Giant Ice Toad}
\DndMonsterType{Large monstrosity, neutral}
\DndMonsterBasics[
armor-class = {14 (natural armor)},
hit-points = {\DndDice{7d10 + 14}},
speed = {30 ft.}
]
\DndMonsterAbilityScores[
 str = 16,
 dex = 13,
 con = 14,
 int = 8,
 wis = 10,
 cha = 6
]
\DndMonsterDetails[
damage-immunities = {cold},
senses = {darkvision 60 ft., passive perception 10},
languages = {Ice Toad},
challenge = 3,
]
\DndMonsterAction{Amphibious}
The toad can breathe air and water.

\DndMonsterAction{Cold Aura}
Any creature that starts its turn within 10 feet of the toad takes 5 (1d10) cold damage.

\DndMonsterAction{Standing Leap}
The toad's long jump is up to 20 feet and its high jump is up to 10 feet, with or without a running start.

\DndMonsterSection{Actions}
\DndMonsterAction{Bite}
\textsl{Melee Weapon Attack:} 5 to hit, reach 5 ft., one target. \textsl{Hit:} 10 (2d6 + 3) piercing damage, and the target is grappled (escape DC 13). Until this grapple ends, the target is restrained, and the toad can't bite another target.

\DndMonsterAction{Swallow}
\textsl{Melee Weapon Attack:} 5 to hit, reach 5 ft., one Medium or smaller creature the toad is grappling. \textsl{Hit:} 10 (2d6 + 3) piercing damage, the target is swallowed, and the grapple ends. The swallowed target is blinded and restrained, it has total cover against attacks and other effects outside the toad, and it takes 10 (3d6) acid damage and 11 (2d10) cold damage at the start of each of the toad's turns. The toad can have only one target swallowed at a time.

If the toad dies, a swallowed creature is no longer restrained by it and can escape from the corpse using 5 feet of movement, exiting prone.

\end{DndMonster}



\begin{DndMonster}[float=t]{Giant Lightning Eel}
\DndMonsterType{Large beast, unaligned}
\DndMonsterBasics[
armor-class = {13},
hit-points = {\DndDice{5d10 + 15}},
speed = {5 ft., swim 30 ft.}
]
\DndMonsterAbilityScores[
 str = 11,
 dex = 17,
 con = 16,
 int = 2,
 wis = 12,
 cha = 3
]
\DndMonsterDetails[
damage-resistances = {lightning},
senses = {blindsight 60 ft., passive perception 11},
challenge = 3,
]
\DndMonsterAction{Water Breathing}
The eel can breathe only underwater.

\DndMonsterSection{Actions}
\DndMonsterAction{Multiattack}
The eel makes two bite attacks.

\DndMonsterAction{Bite}
\textsl{Melee Weapon Attack:} 5 to hit, reach 5 ft., one target. \textsl{Hit:} 10 (2d6 + 3) piercing damage plus 4 (1d8) lightning damage.

\DndMonsterAction{Lightning Jolt (Recharge 5\textendash 6)}
One creature the eel touches within 5 feet of it outside water, or each creature within 15 feet of it in a body of water, must make a DC 12 Constitution saving throw. On failed save, a target takes 13 (3d8) lightning damage. If the target takes any of this damage, the target is stunned until the end of the eel's next turn. On a successful save, a target takes half as much damage and isn't stunned.

\end{DndMonster}



\begin{DndMonster}[float=t]{Giant Scorpion}
\DndMonsterType{Large beast, unaligned}
\DndMonsterBasics[
armor-class = {15 (natural armor)},
hit-points = {\DndDice{7d10 + 14}},
speed = {40 ft.}
]
\DndMonsterAbilityScores[
 str = 15,
 dex = 13,
 con = 15,
 int = 1,
 wis = 9,
 cha = 3
]
\DndMonsterDetails[
senses = {blindsight 60 ft., passive perception 9},
challenge = 3,
]
\DndMonsterSection{Actions}
\DndMonsterAction{Multiattack}
The scorpion makes three attacks: two with its claws and one with its sting.

\DndMonsterAction{Claw}
\textsl{Melee Weapon Attack:} 4 to hit, reach 5 ft., one target. \textsl{Hit:} 6 (1d8 + 2) bludgeoning damage, and the target is grappled (escape DC 12). The scorpion has two claws, each of which can grapple only one target.

\DndMonsterAction{Sting}
\textsl{Melee Weapon Attack:} 4 to hit, reach 5 ft., one creature. \textsl{Hit:} 7 (1d10 + 2) piercing damage, and the target must make a DC 12 Constitution saving throw, taking 22 (4d10) poison damage on a failed save, or half as much damage on a successful one.

\end{DndMonster}



\begin{DndMonster}[float=t]{Giant Skeleton}
\DndMonsterType{Huge undead, neutral evil}
\DndMonsterBasics[
armor-class = {17 (natural armor)},
hit-points = {\DndDice{10d12 + 50}},
speed = {30 ft.}
]
\DndMonsterAbilityScores[
 str = 21,
 dex = 10,
 con = 20,
 int = 4,
 wis = 6,
 cha = 6
]
\DndMonsterDetails[
damage-vulnerabilities = {bludgeoning},
damage-immunities = {poison},
condition-immunities = {exhaustion, poisoned},
senses = {darkvision 60 ft., passive perception 8},
languages = {understands Giant but can't speak},
challenge = 7,
]
\DndMonsterAction{Evasion}
If the skeleton is subjected to an effect that allows it to make a saving throw to take only half damage, it instead takes no damage if it succeeds on the saving throw, and only half damage if it fails.

\DndMonsterAction{Magic Resistance}
The skeleton has advantage on saving throws against spells and other magical effects.

\DndMonsterAction{Turn Immunity}
The skeleton is immune to effects that turn undead.

\DndMonsterSection{Actions}
\DndMonsterAction{Multiattack}
The skeleton makes three scimitar attacks.

\DndMonsterAction{Scimitar}
\textsl{Melee Weapon Attack:} 8 to hit, reach 10 ft., one target. \textsl{Hit:} 15 (3d6 + 5) slashing damage.

\end{DndMonster}



\begin{DndMonster}[float=t]{Giant Subterranean Lizard}
\DndMonsterType{Huge beast, unaligned}
\DndMonsterBasics[
armor-class = {14 (natural armor)},
hit-points = {\DndDice{7d12 + 21}},
speed = {30 ft., swim 50 ft.}
]
\DndMonsterAbilityScores[
 str = 21,
 dex = 9,
 con = 17,
 int = 2,
 wis = 10,
 cha = 7
]
\DndMonsterDetails[
skills = {Stealth +3},
challenge = 4,
]
\DndMonsterSection{Actions}
\DndMonsterAction{Multiattack}
The lizard makes two attacks: one with its bite and one with its tail. One attack can be replaced by Swallow.

\DndMonsterAction{Bite}
\textsl{Melee Weapon Attack:} 7 to hit, reach 5 ft., one target. \textsl{Hit:} 16 (2d10 + 5) piercing damage and the target is grappled (escape DC 15). Until this grapple ends, the target is restrained, and the lizard can't bite another target.

\DndMonsterAction{Tail}
\textsl{Melee Weapon Attack:} 7 to hit, reach 10 ft., one target. \textsl{Hit:} 12 (2d6 + 5) piercing damage. If the target is a creature, it must succeed on a DC 15 Strength saving throw or be knocked prone.

\DndMonsterAction{Swallow}
\textsl{Melee Weapon Attack:} 7 to hit, reach 5 ft., one Medium or smaller creature the lizard is grappling. \textsl{Hit:} 16 (2d10 + 5) piercing damage. The target is swallowed, and the grapple ends. The swallowed target is blinded and restrained, it has total cover against attacks and other effects outside the lizard, and it takes 10 (3d6) acid damage at the start of each of the lizard's turns. The lizard can have only one target swallowed at a time.

If the lizard dies, a swallowed creature is no longer restrained by it and can escape from the corpse using 10 feet of movement, exiting prone.

\end{DndMonster}



\begin{DndMonster}[float=t]{Gray Ooze}
\DndMonsterType{Medium ooze, unaligned}
\DndMonsterBasics[
armor-class = {8},
hit-points = {\DndDice{3d8 + 9}},
speed = {10 ft., climb 10 ft.}
]
\DndMonsterAbilityScores[
 str = 12,
 dex = 6,
 con = 16,
 int = 1,
 wis = 6,
 cha = 2
]
\DndMonsterDetails[
skills = {Stealth +2},
damage-resistances = {acid, cold, fire},
condition-immunities = {blinded, charmed, deafened, exhaustion, frightened, prone},
senses = {blindsight 60 ft. (blind beyond this radius), passive perception 8},
challenge = 1/2,
]
\DndMonsterAction{Amorphous}
The ooze can move through a space as narrow as 1 inch wide without squeezing.

\DndMonsterAction{Corrode Metal}
Any nonmagical weapon made of metal that hits the ooze corrodes. After dealing damage, the weapon takes a permanent and cumulative −1 penalty to damage rolls. If its penalty drops to −5, the weapon is destroyed. Nonmagical ammunition made of metal that hits the ooze is destroyed after dealing damage.

The ooze can eat through 2-inch-thick, nonmagical metal in 1 round.

\DndMonsterAction{False Appearance}
While the ooze remains motionless, it is indistinguishable from an oily pool or wet rock.

\DndMonsterSection{Actions}
\DndMonsterAction{Pseudopod}
\textsl{Melee Weapon Attack:} 3 to hit, reach 5 ft., one target. \textsl{Hit:} 4 (1d6 + 1) bludgeoning damage plus 7 (2d6) acid damage, and if the target is wearing nonmagical metal armor, its armor is partly corroded and takes a permanent and cumulative −1 penalty to the AC it offers. The armor is destroyed if the penalty reduces its AC to 10.

\end{DndMonster}



\begin{DndMonster}[float=t]{Greater Zombie}
\DndMonsterType{Medium undead, neutral evil}
\DndMonsterBasics[
armor-class = {15 (natural armor)},
hit-points = {\DndDice{13d8 + 39}},
speed = {30 ft.}
]
\DndMonsterAbilityScores[
 str = 18,
 dex = 10,
 con = 17,
 int = 4,
 wis = 6,
 cha = 6
]
\DndMonsterDetails[
saving-throws = {Wis +1},
damage-resistances = {cold, necrotic},
damage-immunities = {poison},
condition-immunities = {charmed, exhaustion, frightened, paralyzed, poisoned},
senses = {darkvision 60 ft., passive perception 8},
languages = {understands the languages it knew in life but can't speak},
challenge = 5,
]
\DndMonsterAction{Turn Resistance}
The zombie has advantage on saving throws against any effect that turns undead.

\DndMonsterAction{Undead Fortitude}
If damage reduces the zombie to 0 hit points, it must make a Constitution saving throw with a DC of 5 + the damage taken, unless the damage is radiant or from a critical hit. On a success, the zombie drops to 1 hit point instead.

\DndMonsterSection{Actions}
\DndMonsterAction{Multiattack}
The zombie makes two melee attacks.

\DndMonsterAction{Empowered Slam}
\textsl{Melee Weapon Attack:} 7 to hit, reach 5 ft., one target. \textsl{Hit:} 7 (1d6 + 4) bludgeoning damage and 7 (2d6) necrotic damage.

\end{DndMonster}



\begin{DndMonster}[float*=b,width=\textwidth + 8pt]{Gynosphinx}
\begin{multicols}{2}
\DndMonsterType{Large monstrosity, lawful neutral}
\DndMonsterBasics[
armor-class = {17 (natural armor)},
hit-points = {\DndDice{16d10 + 48}},
speed = {40 ft., fly 60 ft.}
]
\DndMonsterAbilityScores[
 str = 18,
 dex = 15,
 con = 16,
 int = 18,
 wis = 18,
 cha = 18
]
\DndMonsterDetails[
skills = {Arcana +12, History +12, Perception +8, Religion +8},
damage-resistances = {bludgeoning, piercing, slashing from nonmagical attacks},
damage-immunities = {psychic},
condition-immunities = {charmed, frightened},
senses = {truesight 120 ft., passive perception 18},
languages = {Common, Sphinx},
challenge = 11,
]
\DndMonsterAction{Inscrutable}
The sphinx is immune to any effect that would sense its emotions or read its thoughts, as well as any divination spell that it refuses. Wisdom (Insight) checks made to ascertain the sphinx's intentions or sincerity have disadvantage.

\DndMonsterAction{Magic Weapons}
The sphinx's weapon attacks are magical.

\DndMonsterAction{Spellcasting}
The sphinx is a 9th-level spellcaster. Its spellcasting ability is Intelligence (spell save DC 16, 8 to hit with spell attacks). It requires no material components to cast its spells. The sphinx has the following wizard spells prepared:
\begin{DndMonsterSpells}
  \DndMonsterSpellLevel{\textit{mage hand}, \textit{minor illusion}, \textit{prestidigitation}}
   \DndMonsterSpellLevel[1][4]{\textit{detect magic}, \textit{identify}, \textit{shield}}
   \DndMonsterSpellLevel[2][3]{\textit{darkness}, \textit{locate object}, \textit{suggestion}}
   \DndMonsterSpellLevel[3][3]{\textit{dispel magic}, \textit{remove curse}, \textit{tongues}}
   \DndMonsterSpellLevel[4][3]{\textit{banishment}, \textit{greater invisibility}}
   \DndMonsterSpellLevel[5][1]{\textit{legend lore}}
\end{DndMonsterSpells}
\DndMonsterSection{Actions}
\DndMonsterAction{Multiattack}
The sphinx makes two claw attacks.

\DndMonsterAction{Claw}
\textsl{Melee Weapon Attack:} 8 to hit, reach 5 ft., one target. \textsl{Hit:} 13 (2d8 + 4) slashing damage.

\DndMonsterSection{Legendary Actions}
Gynosphinx can take 3 legendary actions, choosing from the options below. Only one legendary action can be used at a time and only at the end of another creature's turn. Gynosphinx regains spent legendary actions at the start of its turn.

\begin{DndMonsterLegendaryActions}
\DndMonsterLegendaryAction{Claw Attack}{The sphinx makes one claw attack.
}
\DndMonsterLegendaryAction{Teleport (Costs 2 Actions)}{The sphinx magically teleports, along with any equipment it is wearing or carrying, up to 120 feet to an unoccupied space it can see.
}
\DndMonsterLegendaryAction{Cast a Spell (Costs 3 Actions)}{The sphinx casts a spell from its list of prepared spells, using a spell slot as normal.
}
\end{DndMonsterLegendaryActions}
\end{multicols}
\end{DndMonster}



\begin{DndMonster}[float=t]{Hell Hound}
\DndMonsterType{Medium fiend, lawful evil}
\DndMonsterBasics[
armor-class = {15 (natural armor)},
hit-points = {\DndDice{7d8 + 14}},
speed = {50 ft.}
]
\DndMonsterAbilityScores[
 str = 17,
 dex = 12,
 con = 14,
 int = 6,
 wis = 13,
 cha = 6
]
\DndMonsterDetails[
skills = {Perception +5},
damage-immunities = {fire},
senses = {darkvision 60 ft., passive perception 15},
languages = {understands Infernal but can't speak it},
challenge = 3,
]
\DndMonsterAction{Keen Hearing and Smell}
The hound has advantage on Wisdom (Perception) checks that rely on hearing or smell.

\DndMonsterAction{Pack Tactics}
The hound has advantage on an attack roll against a creature if at least one of the hound's allies is within 5 feet of the creature and the ally isn't incapacitated.

\DndMonsterSection{Actions}
\DndMonsterAction{Bite}
\textsl{Melee Weapon Attack:} 5 to hit, reach 5 ft., one target. \textsl{Hit:} 7 (1d8 + 3) piercing damage plus 7 (2d6) fire damage.

\DndMonsterAction{Fire Breath (Recharge 5\textendash 6)}
The hound exhales fire in a 15-foot cone. Each creature in that area must make a DC 12 Dexterity saving throw, taking 21 (6d6) fire damage on a failed save, or half as much damage on a successful one.

\end{DndMonster}



\begin{DndMonster}[float=t]{Huge Giant Crab}
\DndMonsterType{Huge beast, unaligned}
\DndMonsterBasics[
armor-class = {15 (natural armor)},
hit-points = {\DndDice{14d12 + 70}},
speed = {30 ft., swim 30 ft.}
]
\DndMonsterAbilityScores[
 str = 20,
 dex = 15,
 con = 20,
 int = 1,
 wis = 9,
 cha = 3
]
\DndMonsterDetails[
skills = {Stealth +5},
condition-immunities = {charmed, frightened, paralyzed},
senses = {blindsight 30 ft., passive perception 9},
challenge = 8,
]
\DndMonsterAction{Banded Claw}
On one of its claws, the crab wears a rune-covered copper band that makes it immune to being charmed, frightened, and paralyzed. The copper band is worthless as a treasure, as the magic is keyed to this crab.

\DndMonsterAction{Amphibious}
The crab can breathe air and water.

\DndMonsterSection{Actions}
\DndMonsterAction{Claw}
\textsl{Melee Weapon Attack:} 9 to hit, reach 10 ft., one target. \textsl{Hit:} 27 (4d10 + 5) bludgeoning damage, and the target is grappled, escape DC 14. The crab has two claws, each of which can grapple only one target.

\end{DndMonster}



\begin{DndMonster}[float=t]{Illusionist}
\DndMonsterType{Medium humanoid (any race), any alignment}
\DndMonsterBasics[
armor-class = {12 (15 with \textit{mage armor})},
hit-points = {\DndDice{7d8 + 7}},
speed = {30 ft.}
]
\DndMonsterAbilityScores[
 str = 9,
 dex = 14,
 con = 13,
 int = 16,
 wis = 11,
 cha = 12
]
\DndMonsterDetails[
saving-throws = {Int +5, Wis +2},
skills = {Arcana +5, History +5},
languages = {any four languages},
challenge = 3,
]
\DndMonsterAction{Displacement (Recharges after the Illusionist Casts an Illusion Spell of 1st Level or Higher)}
As a bonus action, the illusionist projects an illusion that makes the illusionist appear to be standing in a place a few inches from its actual location, causing any creature to have disadvantage on attack rolls against the illusionist. The effect ends if the illusionist takes damage, it is incapacitated, or its speed becomes 0.

\DndMonsterAction{Spellcasting}
The illusionist is a 7th-level spellcaster. its spellcasting ability is Intelligence (spell save DC 13, 5 to hit with spell attacks). The illusionist has the following wizard spells prepared:
\begin{DndMonsterSpells}
  \DndMonsterSpellLevel{\textit{dancing lights}, \textit{mage hand}, \textit{minor illusion}, \textit{poison spray}}
   \DndMonsterSpellLevel[1][4]{\textit{color spray}*, \textit{disguise self}*, \textit{mage armor}, \textit{magic missile}}
   \DndMonsterSpellLevel[2][3]{\textit{invisibility}*, \textit{mirror image}*, \textit{phantasmal force}*}
   \DndMonsterSpellLevel[3][3]{\textit{major image}*, \textit{phantom steed}*}
   \DndMonsterSpellLevel[4][1]{\textit{phantasmal killer}*}
\end{DndMonsterSpells}
\DndMonsterSection{Actions}
\DndMonsterAction{Quarterstaff}
\textsl{Melee Weapon Attack:} 1 to hit, reach 5 ft., one target. \textsl{Hit:} 2 (1d6 - 1) bludgeoning damage, or 3 (1d8 - 1) bludgeoning damage if used with two hands.

\end{DndMonster}



\begin{DndMonster}[float=t]{Invisible Stalker}
\DndMonsterType{Medium elemental, neutral}
\DndMonsterBasics[
armor-class = {14},
hit-points = {\DndDice{16d8 + 32}},
speed = {50 ft., fly 50 ft. \textit{(hover)}}
]
\DndMonsterAbilityScores[
 str = 16,
 dex = 19,
 con = 14,
 int = 10,
 wis = 15,
 cha = 11
]
\DndMonsterDetails[
skills = {Perception +8, Stealth +10},
damage-resistances = {bludgeoning, piercing, slashing from nonmagical attacks},
damage-immunities = {poison},
condition-immunities = {exhaustion, grappled, paralyzed, petrified, poisoned, prone, restrained, unconscious},
senses = {darkvision 60 ft., passive perception 18},
languages = {Auran, understands Common but doesn't speak it},
challenge = 6,
]
\DndMonsterAction{Invisibility}
The stalker is invisible.

\DndMonsterAction{Faultless Tracker}
The stalker is given a quarry by its summoner. The stalker knows the direction and distance to its quarry as long as the two of them are on the same plane of existence. The stalker also knows the location of its summoner.

\DndMonsterSection{Actions}
\DndMonsterAction{Multiattack}
The stalker makes two slam attacks.

\DndMonsterAction{Slam}
\textsl{Melee Weapon Attack:} 6 to hit, reach 5 ft., one target. \textsl{Hit:} 10 (2d6 + 3) bludgeoning damage.

\end{DndMonster}



\begin{DndMonster}[float=t]{Kalka-Kylla}
\DndMonsterType{Large monstrosity, neutral}
\DndMonsterBasics[
armor-class = {15 (natural armor)},
hit-points = {\DndDice{10d10 + 30}},
speed = {30 ft., swim 30 ft.}
]
\DndMonsterAbilityScores[
 str = 17,
 dex = 12,
 con = 16,
 int = 15,
 wis = 16,
 cha = 12
]
\DndMonsterDetails[
skills = {Deception +3, Insight +5, Stealth +3},
senses = {blindsight 30 ft., passive perception 13},
languages = {Olman},
challenge = 3,
]
\DndMonsterAction{Amphibious}
Kalka-Kylla can breathe air and water.

\DndMonsterAction{False Appearance}
While Kalka-Kylla remains motionless and hidden in its shell, it is indistinguishable from a polished boulder.

\DndMonsterAction{Shell}
Kalka-Kylla can use a bonus action to retract into or emerge from its shell. While retracted, Kalka-Kylla gains a +4 bonus to AC, and it has a speed of 0 and can't benefit from bonuses to speed.

\DndMonsterSection{Actions}
\DndMonsterAction{Multiattack}
Kalka-Kylla makes two claw attacks.

\DndMonsterAction{Claw}
\textsl{Melee Weapon Attack:} 5 to hit, reach 5 ft., one target. \textsl{Hit:} 10 (2d6 + 3) bludgeoning damage, and if the target is a Medium or smaller creature, it is grappled (escape DC 13). Until this grapple ends, the target is restrained. Kalka-Kylla has two claws, each of which can grapple only one target.

\end{DndMonster}



\begin{DndMonster}[float=t]{Kelpie}
\DndMonsterType{Medium plant, neutral evil}
\DndMonsterBasics[
armor-class = {14 (natural armor)},
hit-points = {\DndDice{9d8 + 27}},
speed = {10 ft., swim 30 ft.}
]
\DndMonsterAbilityScores[
 str = 14,
 dex = 14,
 con = 16,
 int = 7,
 wis = 12,
 cha = 10
]
\DndMonsterDetails[
skills = {Perception +3, Stealth +4},
damage-resistances = {fire, bludgeoning, piercing},
condition-immunities = {blinded, deafened, exhaustion},
senses = {blindsight 60 ft., passive perception 13},
languages = {Common, Sylvan},
challenge = 4,
]
\DndMonsterAction{Amphibious}
The kelpie can breathe air and water.

\DndMonsterAction{Seaweed Shape}
The kelpie can use its action to reshape its body into the form of a humanoid or beast that is Small, Medium, or Large. Its statistics are otherwise unchanged. The disguise is convincing, unless the kelpie is in bright light or the viewer is within 30 feet of it, in which case the seams between the seaweed strands are visible. The kelpie returns to its true form if takes a bonus action to do so or if it dies.

\DndMonsterAction{False Appearance}
While the kelpie remains motionless in its true form, it is indistinguishable from normal seaweed.

\DndMonsterSection{Actions}
\DndMonsterAction{Multiattack}
The kelpie makes two slam attacks.

\DndMonsterAction{Slam}
\textsl{Melee Weapon Attack:} 4 to hit, reach 10 ft., one target. \textsl{Hit:} 11 (2d8 + 2) piercing damage. If the target is a Medium or smaller creature, it is grappled (escape DC 12).

\DndMonsterAction{Drowning Hypnosis}
The kelpie chooses one humanoid it can see within 150 feet of it. If the target can see the kelpie, the target must succeed on a DC 11 Wisdom saving throw or be magically charmed while the kelpie maintains concentration, up to 10 minutes (as if concentration on a spell). The charmed target is incapacitated, and instead of holding its breath underwater, it tries to breathe normally and immediately runs out of breath, unless it can breathe water. If the charmed target is more than 5 feet away from the kelpie, the target must move on its turn toward the kelpie by the most direct route, trying to get within 5 feet. It doesn't avoid opportunity attacks.

Before moving into damaging terrain, such as lava or a pit, and whenever it takes damage from a source other than the kelpie or drowning, the target can repeat the saving throw. A charmed target can also repeat the saving throw at the end of each of its turns. If the saving throw is successful, the effect ends on it.

A target that successfully saves is immune to this kelpie's hypnosis for the next 24 hours.

\end{DndMonster}



\begin{DndMonster}[float=t]{Knight}
\DndMonsterType{Medium humanoid (any race), any alignment}
\DndMonsterBasics[
armor-class = {18 (\textit{plate armor})},
hit-points = {\DndDice{8d8 + 16}},
speed = {30 ft.}
]
\DndMonsterAbilityScores[
 str = 16,
 dex = 11,
 con = 14,
 int = 11,
 wis = 11,
 cha = 15
]
\DndMonsterDetails[
saving-throws = {Con +4, Wis +2},
languages = {any one language (usually Common)},
challenge = 3,
]
\DndMonsterAction{Brave}
The knight has advantage on saving throws against being frightened.

\DndMonsterSection{Actions}
\DndMonsterAction{Multiattack}
The knight makes two melee attacks.

\DndMonsterAction{Greatsword}
\textsl{Melee Weapon Attack:} 5 to hit, reach 5 ft., one target. \textsl{Hit:} 10 (2d6 + 3) slashing damage.

\DndMonsterAction{Heavy Crossbow}
\textsl{Ranged Weapon Attack:} 2 to hit, range 100/400 ft., one target. \textsl{Hit:} 5 (1d10) piercing damage.

\DndMonsterAction{Leadership (Recharges after a Short or Long Rest)}
For 1 minute, the knight can utter a special command or warning whenever a nonhostile creature that it can see within 30 feet of it makes an attack roll or a saving throw. The creature can add a d4 to its roll provided it can hear and understand the knight. A creature can benefit from only one Leadership die at a time. This effect ends if the knight is incapacitated.

\DndMonsterSection{Reactions}
\DndMonsterAction{Parry}
The knight adds 2 to its AC against one melee attack that would hit it. To do so, the knight must see the attacker and be wielding a melee weapon.

\end{DndMonster}



\begin{DndMonster}[float=t]{Leucrotta}
\DndMonsterType{Large monstrosity, chaotic evil}
\DndMonsterBasics[
armor-class = {14 (natural armor)},
hit-points = {\DndDice{9d10 + 18}},
speed = {50 ft.}
]
\DndMonsterAbilityScores[
 str = 18,
 dex = 14,
 con = 15,
 int = 9,
 wis = 12,
 cha = 6
]
\DndMonsterDetails[
skills = {Deception +2, Perception +3},
senses = {darkvision 60 ft., passive perception 13},
languages = {Abyssal, Gnoll},
challenge = 3,
]
\DndMonsterAction{Keen Smell}
The leucrotta has advantage on Wisdom (Perception) checks that rely on smell.

\DndMonsterAction{Kicking Retreat}
If the leucrotta attacks with its hooves, it can take the Disengage action as a bonus action.

\DndMonsterAction{Mimicry}
The leucrotta can mimic animal sounds and humanoid voices. A creature that hears the sounds can tell they are imitations with a successful DC 14 Wisdom (Insight) check.

\DndMonsterAction{Rampage}
When the leucrotta reduces a creature to 0 hit points with a melee attack on its turn, it can take a bonus action to move up to half its speed and make an attack with its hooves.

\DndMonsterSection{Actions}
\DndMonsterAction{Multiattack}
The leucrotta makes two attacks: one with its bite and one with its hooves.

\DndMonsterAction{Bite}
\textsl{Melee Weapon Attack:} 6 to hit, reach 5 ft., one target. \textsl{Hit:} 8 (1d8 + 4) piercing damage. If the leucrotta scores a critical hit, it rolls the damage dice three times, instead of twice.

\DndMonsterAction{Hooves}
\textsl{Melee Weapon Attack:} 6 to hit, reach 5 ft., one target. \textsl{Hit:} 11 (2d6 + 4) bludgeoning damage.

\end{DndMonster}



\begin{DndMonster}[float=t]{Malformed Kraken}
\DndMonsterType{Huge monstrosity, chaotic evil}
\DndMonsterBasics[
armor-class = {17 (natural armor)},
hit-points = {\DndDice{15d12 + 75}},
speed = {20 ft., swim 40 ft.}
]
\DndMonsterAbilityScores[
 str = 25,
 dex = 11,
 con = 20,
 int = 11,
 wis = 15,
 cha = 15
]
\DndMonsterDetails[
saving-throws = {Str +11, Con +9, Int +4, Wis +6, Cha +6},
damage-resistances = {bludgeoning, piercing, slashing from nonmagical attacks},
damage-immunities = {lightning},
condition-immunities = {frightened, paralyzed},
senses = {truesight 60 ft., passive perception 12},
languages = {understands Common but can't speak, telepathy 60 ft.},
challenge = 10,
]
\DndMonsterAction{Amphibious}
The kraken can breathe air and water.

\DndMonsterAction{Siege Monster}
The kraken deals double damage to objects and structures.

\DndMonsterSection{Actions}
\DndMonsterAction{Multiattack}
The kraken makes three tentacle attacks. One of them can be replaced with a bite attack, and any of them can be replaced with Fling.

\DndMonsterAction{Bite}
\textsl{Melee Weapon Attack:} 11 to hit, reach 5 ft., one target. \textsl{Hit:} 16 (2d8 + 7) piercing damage.

\DndMonsterAction{Tentacle}
\textsl{Melee Weapon Attack:} 11 to hit, reach 20 ft., one target. \textsl{Hit:} 14 (2d6 + 7) bludgeoning damage, and the target is grappled (escape DC 16). Until this grapple ends, the target is restrained. The kraken has ten tentacles, each of which can grapple one target.

\DndMonsterAction{Fling}
One Medium or smaller object held or creature grappled by the kraken's tentacles is thrown up to 60 feet in a random direction and knocked prone. If a thrown target strikes a solid surface, the target takes 3 (1d6) bludgeoning damage for every 10 feet it was thrown. If the target is thrown at another creature, that creature must succeed on a DC 16 Dexterity saving throw or take the same damage and be knocked prone.

\DndMonsterAction{Lightning Storm}
The kraken creates three bolts of lightning, each of which can strike a target the kraken can see within 150 feet of it. A target must make a DC 16 Dexterity saving throw, taking 16 (3d10) lightning damage on a failed save, or half as much damage on a successful one.

\end{DndMonster}



\begin{DndMonster}[float=t]{Manticore}
\DndMonsterType{Large monstrosity, lawful evil}
\DndMonsterBasics[
armor-class = {14 (natural armor)},
hit-points = {\DndDice{8d10 + 24}},
speed = {30 ft., fly 50 ft.}
]
\DndMonsterAbilityScores[
 str = 17,
 dex = 16,
 con = 17,
 int = 7,
 wis = 12,
 cha = 8
]
\DndMonsterDetails[
senses = {darkvision 60 ft., passive perception 11},
languages = {Common},
challenge = 3,
]
\DndMonsterAction{Tail Spike Regrowth}
The manticore has twenty-four tail spikes. Used spikes regrow when the manticore finishes a long rest.

\DndMonsterSection{Actions}
\DndMonsterAction{Multiattack}
The manticore makes three attacks: one with its bite and two with its claws or three with its tail spikes.

\DndMonsterAction{Bite}
\textsl{Melee Weapon Attack:} 5 to hit, reach 5 ft., one target. \textsl{Hit:} 7 (1d8 + 3) piercing damage.

\DndMonsterAction{Claw}
\textsl{Melee Weapon Attack:} 5 to hit, reach 5 ft., one target. \textsl{Hit:} 6 (1d6 + 3) slashing damage.

\DndMonsterAction{Tail Spike}
\textsl{Ranged Weapon Attack:} 5 to hit, range 100/200 ft., one target. \textsl{Hit:} 7 (1d8 + 3) piercing damage.

\end{DndMonster}


\clearpage

\begin{DndMonster}[float=t]{Martial Arts Adept}
\DndMonsterType{Medium humanoid (any race), any alignment}
\DndMonsterBasics[
armor-class = {16},
hit-points = {\DndDice{11d8 + 11}},
speed = {40 ft.}
]
\DndMonsterAbilityScores[
 str = 11,
 dex = 17,
 con = 13,
 int = 11,
 wis = 16,
 cha = 10
]
\DndMonsterDetails[
skills = {Acrobatics +5, Insight +5, Stealth +5},
languages = {any one language (usually Common)},
challenge = 3,
]
\DndMonsterAction{Unarmored Defense}
While the adept is wearing no armor and wielding no shield, its AC includes its Wisdom modifier.

\DndMonsterSection{Actions}
\DndMonsterAction{Multiattack}
The adept makes three unarmed strikes or three dart attacks.

\DndMonsterAction{Unarmed Strike}
\textsl{Melee Weapon Attack:} 5 to hit, reach 5 ft., one target. \textsl{Hit:} 7 (1d8 + 3) bludgeoning damage. If the target is a creature, the adept can choose one of the following additional effects:


\begin{itemize}
\item The target must succeed on a DC 13 Strength saving throw or drop one item it is holding (adept's choice).
\item The target must succeed on a DC 13 Dexterity saving throw or be knocked prone.
\item The target must succeed on a DC 13 Constitution saving throw or be stunned until the end of the adept's next turn.
\end{itemize}

\DndMonsterAction{Dart}
\textsl{Ranged Weapon Attack:} 5 to hit, range 20/60 ft., one target. \textsl{Hit:} 5 (1d4 + 3) piercing damage.

\DndMonsterSection{Reactions}
\DndMonsterAction{Deflect Missile}
In response to being hit by a ranged weapon attack, the adept deflects the missile. The damage it takes from the attack is reduced by 1d10 + 3. If the damage is reduced to 0, the adept catches the missile if it's small enough to hold in one hand and the adept has a hand free.

\end{DndMonster}



\begin{DndMonster}[float=t]{Necromancer}
\DndMonsterType{Medium humanoid (any race), any alignment}
\DndMonsterBasics[
armor-class = {12 (15 with \textit{mage armor})},
hit-points = {\DndDice{12d8 + 12}},
speed = {30 ft.}
]
\DndMonsterAbilityScores[
 str = 9,
 dex = 14,
 con = 12,
 int = 17,
 wis = 12,
 cha = 11
]
\DndMonsterDetails[
saving-throws = {Int +7, Wis +5},
skills = {Arcana +7, History +7},
damage-resistances = {necrotic},
languages = {any four languages},
challenge = 9,
]
\DndMonsterAction{Grim Harvest (1/Turn)}
When necromancer kills a creature that is neither a construct nor undead with a spell of 1st level or higher, the necromancer regains hit points equal to twice the spell's level, or three times if it is a necromancy spell.

\DndMonsterAction{Spellcasting}
The necromancer is a 12th-level spellcaster. Its spellcasting ability is Intelligence (spell save DC 15, 7 to hit with spell attacks). The necromancer has the following wizard spells prepared:
\begin{DndMonsterSpells}
  \DndMonsterSpellLevel{\textit{chill touch}, \textit{dancing lights}, \textit{mage hand}, \textit{mending}}
   \DndMonsterSpellLevel[1][4]{\textit{false life}*, \textit{mage armor}, \textit{ray of sickness}*}
   \DndMonsterSpellLevel[2][3]{\textit{blindness/deafness}*, \textit{ray of enfeeblement}*, \textit{web}}
   \DndMonsterSpellLevel[3][3]{\textit{animate dead}*, \textit{bestow curse}*, \textit{vampiric touch}*}
   \DndMonsterSpellLevel[4][3]{\textit{blight}*, \textit{dimension door}, \textit{stoneskin}}
   \DndMonsterSpellLevel[5][2]{\textit{Bigby's hand}, \textit{cloudkill}}
   \DndMonsterSpellLevel[6][1]{\textit{circle of death}*}
\end{DndMonsterSpells}
\DndMonsterSection{Actions}
\DndMonsterAction{Withering Touch}
\textsl{Melee Spell Attack:} 7 to hit, reach 5 ft., one creature. \textsl{Hit:} 5 (2d4) necrotic damage.

\end{DndMonster}



\begin{DndMonster}[float=t]{Nereid}
\DndMonsterType{Medium fey, any chaotic alignment}
\DndMonsterBasics[
armor-class = {13},
hit-points = {\DndDice{8d8 + 8}},
speed = {30 ft., swim 60 ft.}
]
\DndMonsterAbilityScores[
 str = 10,
 dex = 17,
 con = 12,
 int = 13,
 wis = 14,
 cha = 16
]
\DndMonsterDetails[
skills = {Acrobatics +5, Nature +3, Stealth +5, Survival +4},
senses = {darkvision 60 ft., passive perception 12},
languages = {Aquan, Common, Elvish, Sylvan},
challenge = 2,
]
\DndMonsterAction{Amphibious}
The nereid can breathe air and water.

\DndMonsterAction{Aquatic Invisibility}
If immersed in water, the nereid can make itself invisible as a bonus action. It remains invisible until it leaves the water, ends the invisibility as a bonus action, or dies.

\DndMonsterAction{Mantle Dependent}
The nereid wears a mantle of silky cloth the color of sea foam, which holds the creature's spirit. The mantle has an AC and hit points equal to that of the nereid, but the garment can't be directly harmed while the nereid wears it. If the mantle is destroyed, the nereid becomes poisoned and dies within 1 hour. A nereid is willing to do anything in its power to recover the mantle if it is stolen, including serving the thief.

\DndMonsterAction{Shape Water}
The nereid can cast \textit{control water} at will, requiring no components. Its spellcasting ability for it is Charisma. This use of the spell has a range of 30 feet and can affect a cube of water no larger than 30 feet on a side.

\DndMonsterAction{Speak With Animals}
The nereid can comprehend and verbally communicate with beasts.

\DndMonsterSection{Actions}
\DndMonsterAction{Blinding Acid}
\textsl{Melee or Ranged Spell Attack:} 5 to hit, reach 5 ft. or range 30 ft., one target. \textsl{Hit:} 16 (2d12 + 3) acid damage, and the target is blinded until the start of the nereid's next turn.

\DndMonsterAction{Drowning Kiss (Recharge 5\textendash 6)}
The nereid touches one creature it can see within 5 feet of it. The target must succeed on a DC 13 Constitution saving throw or take 22 (3d12 + 3) acid damage. On a failure, it also runs out of breath and can't speak for 1 minute. At the end of each of its turns, it can repeat the save, ending the effect on itself on a success.

\DndMonsterAction{Water Lash}
The nereid causes a 5-foot cube of water within 60 feet of it to take a shape of its choice and strike one target it can see within 5 feet of that water. The target must make a DC 13 Strength saving throw. On a failed save, it takes 17 (4d6 + 3) bludgeoning damage, and if it is a Large or smaller creature, it is pushed up to 15 feet in a straight line or is knocked prone (nereid's choice). On a successful save, the target takes half as much damage and isn't pushed or knocked prone.

\end{DndMonster}



\begin{DndMonster}[float=t]{Ogre}
\DndMonsterType{Large giant, chaotic evil}
\DndMonsterBasics[
armor-class = {11 (\textit{hide armor})},
hit-points = {\DndDice{7d10 + 21}},
speed = {40 ft.}
]
\DndMonsterAbilityScores[
 str = 19,
 dex = 8,
 con = 16,
 int = 5,
 wis = 7,
 cha = 7
]
\DndMonsterDetails[
senses = {darkvision 60 ft., passive perception 8},
languages = {Common, Giant},
challenge = 2,
]
\DndMonsterSection{Actions}
\DndMonsterAction{Greatclub}
\textsl{Melee Weapon Attack:} 6 to hit, reach 5 ft., one target. \textsl{Hit:} 13 (2d8 + 4) bludgeoning damage.

\DndMonsterAction{Javelin}
\textsl{Melee or Ranged Weapon Attack:} 6 to hit, reach 5 ft. or range 30/120 ft., one target. \textsl{Hit:} 11 (2d6 + 4) piercing damage.

\end{DndMonster}



\begin{DndMonster}[float=t]{Oni}
\DndMonsterType{Large giant, lawful evil}
\DndMonsterBasics[
armor-class = {16 (\textit{chain mail})},
hit-points = {\DndDice{13d10 + 39}},
speed = {30 ft., fly 30 ft.}
]
\DndMonsterAbilityScores[
 str = 19,
 dex = 11,
 con = 16,
 int = 14,
 wis = 12,
 cha = 15
]
\DndMonsterDetails[
saving-throws = {Dex +3, Con +6, Wis +4, Cha +5},
skills = {Arcana +5, Deception +8, Perception +4},
senses = {darkvision 60 ft., passive perception 14},
languages = {Common, Giant},
challenge = 7,
]
\DndMonsterAction{Magic Weapons}
The oni's weapon attacks are magical.

\DndMonsterAction{Regeneration}
The oni regains 10 hit points at the start of its turn if it has at least 1 hit point.

\DndMonsterAction{Innate Spellcasting}
The oni's innate spellcasting ability is Charisma (spell save DC 13). The oni can innately cast the following spells, requiring no material components:
\begin{DndMonsterSpells}
  \DndInnateSpellLevel{\textit{darkness}, \textit{invisibility}}
  \DndInnateSpellLevel[1e]{\textit{charm person}, \textit{cone of cold}, \textit{gaseous form}, \textit{sleep}}
\end{DndMonsterSpells}
\DndMonsterSection{Actions}
\DndMonsterAction{Multiattack}
The oni makes two attacks, either with its claws or its glaive.

\DndMonsterAction{Claw (Oni Form Only)}
\textsl{Melee Weapon Attack:} 7 to hit, reach 5 ft., one target. \textsl{Hit:} 8 (1d8 + 4) slashing damage.

\DndMonsterAction{Glaive}
\textsl{Melee Weapon Attack:} 7 to hit, reach 10 ft., one target. \textsl{Hit:} 15 (2d10 + 4) slashing damage, or 9 (1d10 + 4) slashing damage in Small or Medium form.

\DndMonsterAction{Change Shape}
The oni magically polymorphs into a Small or Medium humanoid, into a Large giant, or back into its true form. Other than its size, its statistics are the same in each form. The only equipment that is transformed is its glaive, which shrinks so that it can be wielded in humanoid form. If the oni dies, it reverts to its true form, and its glaive reverts to its normal size.

\end{DndMonster}



\begin{DndMonster}[float=t]{Ooze Master}
\DndMonsterType{Huge undead, lawful evil}
\DndMonsterBasics[
armor-class = {9 (natural armor)},
hit-points = {\DndDice{12d12 + 60}},
speed = {30 ft., climb 30 ft.}
]
\DndMonsterAbilityScores[
 str = 16,
 dex = 1,
 con = 20,
 int = 17,
 wis = 10,
 cha = 16
]
\DndMonsterDetails[
saving-throws = {Int +7, Wis +4},
skills = {Arcana +7, Insight +4},
damage-resistances = {lightning; necrotic; bludgeoning, piercing, slashing from nonmagical attacks},
damage-immunities = {acid, cold, poison},
condition-immunities = {blinded, charmed, deafened, exhaustion, frightened, paralyzed, poisoned, prone},
senses = {blindsight 120 ft., passive perception 10},
languages = {Common, Primordial, Thayan},
challenge = 10,
]
\DndMonsterAction{Corrosive Form}
A creature that touches the Ooze Master or hits it with a melee attack while within 5 feet of it takes 9 (2d8) acid damage. Any nonmagical weapon that hits the Ooze Master corrodes. After dealing damage, the weapon takes a permanent and cumulative −1 penalty to damage rolls. If its penalty drops to −5, the weapon is destroyed. Nonmagical ammunition that hits the Ooze Master is destroyed after dealing damage.

The Ooze Master can eat through 2-inch-thick, nonmagical wood or metal in 1 round.

\DndMonsterAction{Instinctive Attack}
When the Ooze Master casts a spell with a casting time of 1 action, it can make one pseudopod attack as a bonus action.

\DndMonsterAction{Spider Climb}
The Ooze Master can climb difficult surfaces, including upside down on ceilings, without needing to make an ability check.

\DndMonsterAction{Spellcasting}
The Ooze Master is a 9th-level spellcaster. Its spellcasting ability is Intelligence (spell save DC 15, 7 to hit with spell attacks). It has the following wizard spells prepared:
\begin{DndMonsterSpells}
  \DndMonsterSpellLevel{\textit{acid splash}, \textit{friends}, \textit{mage hand}, \textit{poison spray}}
   \DndMonsterSpellLevel[1][4]{\textit{charm person}, \textit{detect magic}, \textit{magic missile}, \textit{ray of sickness}}
   \DndMonsterSpellLevel[2][3]{\textit{detect thoughts}, \textit{Melf's acid arrow}, \textit{suggestion}}
   \DndMonsterSpellLevel[3][3]{\textit{fear}, \textit{slow}, \textit{stinking cloud}}
   \DndMonsterSpellLevel[4][3]{\textit{confusion}, \textit{Evard's black tentacles}}
   \DndMonsterSpellLevel[5][1]{\textit{cloudkill}}
\end{DndMonsterSpells}
\DndMonsterSection{Actions}
\DndMonsterAction{Pseudopod}
\textsl{Melee Weapon Attack:} 7 to hit, reach 10 ft., one target. \textsl{Hit:} 13 (3d6 + 3) bludgeoning damage plus 10 (3d6) acid damage.

\DndMonsterSection{Reactions}
\DndMonsterAction{Instinctive Charm}
If a creature the Ooze Master can see makes an attack roll against it while within 30 feet of it, the Ooze Master can use a reaction to divert the attack if another creature is within the attack's range. The attacker must make a DC 15 Wisdom saving throw. On a failed save, the attacker targets the creature that is closest to it, not including itself or the Ooze Master. If multiple creatures are closest, the attacker chooses which one to target. On a successful save, the attacker is immune to this Instinctive Charm for 24 hours. Creatures that can't be charmed are immune to this effect.

\end{DndMonster}



\begin{DndMonster}[float=t]{Sea Lion}
\DndMonsterType{Large monstrosity, unaligned}
\DndMonsterBasics[
armor-class = {15 (natural armor)},
hit-points = {\DndDice{12d10 + 24}},
speed = {10 ft., swim 40 ft.}
]
\DndMonsterAbilityScores[
 str = 17,
 dex = 15,
 con = 15,
 int = 3,
 wis = 12,
 cha = 8
]
\DndMonsterDetails[
skills = {Perception +4, Stealth +5},
challenge = 5,
]
\DndMonsterAction{Amphibious}
The sea lion can breathe air and water.

\DndMonsterAction{Keen Smell}
The sea lion has advantage on Wisdom (Perception) checks that rely on smell.

\DndMonsterAction{Pack Tactics}
The sea lion has advantage on an attack roll against a creature if at least one of the sea lion's allies is within 5 feet of the creature and the ally isn't incapacitated.

\DndMonsterAction{Swimming Leap}
With a 10-foot swimming start, the sea lion can long jump out of or across the water up to 25 feet.

\DndMonsterSection{Actions}
\DndMonsterAction{Multiattack}
The sea lion makes three attacks: one bite attack and two claw attacks.

\DndMonsterAction{Bite}
\textsl{Melee Weapon Attack:} 6 to hit, reach 5 ft., one target. \textsl{Hit:} 12 (2d8 + 3) piercing damage.

\DndMonsterAction{Claw}
\textsl{Melee Weapon Attack:} 6 to hit, reach 5 ft., one target. \textsl{Hit:} 12 (2d8 + 3) piercing damage.

\end{DndMonster}



\begin{DndMonster}[float=t]{Shadow}
\DndMonsterType{Medium undead, chaotic evil}
\DndMonsterBasics[
armor-class = {12},
hit-points = {\DndDice{3d8 + 3}},
speed = {40 ft.}
]
\DndMonsterAbilityScores[
 str = 6,
 dex = 14,
 con = 13,
 int = 6,
 wis = 10,
 cha = 8
]
\DndMonsterDetails[
skills = {Stealth +4},
damage-vulnerabilities = {radiant},
damage-resistances = {acid; cold; fire; lightning; thunder; bludgeoning, piercing, slashing from nonmagical attacks},
damage-immunities = {necrotic, poison},
condition-immunities = {exhaustion, frightened, grappled, paralyzed, petrified, poisoned, prone, restrained},
senses = {darkvision 60 ft., passive perception 10},
challenge = 1/2,
]
\DndMonsterAction{Amorphous}
The shadow can move through a space as narrow as 1 inch wide without squeezing.

\DndMonsterAction{Shadow Stealth}
While in dim light or darkness, the shadow can take the Hide action as a bonus action. Its stealth bonus is also improved to +6.

\DndMonsterAction{Sunlight Weakness}
While in sunlight, the shadow has disadvantage on attack rolls, ability checks, and saving throws.

\DndMonsterSection{Actions}
\DndMonsterAction{Strength Drain}
\textsl{Melee Weapon Attack:} 4 to hit, reach 5 ft., one creature. \textsl{Hit:} 9 (2d6 + 2) necrotic damage, and the target's Strength score is reduced by 1d4. The target dies if this reduces its Strength to 0. Otherwise, the reduction lasts until the target finishes a short or long rest.

If a non-evil humanoid dies from this attack, a new shadow rises from the corpse 1d4 hours later.

\end{DndMonster}



\begin{DndMonster}[float=t]{Sharwyn Hucrele}
\DndMonsterType{Medium humanoid (human), neutral evil}
\DndMonsterBasics[
armor-class = {16 (Barkskin trait)},
hit-points = {\DndDice{2d8 + 4}},
speed = {30 ft.}
]
\DndMonsterAbilityScores[
 str = 11,
 dex = 13,
 con = 14,
 int = 16,
 wis = 14,
 cha = 9
]
\DndMonsterDetails[
skills = {Arcana +5, Insight +4, Persuasion +1},
languages = {Common, Draconic, Goblin},
challenge = 1/2,
]
\DndMonsterAction{Barkskin}
Sharwyn's AC can't be lower than 16.

\DndMonsterAction{Special Equipment}
Sharwyn has a \textit{spellbook} that contains the spells listed in her Spellcasting trait, plus \textit{detect magic} and \textit{silent image}.

\DndMonsterAction{Tree Thrall}
If the Gulthias Tree dies, Sharwyn dies 24 hours later.

\DndMonsterAction{Spellcasting}
Sharwyn is a 1st-level spellcaster. Her spellcasting ability is Intelligence (spell save DC 13, 5 to hit with spell attacks). She has the following wizard spells prepared:
\begin{DndMonsterSpells}
  \DndMonsterSpellLevel{\textit{light}, \textit{prestidigitation}, \textit{ray of frost}}
   \DndMonsterSpellLevel[1][2]{\textit{color spray}, \textit{magic missile}, \textit{shield}, \textit{sleep}}
\end{DndMonsterSpells}
\DndMonsterSection{Actions}
\DndMonsterAction{Dagger}
\textsl{Melee Weapon Attack:} 4, reach 5 ft., one target. \textsl{Hit:} 4 (1d4 + 2) piercing damage.

\end{DndMonster}



\begin{DndMonster}[float=t]{Sir Braford}
\DndMonsterType{Medium humanoid (human), neutral evil}
\DndMonsterBasics[
armor-class = {18 (\textit{chain mail})},
hit-points = {\DndDice{3d8 + 6}},
speed = {30 ft.}
]
\DndMonsterAbilityScores[
 str = 16,
 dex = 9,
 con = 14,
 int = 10,
 wis = 13,
 cha = 14
]
\DndMonsterDetails[
skills = {Athletics +5, Perception +3},
languages = {Common},
challenge = 1/2,
]
\DndMonsterAction{Barkskin}
Sir Braford's AC can't be lower than 16.

\DndMonsterAction{Special Equipment}
Sir Braford wields \textit{Shatterspike}, a magic longsword that grants a +1 bonus to attack and damage rolls made with it (included in his attack). See the Shatterspike handout for the item's other properties.

\DndMonsterAction{Tree Thrall}
If the Gulthias Tree dies, Sir Braford dies 24 hours later.

\DndMonsterSection{Actions}
\DndMonsterAction{Longsword}
\textsl{Melee Weapon Attack:} 6 to hit, reach 5 ft., one target. \textsl{Hit:} 8 (1d8 + 4) slashing damage, or 9 (1d10 + 4) slashing damage if used with two hands.

\DndMonsterSection{Reactions}
\DndMonsterAction{Protection}
When a creature Sir Braford can see attacks a target other than him that is within 5 feet of him, he can use a reaction to use his shield to impose disadvantage on the attack roll.

\end{DndMonster}



\begin{DndMonster}[float=t]{Siren}
\DndMonsterType{Medium fey, chaotic good}
\DndMonsterBasics[
armor-class = {14},
hit-points = {\DndDice{7d8 + 7}},
speed = {30 ft., swim 30 ft.}
]
\DndMonsterAbilityScores[
 str = 10,
 dex = 18,
 con = 12,
 int = 13,
 wis = 14,
 cha = 16
]
\DndMonsterDetails[
skills = {Medicine +4, Nature +3, Stealth +6, Survival +4},
senses = {darkvision 60 ft., passive perception 12},
languages = {Common, Elvish, Sylvan},
challenge = 3,
]
\DndMonsterAction{Amphibious}
Siren can breathe air and water.

\DndMonsterAction{Magic Resistance}
Siren has advantage on saving throws against spells and other magical effects.

\DndMonsterAction{Innate Spellcasting}
Siren's innate spellcasting ability is Charisma (spell save DC 13). She can innately cast the following spells, requiring no material components:
\begin{DndMonsterSpells}
  \DndInnateSpellLevel[1e]{\textit{charm person}, \textit{fog cloud}, \textit{greater invisibility}, \textit{polymorph} (self only)}
\end{DndMonsterSpells}
\DndMonsterSection{Actions}
\DndMonsterAction{Shortsword}
\textsl{Melee Weapon Attack:} 6 to hit, reach 5 ft., one target. \textsl{Hit:} 7 (1d6 + 4) piercing damage.

\DndMonsterAction{Stupefying Touch}
Siren touches one creature she can see within 5 feet of her. The creature must succeed on a DC 13 Intelligence saving throw or take 13 (3d6 + 3) psychic damage and be stunned until the start of Siren's next turn.

\end{DndMonster}



\begin{DndMonster}[float=t]{Snarla}
\DndMonsterType{Medium humanoid (human, shapechanger), chaotic evil}
\DndMonsterBasics[
armor-class = {11 in humanoid form (12 (natural armor) in wolf and hybrid forms)},
hit-points = {\DndDice{9d8 + 18}},
speed = {30 ft. (40 ft. in wolf form)}
]
\DndMonsterAbilityScores[
 str = 15,
 dex = 13,
 con = 14,
 int = 16,
 wis = 11,
 cha = 10
]
\DndMonsterDetails[
skills = {Perception +3},
damage-immunities = {bludgeoning, piercing, slashing from nonmagical attacks that aren't silvered},
languages = {Common (can't speak in wolf form)},
challenge = 5,
]
\DndMonsterAction{Shapechanger}
The werewolf can use its action to polymorph into a wolf-humanoid hybrid or into a wolf, or back into its true form, which is humanoid. Its statistics, other than its AC, are the same in each form. Any equipment it is wearing or carrying isn't transformed. It reverts to its true form if it dies.

\DndMonsterAction{Keen Hearing and Smell}
The werewolf has advantage on Wisdom (Perception) checks that rely on hearing or smell.

\DndMonsterAction{Spellcasting}
Snarla is a 6th-level spellcaster. Her spellcasting ability is Intelligence (spell save DC 14, +6 to hit with spell attacks). She has the following wizard spells prepared:
\begin{DndMonsterSpells}
  \DndMonsterSpellLevel{\textit{fire bolt}, \textit{light}, \textit{mage hand}, \textit{shocking grasp}}
   \DndMonsterSpellLevel[1][4]{\textit{magic missile}, \textit{shield}, \textit{thunderwave}}
   \DndMonsterSpellLevel[2][3]{\textit{mirror image}, \textit{web}}
   \DndMonsterSpellLevel[3][3]{\textit{dispel magic}, \textit{fear}, \textit{haste}, \textit{stinking cloud}}
\end{DndMonsterSpells}
\DndMonsterSection{Actions}
\DndMonsterAction{Multiattack (Humanoid or Hybrid Form Only)}
The werewolf makes two attacks: one with its bite and one with its claws or spear.

\DndMonsterAction{Bite (Wolf or Hybrid Form Only)}
\textsl{Melee Weapon Attack:} 4 to hit, reach 5 ft., one target. \textsl{Hit:} 6 (1d8 + 2) piercing damage. If the target is a humanoid, it must succeed on a DC 12 Constitution saving throw or be cursed with werewolf lycanthropy.

\DndMonsterAction{Claws (Hybrid Form Only)}
\textsl{Melee Weapon Attack:} 4 to hit, reach 5 ft., one creature. \textsl{Hit:} 7 (2d4 + 2) slashing damage.

\DndMonsterAction{Spear (Humanoid Form Only)}
\textsl{Melee or Ranged Weapon Attack:} 4 to hit, reach 5 ft. or range 20/60 ft., one creature. \textsl{Hit:} 5 (1d6 + 2) piercing damage, or 6 (1d8 + 2) piercing damage if used with two hands to make a melee attack.

\end{DndMonster}



\begin{DndMonster}[float*=b,width=\textwidth + 8pt]{Tarul Var}
\begin{multicols}{2}
\DndMonsterType{Medium undead, neutral evil}
\DndMonsterBasics[
armor-class = {16 (natural armor)},
hit-points = {\DndDice{14d8 + 42}},
speed = {30 ft.}
]
\DndMonsterAbilityScores[
 str = 11,
 dex = 16,
 con = 16,
 int = 19,
 wis = 14,
 cha = 16
]
\DndMonsterDetails[
saving-throws = {Con +8, Int +9, Wis +7},
skills = {Arcana +9, History +9, Insight +7, Perception +7},
damage-resistances = {cold; lightning; necrotic; bludgeoning, piercing, slashing from nonmagical attacks},
damage-immunities = {poison},
condition-immunities = {charmed, exhaustion, frightened, paralyzed, poisoned},
senses = {darkvision 60 ft., passive perception 17},
languages = {Abyssal, Common, Infernal, Primordial, Thayan},
challenge = 13,
]
\DndMonsterAction{Focused Conjuration}
While Var is concentration on a conjuration spell, his concentration can't be broken as a result of taking damage.

\DndMonsterAction{Legendary Resistance (3/Day)}
If Var fails a saving throw, he can choose to succeed instead.

\DndMonsterAction{Rejuvenation}
If Var is destroyed but his phylactery remains intact, Var gains a new body in 1d10 days, regaining all his hit points and becoming active again. The new body appears within 5 feet of the phylactery.

\DndMonsterAction{Turn Resistance}
Var has advantage on saving throws against any effect that turns undead.

\DndMonsterAction{Spellcasting}
Var is a 12th-level spellcaster. His spellcasting ability is Intelligence (spell save DC 17, 9 to hit with spell attacks). He has the following wizard spells prepared:
*Conjuration spell of 1st level or higher
\begin{DndMonsterSpells}
  \DndMonsterSpellLevel{\textit{fire bolt}, \textit{mage hand}, \textit{minor illusion}, \textit{prestidigitation}, \textit{ray of frost}}
   \DndMonsterSpellLevel[1][4]{\textit{detect magic}, \textit{magic missile}, \textit{shield}, \textit{unseen servant}*}
   \DndMonsterSpellLevel[2][3]{\textit{detect thoughts}, \textit{flaming sphere}*, \textit{mirror image}, \textit{scorching ray}}
   \DndMonsterSpellLevel[3][3]{\textit{counterspell}, \textit{dispel magic}, \textit{fireball}}
   \DndMonsterSpellLevel[4][3]{\textit{dimension door}*, \textit{Evard's black tentacles}*}
   \DndMonsterSpellLevel[5][3]{\textit{cloudkill}*, \textit{scrying}}
   \DndMonsterSpellLevel[6][1]{\textit{circle of death}}
\end{DndMonsterSpells}
\DndMonsterSection{Actions}
\DndMonsterAction{Paralyzing Touch}
\textsl{Melee Spell Attack:} 9 to hit, reach 5 ft., one creature. \textsl{Hit:} 10 (3d6) cold damage. The target must succeed on a DC 17 Constitution saving throw or be paralyzed for 1 minute. The target can repeat the saving throw at the end of each of its turns, ending the effect on itself on a success.

\DndMonsterAction{Benign Transposition}
Var teleports up to 30 feet to an unoccupied space he can see. Alternatively, he can choose a space within range that is occupied by a Small or Medium creature. If that creature is willing, both creatures teleport, swapping places. Var can use this feature again only after he finishes a long rest or casts a conjuration spell of 1st level or higher.

\DndMonsterSection{Legendary Actions}
Tarul Var can take 3 legendary actions, choosing from the options below. Only one legendary action can be used at a time and only at the end of another creature's turn. Tarul Var regains spent legendary actions at the start of its turn.

\begin{DndMonsterLegendaryActions}
\DndMonsterLegendaryAction{Cantrip}{Var casts a cantrip.
}
\DndMonsterLegendaryAction{Paralyzing Touch (Costs 2 Actions)}{Var uses Paralyzing Touch.
}
\DndMonsterLegendaryAction{Frightening Gaze (Costs 2 Actions)}{Var fixes his gaze on one creature he can see within 10 feet of him. The target must succeed on a DC 17 Wisdom saving throw against this magic or become frightened for 1 minute. The frightened target can repeat the saving throw at the end of each of its turns, ending the effect on itself on a success. If a target's saving throw is successful or the effect ends for it, the target is immune to Var's gaze for the next 24 hours.
}
\end{DndMonsterLegendaryActions}
\end{multicols}
\end{DndMonster}



\begin{DndMonster}[float=t]{Tecuziztecatl}
\DndMonsterType{Large monstrosity, neutral}
\DndMonsterBasics[
armor-class = {13 (natural armor)},
hit-points = {\DndDice{12d10 + 36}},
speed = {30 ft., climb 30 ft., swim 30 ft.}
]
\DndMonsterAbilityScores[
 str = 17,
 dex = 10,
 con = 16,
 int = 15,
 wis = 16,
 cha = 13
]
\DndMonsterDetails[
skills = {Deception +3, Stealth +2},
damage-resistances = {bludgeoning from nonmagical attacks},
damage-immunities = {acid},
senses = {blindsight 30 ft., passive perception 13},
languages = {Olman, Primordial},
challenge = 4,
]
\DndMonsterAction{Amphibious}
Tecuziztecatl can breathe air and water.

\DndMonsterAction{Glowing}
Tecuziztecatl sheds dim light within 20 feet of itself.

\DndMonsterAction{Flexible}
Tecuziztecatl can enter a space large enough for a Medium creature without squeezing.

\DndMonsterAction{Spider Climb}
Tecuziztecatl can climb difficult surfaces, including upside down on ceilings, without needing to make an ability check.

\DndMonsterSection{Actions}
\DndMonsterAction{Multiattack}
Tecuziztecatl makes two pseudopod attacks.

\DndMonsterAction{Pseudopod}
\textsl{Melee Weapon Attack:} 5 to hit, reach 10 ft., one target. \textsl{Hit:} 12 (2d8 + 3) bludgeoning damage.

\DndMonsterAction{Spit Acid (Recharge 4\textendash 6)}
Tecuziztecatl exhales acid in a 30-foot line that is 5 feet wide. Each creature in that line must make a DC 13 Dexterity saving throw, taking 18 (4d8) acid damage on a failed save, or half as much damage on a successful one.

\end{DndMonster}



\begin{DndMonster}[float=t]{Thayan Apprentice}
\DndMonsterType{Medium humanoid (human), any non-good alignment}
\DndMonsterBasics[
armor-class = {12 (15 with \textit{mage armor})},
hit-points = {\DndDice{5d8 + 5}},
speed = {30 ft.}
]
\DndMonsterAbilityScores[
 str = 10,
 dex = 14,
 con = 12,
 int = 15,
 wis = 13,
 cha = 11
]
\DndMonsterDetails[
skills = {Arcana +4},
languages = {Common, Thayan},
challenge = 2,
]
\DndMonsterAction{Doomvault Devotion}
Within the Doomvault, the apprentice has advantage on saving throws against being charmed or frightened.

\DndMonsterAction{Spellcasting}
The apprentice is a 4th-level spellcaster. Its spellcasting ability is Intelligence (spell save DC 12, 4 to hit with spell attacks). It has the following wizard spells prepared:
\begin{DndMonsterSpells}
  \DndMonsterSpellLevel{\textit{fire bolt}, \textit{mage hand}, \textit{prestidigitation}, \textit{shocking grasp}}
   \DndMonsterSpellLevel[1][4]{\textit{burning hands}, \textit{detect magic}, \textit{mage armor}, \textit{shield}}
   \DndMonsterSpellLevel[2][3]{\textit{blur}, \textit{scorching ray}}
\end{DndMonsterSpells}
\DndMonsterSection{Actions}
\DndMonsterAction{Dagger}
\textsl{Melee or Ranged Weapon Attack:} 4 to hit, reach 5 ft. or range 20/60 ft., one target. \textsl{Hit:} 4 (1d4 + 2) piercing damage.

\end{DndMonster}



\begin{DndMonster}[float=t]{Thayan Warrior}
\DndMonsterType{Medium humanoid (human), any non-good alignment}
\DndMonsterBasics[
armor-class = {16 (\textit{chain shirt})},
hit-points = {\DndDice{8d8 + 16}},
speed = {30 ft.}
]
\DndMonsterAbilityScores[
 str = 16,
 dex = 13,
 con = 14,
 int = 10,
 wis = 11,
 cha = 11
]
\DndMonsterDetails[
skills = {Perception +2},
languages = {Common, Thayan},
challenge = 2,
]
\DndMonsterAction{Doomvault Devotion}
Within the Doomvault, the warrior has advantage on saving throws against being charmed or frightened.

\DndMonsterAction{Pack Tactics}
The warrior has advantage on an attack roll against a creature if at least one of the warrior's allies is within 5 feet of the creature and the ally isn't incapacitated.

\DndMonsterSection{Actions}
\DndMonsterAction{Multiattack}
The warrior makes two melee attacks.

\DndMonsterAction{Longsword}
\textsl{Melee Weapon Attack:} 5 to hit, reach 5 ft., one target. \textsl{Hit:} 7 (1d8 + 3) slashing damage, or 8 (1d10 + 3) slashing damage if used with two hands.

\DndMonsterAction{Javelin}
\textsl{Melee or Ranged Weapon Attack:} 5 to hit, reach 5 ft. or range 30/120 ft., one target. \textsl{Hit:} 6 (1d6 + 3) piercing damage.

\end{DndMonster}



\begin{DndMonster}[float=t]{Thorn Slinger}
\DndMonsterType{Large plant, unaligned}
\DndMonsterBasics[
armor-class = {11},
hit-points = {\DndDice{5d10 + 5}},
speed = {10 ft.}
]
\DndMonsterAbilityScores[
 str = 13,
 dex = 12,
 con = 12,
 int = 1,
 wis = 10,
 cha = 1
]
\DndMonsterDetails[
condition-immunities = {blinded, deafened, frightened},
senses = {blindsight 60 ft. (blind beyond this radius), passive perception 10},
challenge = 1/2,
]
\DndMonsterAction{Adhesive Blossoms}
The thorn slinger adheres to anything that touches it. A Medium or smaller creature adhered to the thorn slinger is also grappled by it (escape DC 11). Ability checks made to escape this grapple have disadvantage.

At the end of each of the thorn slinger's turns, anything grappled by it takes 3 (1d6) acid damage.

\DndMonsterAction{False Appearance}
While the thorn slinger remains motionless, it is indistinguishable from an inanimate bush.

\DndMonsterSection{Actions}
\DndMonsterAction{Thorns}
\textsl{Melee or Ranged Weapon Attack:} 3 to hit, reach 5 ft. or range 30 ft., one target. \textsl{Hit:} 8 (2d6 + 1) piercing damage.

\end{DndMonster}



\begin{DndMonster}[float=t]{Transmuter}
\DndMonsterType{Medium humanoid (any race), any alignment}
\DndMonsterBasics[
armor-class = {12 (15 with \textit{mage armor})},
hit-points = {\DndDice{9d8}},
speed = {30 ft.}
]
\DndMonsterAbilityScores[
 str = 9,
 dex = 14,
 con = 11,
 int = 17,
 wis = 12,
 cha = 11
]
\DndMonsterDetails[
saving-throws = {Int +6, Wis +4},
skills = {Arcana +6, History +6},
languages = {any four languages},
challenge = 5,
]
\DndMonsterAction{Transmuter's Stone}
The transmuter carries a magic stone it crafted that grants its bearer one of the following effects:

Darkvision out to a range of 60 feet

An extra 10 feet of speed while the bearer is unencumbered

Proficiency with Constitution saving throws

Resistance to acid, cold, fire, lightning, or thunder damage (transmuter's choice whenever the transmuter chooses this benefit)

If the transmuter has the stone and casts a transmutation spell of 1st level or higher, it can change the effect of the stone.

\DndMonsterAction{Spellcasting}
The transmuter is a 9th-level spellcaster. Its spellcasting ability is Intelligence (spell save DC 14, 6 to hit with spell attacks). The transmuter has the following wizard spells prepared:
\begin{DndMonsterSpells}
  \DndMonsterSpellLevel{\textit{light}, \textit{mending}, \textit{prestidigitation}, \textit{ray of frost}}
   \DndMonsterSpellLevel[1][4]{\textit{chromatic orb}, \textit{expeditious retreat}*, \textit{mage armor}}
   \DndMonsterSpellLevel[2][3]{\textit{alter self}*, \textit{hold person}, \textit{knock}*}
   \DndMonsterSpellLevel[3][3]{\textit{blink}*, \textit{fireball}, \textit{slow}*}
   \DndMonsterSpellLevel[4][3]{\textit{polymorph}*, \textit{stoneskin}}
   \DndMonsterSpellLevel[5][1]{\textit{telekinesis}*}
\end{DndMonsterSpells}
\DndMonsterSection{Actions}
\DndMonsterAction{Quarterstaff}
\textsl{Melee Weapon Attack:} 2 to hit, reach 5 ft., one target. \textsl{Hit:} 2 (1d6 - 1) bludgeoning damage, or 3 (1d8 - 1) bludgeoning damage if used with two hands.

\end{DndMonster}



\begin{DndMonster}[float*=b,width=\textwidth + 8pt]{Vampire}
\begin{multicols}{2}
\DndMonsterType{Medium undead (shapechanger), lawful evil}
\DndMonsterBasics[
armor-class = {16 (natural armor)},
hit-points = {\DndDice{17d8 + 68}},
speed = {30 ft.}
]
\DndMonsterAbilityScores[
 str = 18,
 dex = 18,
 con = 18,
 int = 17,
 wis = 15,
 cha = 18
]
\DndMonsterDetails[
saving-throws = {Dex +9, Wis +7, Cha +9},
skills = {Perception +7, Stealth +9},
damage-resistances = {necrotic; bludgeoning, piercing, slashing from nonmagical attacks},
senses = {darkvision 120 ft., passive perception 17},
languages = {the languages it knew in life},
challenge = 13,
]
\DndMonsterAction{Shapechanger}
If the vampire isn't in sunlight or running water, it can use its action to polymorph into a Tiny bat or a Medium cloud of mist, or back into its true form.

While in bat form, the vampire can't speak, its walking speed is 5 feet, and it has a flying speed of 30 feet. Its statistics, other than its size and speed, are unchanged. Anything it is wearing transforms with it, but nothing it is carrying does. It reverts to its true form if it dies.

While in mist form, the vampire can't take any actions, speak, or manipulate objects. It is weightless, has a flying speed of 20 feet, can hover, and can enter a hostile creature's space and stop there. In addition, if air can pass through a space, the mist can do so without squeezing, and it can't pass through water. It has advantage on Strength, Dexterity, and Constitution saving throws, and it is immune to all nonmagical damage, except the damage it takes from sunlight.

\DndMonsterAction{Legendary Resistance (3/Day)}
If the vampire fails a saving throw, it can choose to succeed instead.

\DndMonsterAction{Misty Escape}
When it drops to 0 hit points outside its resting place, the vampire transforms into a cloud of mist (as in the Shapechanger trait) instead of falling unconscious, provided that it isn't in sunlight or running water. If it can't transform, it is destroyed.

While it has 0 hit points in mist form, it can't revert to its vampire form, and it must reach its resting place within 2 hours or be destroyed. Once in its resting place, it reverts to its vampire form. It is then paralyzed until it regains at least 1 hit point. After spending 1 hour in its resting place with 0 hit points, it regains 1 hit point.

\DndMonsterAction{Regeneration}
The vampire regains 20 hit points at the start of its turn if it has at least 1 hit point and isn't in sunlight or running water. If the vampire takes radiant damage or damage from holy water, this trait doesn't function at the start of the vampire's next turn.

\DndMonsterAction{Spider Climb}
The vampire can climb difficult surfaces, including upside down on ceilings, without needing to make an ability check.

\DndMonsterAction{Vampire Weaknesses}
The vampire has the following flaws:

\textit{Forbiddance.} The vampire can't enter a residence without an invitation from one of the occupants.

\textit{Harmed by Running Water.} The vampire takes 20 acid damage if it ends its turn in running water.

\textit{Stake to the Heart.} If a piercing weapon made of wood is driven into the vampire's heart while the vampire is incapacitated in its resting place, the vampire is paralyzed until the stake is removed.

\textit{Sunlight Hypersensitivity.} The vampire takes 20 radiant damage when it starts its turn in sunlight. While in sunlight, it has disadvantage on attack rolls and ability checks.

\DndMonsterSection{Actions}
\DndMonsterAction{Multiattack (Vampire Form Only)}
The vampire makes two attacks, only one of which can be a bite attack.

\DndMonsterAction{Unarmed Strike (Vampire Form Only)}
\textsl{Melee Weapon Attack:} 9 to hit, reach 5 ft., one creature. \textsl{Hit:} 8 (1d8 + 4) bludgeoning damage. Instead of dealing damage, the vampire can grapple the target (escape DC 18).

\DndMonsterAction{Bite (Bat or Vampire Form Only)}
\textsl{Melee Weapon Attack:} 9 to hit, reach 5 ft., one willing creature, or a creature that is grappled by the vampire, incapacitated, or restrained. \textsl{Hit:} 7 (1d6 + 4) piercing damage plus 10 (3d6) necrotic damage. The target's hit point maximum is reduced by an amount equal to the necrotic damage taken, and the vampire regains hit points equal to that amount. The reduction lasts until the target finishes a long rest. The target dies if this effect reduces its hit point maximum to 0. A humanoid slain in this way and then buried in the ground rises the following night as a \textbf{vampire spawn} under the vampire's control.

\DndMonsterAction{Charm}
The vampire targets one humanoid it can see within 30 feet of it. If the target can see the vampire, the target must succeed on a DC 17 Wisdom saving throw against this magic or be charmed by the vampire. The charmed target regards the vampire as a trusted friend to be heeded and protected. Although the target isn't under the vampire's control, it takes the vampire's requests or actions in the most favorable way it can, and it is a willing target for the vampire's bite attack.

Each time the vampire or the vampire's companions do anything harmful to the target, it can repeat the saving throw, ending the effect on itself on a success. Otherwise, the effect lasts 24 hours or until the vampire is destroyed, is on a different plane of existence than the target, or takes a bonus action to end the effect.

\DndMonsterAction{Children of the Night (1/Day)}
The vampire magically calls 2d4 swarms of bats or rats, provided that the sun isn't up. While outdoors, the vampire can call 3d6 wolves instead. The called creatures arrive in 1d4 rounds, acting as allies of the vampire and obeying its spoken commands. The beasts remain for 1 hour, until the vampire dies, or until the vampire dismisses them as a bonus action.

\DndMonsterSection{Legendary Actions}
Vampire can take 3 legendary actions, choosing from the options below. Only one legendary action can be used at a time and only at the end of another creature's turn. Vampire regains spent legendary actions at the start of its turn.

\begin{DndMonsterLegendaryActions}
\DndMonsterLegendaryAction{Move}{The vampire moves up to its speed without provoking opportunity attacks.
}
\DndMonsterLegendaryAction{Unarmed Strike}{The vampire makes one unarmed strike.
}
\DndMonsterLegendaryAction{Bite (Costs 2 Actions)}{The vampire makes one bite attack.
}
\end{DndMonsterLegendaryActions}
\DndMonsterSection{Regional Effects}
The region surrounding a vampire's lair is warped by the creature's unnatural presence, creating any of the following effects:


\begin{itemize}
\item There's a noticeable increase in the populations of bats, rats, and wolves in the region.
\item Plants within 500 feet of the lair wither, and their stems and branches become twisted and thorny.
\item Shadows cast within 500 feet of the lair seem abnormally gaunt and sometimes move as though alive.
\item A creeping fog clings to the ground within 500 feet of the vampire's lair. The fog occasionally takes eerie forms, such as grasping claws and writhing serpents.
\end{itemize}

If the vampire is destroyed, these effects end after 2d6 days.

\end{multicols}
\end{DndMonster}



\begin{DndMonster}[float=t]{Vampiric Mist}
\DndMonsterType{Medium undead, chaotic evil}
\DndMonsterBasics[
armor-class = {13},
hit-points = {\DndDice{4d8 + 12}},
speed = {0 ft., fly 30 ft. \textit{(hover)}}
]
\DndMonsterAbilityScores[
 str = 6,
 dex = 16,
 con = 16,
 int = 6,
 wis = 12,
 cha = 7
]
\DndMonsterDetails[
saving-throws = {Wis +3},
damage-resistances = {acid; cold; lightning; necrotic; thunder; bludgeoning, piercing, slashing from nonmagical attacks},
damage-immunities = {poison},
condition-immunities = {charmed, exhaustion, grappled, paralyzed, petrified, poisoned, prone, restrained},
senses = {darkvision 60 ft., passive perception 11},
challenge = 3,
]
\DndMonsterAction{Life Sense}
The mist can sense the location of any creature within 60 feet of it, unless that creature's type is construct or undead.

\DndMonsterAction{Forbiddance}
The mist can't enter a residence without an invitation from one of the occupants.

\DndMonsterAction{Misty Form}
The mist can occupy another creature's space and vice versa. In addition, if air can pass through a space, the mist can pass through it without squeezing. Each foot of movement in water costs it 2 extra feet, rather than 1 extra foot. The mist can't manipulate objects in any way that requires fingers or manual dexterity.

\DndMonsterAction{Sunlight Hypersensitivity}
The mist takes 10 radiant damage whenever it starts its turn in sunlight. While in sunlight, the mist has disadvantage on attack rolls and ability checks.

\DndMonsterSection{Actions}
\DndMonsterAction{Life Drain}
The mist touches one creature in its space. The target must succeed on a DC 13 Constitution saving throw (undead and constructs automatically succeed), or it takes 10 (2d6 + 3) necrotic damage, the mist regains 10 hit points, and the target's hit point maximum is reduced by an amount equal to the necrotic damage taken. This reduction lasts until the target finishes a long rest. The target dies if its hit point maximum is reduced to 0.

\end{DndMonster}


\clearpage

\begin{DndMonster}[float=t]{Veteran}
\DndMonsterType{Medium humanoid (any race), any alignment}
\DndMonsterBasics[
armor-class = {17 (\textit{splint armor})},
hit-points = {\DndDice{9d8 + 18}},
speed = {30 ft.}
]
\DndMonsterAbilityScores[
 str = 16,
 dex = 13,
 con = 14,
 int = 10,
 wis = 11,
 cha = 10
]
\DndMonsterDetails[
skills = {Athletics +5, Perception +2},
languages = {any one language (usually Common)},
challenge = 3,
]
\DndMonsterSection{Actions}
\DndMonsterAction{Multiattack}
The veteran makes two longsword attacks. If it has a shortsword drawn, it can also make a shortsword attack.

\DndMonsterAction{Longsword}
\textsl{Melee Weapon Attack:} 5 to hit, reach 5 ft., one target. \textsl{Hit:} 7 (1d8 + 3) slashing damage, or 8 (1d10 + 3) slashing damage if used with two hands.

\DndMonsterAction{Shortsword}
\textsl{Melee Weapon Attack:} 5 to hit, reach 5 ft., one target. \textsl{Hit:} 6 (1d6 + 3) piercing damage.

\DndMonsterAction{Heavy Crossbow}
\textsl{Ranged Weapon Attack:} 3 to hit, range 100/400 ft., one target. \textsl{Hit:} 6 (1d10 + 1) piercing damage.

\end{DndMonster}



\begin{DndMonster}[float=t]{Werewolf}
\DndMonsterType{Medium humanoid (human, shapechanger), chaotic evil}
\DndMonsterBasics[
armor-class = {11 in humanoid form (12 (natural armor) in wolf or hybrid form)},
hit-points = {\DndDice{9d8 + 18}},
speed = {30 ft. (40 ft. in wolf form)}
]
\DndMonsterAbilityScores[
 str = 15,
 dex = 13,
 con = 14,
 int = 10,
 wis = 11,
 cha = 10
]
\DndMonsterDetails[
skills = {Perception +4, Stealth +3},
damage-immunities = {bludgeoning, piercing, slashing from nonmagical attacks that aren't silvered},
languages = {Common (can't speak in wolf form)},
challenge = 3,
]
\DndMonsterAction{Shapechanger}
The werewolf can use its action to polymorph into a wolf-humanoid hybrid or into a wolf, or back into its true form, which is humanoid. Its statistics, other than its AC, are the same in each form. Any equipment it is wearing or carrying isn't transformed. It reverts to its true form if it dies.

\DndMonsterAction{Keen Hearing and Smell}
The werewolf has advantage on Wisdom (Perception) checks that rely on hearing or smell.

\DndMonsterSection{Actions}
\DndMonsterAction{Multiattack (Humanoid or Hybrid Form Only)}
The werewolf makes two attacks: two with its spear (humanoid form) or one with its bite and one with its claws (hybrid form).

\DndMonsterAction{Bite (Wolf or Hybrid Form Only)}
\textsl{Melee Weapon Attack:} 4 to hit, reach 5 ft., one target. \textsl{Hit:} 6 (1d8 + 2) piercing damage. If the target is a humanoid, it must succeed on a DC 12 Constitution saving throw or be cursed with werewolf lycanthropy.

\DndMonsterAction{Claws (Hybrid Form Only)}
\textsl{Melee Weapon Attack:} 4 to hit, reach 5 ft., one creature. \textsl{Hit:} 7 (2d4 + 2) slashing damage.

\DndMonsterAction{Spear (Humanoid Form Only)}
\textsl{Melee or Ranged Weapon Attack:} 4 to hit, reach 5 ft. or range 20/60 ft., one creature. \textsl{Hit:} 5 (1d6 + 2) piercing damage, or 6 (1d8 + 2) piercing damage if used with two hands to make a melee attack.

\end{DndMonster}



\begin{DndMonster}[float=t]{White Maw}
\DndMonsterType{Gargantuan ooze, chaotic neutral}
\DndMonsterBasics[
armor-class = {5},
hit-points = {\DndDice{14d20 + 70}},
speed = {10 ft.}
]
\DndMonsterAbilityScores[
 str = 18,
 dex = 1,
 con = 20,
 int = 12,
 wis = 10,
 cha = 3
]
\DndMonsterDetails[
damage-resistances = {acid, cold, fire},
damage-immunities = {poison},
condition-immunities = {blinded, charmed, deafened, exhaustion, frightened, poisoned, prone},
senses = {blindsight 60 ft. (blind beyond this radius), passive perception 10},
languages = {telepathy 50 ft.},
challenge = 10,
]
\DndMonsterAction{Amorphous Form}
White Maw can occupy another creature's space and vice versa.

\DndMonsterAction{Corrode Metal}
Any nonmagical weapon made of metal that hits White Maw corrodes. After dealing damage, the weapon takes a permanent and cumulative −1 penalty to damage rolls. if its penalty drops to −5, the weapon is destroyed. Nonmagical ammunition made of metal that hits White Maw is destroyed after dealing damage.

White Maw can eat through 2-inch-thick, nonmagical metal in 1 round.

\DndMonsterAction{False Appearance}
While White Maw remains motionless, it is indistinguishable from white stone.

\DndMonsterAction{Killer Response}
Any creature that starts its turn in White Maw's space is targeted by a pseudopod attack if White Maw isn't incapacitated.

\DndMonsterSection{Actions}
\DndMonsterAction{Pseudopod}
\textsl{Melee Weapon Attack:} 8 to hit, reach 10 ft., one target. \textsl{Hit:} 22 (4d8 + 4) bludgeoning damage plus 9 (2d8) acid damage. If the target is wearing nonmagical metal armor, its armor is partly corroded and takes a permanent and cumulative −1 penalty to the AC it offers. The armor is destroyed if the penalty reduces its AC to 10.

\end{DndMonster}



\begin{DndMonster}[float=t]{Wight}
\DndMonsterType{Medium undead, neutral evil}
\DndMonsterBasics[
armor-class = {14 (studded leather armor)},
hit-points = {\DndDice{6d8 + 18}},
speed = {30 ft.}
]
\DndMonsterAbilityScores[
 str = 15,
 dex = 14,
 con = 16,
 int = 10,
 wis = 13,
 cha = 15
]
\DndMonsterDetails[
skills = {Perception +3, Stealth +4},
damage-resistances = {necrotic; bludgeoning, piercing, slashing from nonmagical attacks that aren't silvered},
damage-immunities = {poison},
condition-immunities = {exhaustion, poisoned},
senses = {darkvision 60 ft., passive perception 13},
languages = {the languages it knew in life},
challenge = 3,
]
\DndMonsterAction{Sunlight Sensitivity}
While in sunlight, the wight has disadvantage on attack rolls, as well as on Wisdom (Perception) checks that rely on sight.

\DndMonsterSection{Actions}
\DndMonsterAction{Multiattack}
The wight makes two longsword attacks or two longbow attacks. It can use its Life Drain in place of one longsword attack.

\DndMonsterAction{Life Drain}
\textsl{Melee Weapon Attack:} 4 to hit, reach 5 ft., one creature. \textsl{Hit:} 5 (1d6 + 2) necrotic damage. The target must succeed on a DC 13 Constitution saving throw or its hit point maximum is reduced by an amount equal to the damage taken. This reduction lasts until the target finishes a long rest. The target dies if this effect reduces its hit point maximum to 0.

A humanoid slain by this attack rises 24 hours later as a \textbf{zombie} under the wight's control, unless the humanoid is restored to life or its body is destroyed. The wight can have no more than twelve zombies under its control at one time.

\DndMonsterAction{Longsword}
\textsl{Melee Weapon Attack:} 4 to hit, reach 5 ft., one target. \textsl{Hit:} 6 (1d8 + 2) slashing damage, or 7 (1d10 + 2) slashing damage if used with two hands.

\DndMonsterAction{Longbow}
\textsl{Ranged Weapon Attack:} 4 to hit, range 150/600 ft., one target. \textsl{Hit:} 6 (1d8 + 2) piercing damage.

\end{DndMonster}



\begin{DndMonster}[float=t]{Yusdrayl}
\DndMonsterType{Small humanoid (kobold), lawful evil}
\DndMonsterBasics[
armor-class = {12 (15 with \textit{mage armor})},
hit-points = {\DndDice{3d6 + 6}},
speed = {30 ft.}
]
\DndMonsterAbilityScores[
 str = 8,
 dex = 15,
 con = 14,
 int = 10,
 wis = 10,
 cha = 16
]
\DndMonsterDetails[
skills = {Arcana +2, Insight +2, Stealth +4},
senses = {darkvision 60 ft., passive perception 10},
languages = {Common, Draconic},
challenge = 1,
]
\DndMonsterAction{Sunlight Sensitivity}
While in sunlight, Yusdrayl has disadvantage on attack rolls, as well as on Wisdom (Perception) checks that rely on sight.

\DndMonsterAction{Pack Tactics}
Yusdrayl has advantage on an attack roll against a creature if at least one of her allies is within 5 feet of the creature and the ally isn't incapacitated.

\DndMonsterAction{Spellcasting}
Yusdrayl is a 2nd-level spellcaster. Her spellcasting ability is Charisma (spell save DC 13, 5 to hit with spell attacks). She knows the following sorcerer spells:
\begin{DndMonsterSpells}
  \DndMonsterSpellLevel{\textit{mage hand}, \textit{prestidigitation}, \textit{ray of frost}, \textit{shocking grasp}}
   \DndMonsterSpellLevel[1][4]{\textit{burning hands}, \textit{chromatic orb}, \textit{mage armor}}
\end{DndMonsterSpells}
\DndMonsterSection{Actions}
\DndMonsterAction{Dagger}
\textsl{Melee Weapon Attack:} 4 to hit, reach 5 ft., one target. \textsl{Hit:} 4 (1d4 + 2) piercing damage.

\end{DndMonster}


\chapter{Spells}
\addcontentsline{toc}{section}{Acid Splash}


\DndSpellHeader%
{Acid Splash}
{Conjuration cantrip}
{1 action}
{60 feet}
{V, S}
{Instantaneous}
You hurl a bubble of acid. Choose one creature you can see within range, or choose two creatures you can see within range that are within 5 feet of each other. A target must succeed on a Dexterity saving throw or take 1d6 acid damage.

This spell's damage increases by 1d6 when you reach 5th level (2d6), 11th level (3d6), and 17th level (4d6).

\addcontentsline{toc}{section}{Alter Self}


\DndSpellHeader%
{Alter Self}
{2nd level Transmutation}
{1 action}
{Self}
{V, S}
{Concentration, up to 1 hour}
You assume a different form. When you cast the spell, choose one of the following options, the effects of which last for the duration of the spell. While the spell lasts, you can end one option as an action to gain the benefits of a different one.

\subparagraph{Aquatic Adaptation}
You adapt your body to an aquatic environment, sprouting gills and growing webbing between your fingers. You can breathe underwater and gain a swimming speed equal to your walking speed.

\subparagraph{Change Appearance}
You transform your appearance. You decide what you look like, including your height, weight, facial features, sound of your voice, hair length, coloration, and distinguishing characteristics, if any. You can make yourself appear as a member of another race, though none of your statistics change. You also can't appear as a creature of a different size than you, and your basic shape stays the same; if you're bipedal, you can't use this spell to become quadrupedal, for instance. At any time for the duration of the spell, you can use your action to change your appearance in this way again.

\subparagraph{Natural Weapons}
You grow claws, fangs, spines, horns, or a different natural weapon of your choice. Your unarmed strikes deal 1d6 bludgeoning, piercing, or slashing damage, as appropriate to the natural weapon you chose, and you are proficient with your unarmed strikes. Finally, the natural weapon is magic and you have a +1 bonus to the attack and damage rolls you make using it.

\addcontentsline{toc}{section}{Animate Dead}


\DndSpellHeader%
{Animate Dead}
{3rd level Necromancy}
{1 minute}
{10 feet}
{V, S, a drop of blood, a piece of flesh, and a pinch of bone dust}
{Instantaneous}
This spell creates an undead servant. Choose a pile of bones or a corpse of a Medium or Small humanoid within range. Your spell imbues the target with a foul mimicry of life, raising it as an undead creature. The target becomes a \textbf{skeleton} if you chose bones or a \textbf{zombie} if you chose a corpse (the DM has the creature's game statistics).

On each of your turns, you can use a bonus action to mentally command any creature you made with this spell if the creature is within 60 feet of you (if you control multiple creatures, you can command any or all of them at the same time, issuing the same command to each one). You decide what action the creature will take and where it will move during its next turn, or you can issue a general command, such as to guard a particular chamber or corridor. If you issue no commands, the creature only defends itself against hostile creatures. Once given an order, the creature continues to follow it until its task is complete.

The creature is under your control for 24 hours, after which it stops obeying any command you've given it. To maintain control of the creature for another 24 hours, you must cast this spell on the creature again before the current 24-hour period ends. This use of the spell reasserts your control over up to four creatures you have animated with this spell, rather than animating a new one.

\addcontentsline{toc}{section}{Antimagic Field}


\DndSpellHeader%
{Antimagic Field}
{8th level Abjuration}
{1 action}
{Self (10 feet sphere)}
{V, S, a pinch of powdered iron or iron filings}
{Concentration, up to 1 hour}
A 10-foot-radius invisible sphere of antimagic surrounds you. This area is divorced from the magical energy that suffuses the multiverse. Within the sphere, spells can't be cast, summoned creatures disappear, and even magic items become mundane. Until the spell ends, the sphere moves with you, centered on you.

Spells and other magical effects, except those created by an artifact or a deity, are suppressed in the sphere and can't protrude into it. A slot expended to cast a suppressed spell is consumed. While an effect is suppressed, it doesn't function, but the time it spends suppressed counts against its duration.

\subparagraph{Targeted Effects}
Spells and other magical effects, such as \textit{magic missile} and \textit{charm person}, that target a creature or an object in the sphere have no effect on that target.

\subparagraph{Areas of Magic}
The area of another spell or magical effect, such as \textit{fireball}, can't extend into the sphere. If the sphere overlaps an area of magic, the part of the area that is covered by the sphere is suppressed. For example, the flames created by a wall of fire are suppressed within the sphere, creating a gap in the wall if the overlap is large enough.

\subparagraph{Spells}
Any active spell or other magical effect on a creature or an object in the sphere is suppressed while the creature or object is in it.

\subparagraph{Magic Items}
The properties and powers of magic items are suppressed in the sphere. For example, a +1 longsword in the sphere functions as a nonmagical longsword.

A magic weapon's properties and powers are suppressed if it is used against a target in the sphere or wielded by an attacker in the sphere. If a magic weapon or a piece of magic ammunition fully leaves the sphere (for example, if you fire a magic arrow or throw a magic spear at a target outside the sphere), the magic of the item ceases to be suppressed as soon as it exits.

\subparagraph{Magical Travel}
Teleportation and planar travel fail to work in the sphere, whether the sphere is the destination or the departure point for such magical travel. A portal to another location, world, or plane of existence, as well as an opening to an extradimensional space such as that created by the rope trick spell, temporarily closes while in the sphere.

\subparagraph{Creatures and Objects}
A creature or object summoned or created by magic temporarily winks out of existence in the sphere. Such a creature instantly reappears once the space the creature occupied is no longer within the sphere.

\subparagraph{Dispel Magic}
Spells and magical effects such as \textit{dispel magic} have no effect on the sphere. Likewise, the spheres created by different \textit{antimagic field} spells don't nullify each other.

\addcontentsline{toc}{section}{Banishment}


\DndSpellHeader%
{Banishment}
{4th level Abjuration}
{1 action}
{60 feet}
{V, S, an item distasteful to the target}
{Concentration, up to 1 minute}
You attempt to send one creature that you can see within range to another plane of existence. The target must succeed on a Charisma saving throw or be banished.

If the target is native to the plane of existence you're on, you banish the target to a harmless demiplane. While there, the target is incapacitated. The target remains there until the spell ends, at which point the target reappears in the space it left or in the nearest unoccupied space if that space is occupied.

If the target is native to a different plane of existence than the one you're on, the target is banished with a faint popping noise, returning to its home plane. If the spell ends before 1 minute has passed, the target reappears in the space it left or in the nearest unoccupied space if that space is occupied. Otherwise, the target doesn't return.

\addcontentsline{toc}{section}{Bestow Curse}


\DndSpellHeader%
{Bestow Curse}
{3rd level Necromancy}
{1 action}
{Touch}
{V, S}
{Concentration, up to 1 minute}
You touch a creature, and that creature must succeed on a Wisdom saving throw or become cursed for the duration of the spell. When you cast this spell, choose the nature of the curse from the following options:


\begin{itemize}
\item Choose one ability score. While cursed, the target has disadvantage on ability checks and saving throws made with that ability score.
\item While cursed, the target has disadvantage on attack rolls against you.
\item While cursed, the target must make a Wisdom saving throw at the start of each of its turns. If it fails, it wastes its action that turn doing nothing.
\item While the target is cursed, your attacks and spells deal an extra 1d8 necrotic damage to the target.
\end{itemize}

A \textit{remove curse} spell ends this effect. At the DM's option, you may choose an alternative curse effect, but it should be no more powerful than those described above. The DM has final say on such a curse's effect.

\addcontentsline{toc}{section}{Bigby's Hand}


\DndSpellHeader%
{Bigby's Hand}
{5th level Evocation}
{1 action}
{120 feet}
{V, S, an eggshell and a snakeskin glove}
{Concentration, up to 1 minute}
You create a Large hand of shimmering, translucent force in an unoccupied space that you can see within range. The hand lasts for the spell's duration, and it moves at your command, mimicking the movements of your own hand.

The hand is an object that has AC 20 and hit points equal to your hit point maximum. If it drops to 0 hit points, the spell ends. It has a Strength of 26 (8) and a Dexterity of 10 (0). The hand doesn't fill its space.

When you cast the spell and as a bonus action on your subsequent turns, you can move the hand up to 60 feet and then cause one of the following effects with it.

\subparagraph{Clenched Fist}
The hand strikes one creature or object within 5 feet of it. Make a melee spell attack for the hand using your game statistics. On a hit, the target takes 4d8 force damage.

\subparagraph{Forceful Hand}
The hand attempts to push a creature within 5 feet of it in a direction you choose. Make a check with the hand's Strength contested by the Strength (Athletics) check of the target. If the target is Medium or smaller, you have advantage on the check. If you succeed, the hand pushes the target up to 5 feet plus a number of feet equal to five times your spellcasting ability modifier. The hand moves with the target to remain within 5 feet of it.

\subparagraph{Grasping Hand}
The hand attempts to grapple a Huge or smaller creature within 5 feet of it. You use the hand's Strength score to resolve the grapple. If the target is Medium or smaller, you have advantage on the check. While the hand is grappling the target, you can use a bonus action to have the hand crush it. When you do so, the target takes bludgeoning damage equal to 2d6 + your spellcasting ability modifier.

\subparagraph{Interposing Hand}
The hand interposes itself between you and a creature you choose until you give the hand a different command. The hand moves to stay between you and the target, providing you with half cover against the target. The target can't move through the hand's space if its Strength score is less than or equal to the hand's Strength score. If its Strength score is higher than the hand's Strength score, the target can move toward you through the hand's space, but that space is difficult terrain for the target.

\addcontentsline{toc}{section}{Blight}


\DndSpellHeader%
{Blight}
{4th level Necromancy}
{1 action}
{30 feet}
{V, S}
{Instantaneous}
Necromantic energy washes over a creature of your choice that you can see within range, draining moisture and vitality from it. The target must make a Constitution saving throw. The target takes 8d8 necrotic damage on a failed save, or half as much damage on a successful one. This spell has no effect on undead or constructs.

If you target a plant creature or a magical plant, it makes the saving throw with disadvantage, and the spell deals maximum damage to it.

If you target a nonmagical plant that isn't a creature, such as a tree or shrub, it doesn't make a saving throw, it simply withers and dies.

\addcontentsline{toc}{section}{Blindness/Deafness}


\DndSpellHeader%
{Blindness/Deafness}
{2nd level Necromancy}
{1 action}
{30 feet}
{V}
{1 minute}
You can blind or deafen a foe. Choose one creature that you can see within range to make a Constitution saving throw. If it fails, the target is either blinded or deafened (your choice) for the duration. At the end of each of its turns, the target can make a Constitution saving throw. On a success, the spell ends.

\addcontentsline{toc}{section}{Blink}


\DndSpellHeader%
{Blink}
{3rd level Transmutation}
{1 action}
{Self}
{V, S}
{1 minute}
Roll a d20 at the end of each of your turns for the duration of the spell. On a roll of 11 or higher, you vanish from your current plane of existence and appear in the Ethereal Plane (the spell fails and the casting is wasted if you were already on that plane). At the start of your next turn, and when the spell ends if you are on the Ethereal Plane, you return to an unoccupied space of your choice that you can see within 10 feet of the space you vanished from. If no unoccupied space is available within that range, you appear in the nearest unoccupied space (chosen at random if more than one space is equally near). You can dismiss this spell as an action.

While on the Ethereal Plane, you can see and hear the plane you originated from, which is cast in shades of gray, and you can't see anything there more than 60 feet away. You can only affect and be affected by other creatures on the Ethereal Plane. Creatures that aren't there can't perceive you or interact with you, unless they have the ability to do so.

\addcontentsline{toc}{section}{Blur}


\DndSpellHeader%
{Blur}
{2nd level Illusion}
{1 action}
{Self}
{V}
{Concentration, up to 1 minute}
Your body becomes blurred, shifting and wavering to all who can see you. For the duration, any creature has disadvantage on attack rolls against you. An attacker is immune to this effect if it doesn't rely on sight, as with blindsight, or can see through illusions, as with truesight.

\addcontentsline{toc}{section}{Burning Hands}


\DndSpellHeader%
{Burning Hands}
{1st level Evocation}
{1 action}
{Self (15 feet cone)}
{V, S}
{Instantaneous}
As you hold your hands with thumbs touching and fingers spread, a thin sheet of flames shoots forth from your outstretched fingertips. Each creature in a 15-foot cone must make a Dexterity saving throw. A creature takes 3d6 fire damage on a failed save, or half as much damage on a successful one.

The fire ignites any flammable objects in the area that aren't being worn or carried.

\addcontentsline{toc}{section}{Chain Lightning}


\DndSpellHeader%
{Chain Lightning}
{6th level Evocation}
{1 action}
{150 feet}
{V, S, a bit of fur; a piece of amber, glass, or a crystal rod; and three silver pins}
{Instantaneous}
You create a bolt of lightning that arcs toward a target of your choice that you can see within range. Three bolts then leap from that target to as many as three other targets, each of which must be within 30 feet of the first target. A target can be a creature or an object and can be targeted by only one of the bolts.

A target must make a Dexterity saving throw. The target takes 10d8 lightning damage on a failed save, or half as much damage on a successful one.

\addcontentsline{toc}{section}{Charm Person}


\DndSpellHeader%
{Charm Person}
{1st level Enchantment}
{1 action}
{30 feet}
{V, S}
{1 hour}
You attempt to charm a humanoid you can see within range. It must make a Wisdom saving throw, and does so with advantage if you or your companions are fighting it. If it fails the saving throw, it is charmed by you until the spell ends or until you or your companions do anything harmful to it. The charmed creature regards you as a friendly acquaintance. When the spell ends, the creature knows it was charmed by you.

\addcontentsline{toc}{section}{Chill Touch}


\DndSpellHeader%
{Chill Touch}
{Necromancy cantrip}
{1 action}
{120 feet}
{V, S}
{1 round}
You create a ghostly, skeletal hand in the space of a creature within range. Make a ranged spell attack against the creature to assail it with the chill of the grave. On a hit, the target takes 1d8 necrotic damage, and it can't regain hit points until the start of your next turn. Until then, the hand clings to the target.

If you hit an undead target, it also has disadvantage on attack rolls against you until the end of your next turn.

This spell's damage increases by 1d8 when you reach 5th level (2d8), 11th level (3d8), and 17th level (4d8).

\addcontentsline{toc}{section}{Chromatic Orb}


\DndSpellHeader%
{Chromatic Orb}
{1st level Evocation}
{1 action}
{90 feet}
{V, S, a diamond worth at least 50 gp}
{Instantaneous}
You hurl a 4-inch-diameter sphere of energy at a creature that you can see within range. You choose acid, cold, fire, lightning, poison, or thunder for the type of orb you create, and then make a ranged spell attack against the target. If the attack hits, the creature takes 3d8 damage of the type you chose.

\addcontentsline{toc}{section}{Circle of Death}


\DndSpellHeader%
{Circle of Death}
{6th level Necromancy}
{1 action}
{150 feet}
{V, S, the powder of a crushed black pearl worth at least 500 gp}
{Instantaneous}
A sphere of negative energy ripples out in a 60-foot-radius sphere from a point within range. Each creature in that area must make a Constitution saving throw. A target takes 8d6 necrotic damage on a failed save, or half as much damage on a successful one.

\addcontentsline{toc}{section}{Clairvoyance}


\DndSpellHeader%
{Clairvoyance}
{3rd level Divination}
{10 minute}
{1 mile}
{V, S, a focus worth at least 100 gp, either a jeweled horn for hearing or a glass eye for seeing}
{Concentration, up to 10 minutes}
You create an invisible sensor within range in a location familiar to you (a place you have visited or seen before) or in an obvious location that is unfamiliar to you (such as behind a door, around a corner, or in a grove of trees). The sensor remains in place for the duration, and it can't be attacked or otherwise interacted with.

When you cast the spell, you choose seeing or hearing. You can use the chosen sense through the sensor as if you were in its space. As your action, you can switch between seeing and hearing.

A creature that can see the sensor (such as a creature benefiting from \textit{see invisibility} or truesight) sees a luminous, intangible orb about the size of your fist.

\addcontentsline{toc}{section}{Cloud of Daggers}


\DndSpellHeader%
{Cloud of Daggers}
{2nd level Conjuration}
{1 action}
{60 feet}
{V, S, a sliver of glass}
{Concentration, up to 1 minute}
You fill the air with spinning daggers in a cube 5 feet on each side, centered on a point you choose within range. A creature takes 4d4 slashing damage when it enters the spell's area for the first time on a turn or starts its turn there.

\addcontentsline{toc}{section}{Cloudkill}


\DndSpellHeader%
{Cloudkill}
{5th level Conjuration}
{1 action}
{120 feet}
{V, S}
{Concentration, up to 10 minutes}
You create a 20-foot-radius sphere of poisonous, yellow-green fog centered on a point you choose within range. The fog spreads around corners. It lasts for the duration or until strong wind disperses the fog, ending the spell. Its area is Vision and Light.

When a creature enters the spell's area for the first time on a turn or starts its turn there, that creature must make a Constitution saving throw. The creature takes 5d8 poison damage on a failed save, or half as much damage on a successful one. Creatures are affected even if they hold their breath or don't need to breathe.

The fog moves 10 feet away from you at the start of each of your turns, rolling along the surface of the ground. The vapors, being heavier than air, sink to the lowest level of the land, even pouring down openings.

\addcontentsline{toc}{section}{Color Spray}


\DndSpellHeader%
{Color Spray}
{1st level Illusion}
{1 action}
{Self (15 feet cone)}
{V, S, a pinch of powder or sand that is colored red, yellow, and blue}
{1 round}
A dazzling array of flashing, colored light springs from your hand. Roll 6d10; the total is how many hit points of creatures this spell can effect. Creatures in a 15-foot cone originating from you are affected in ascending order of their current hit points (ignoring unconscious creatures and creatures that can't see).

Starting with the creature that has the lowest current hit points, each creature affected by this spell is blinded until the end of your next turn. Subtract each creature's hit points from the total before moving on to the creature with the next lowest hit points. A creature's hit points must be equal to or less than the remaining total for that creature to be affected.

\addcontentsline{toc}{section}{Cone of Cold}


\DndSpellHeader%
{Cone of Cold}
{5th level Evocation}
{1 action}
{Self (60 feet cone)}
{V, S, a small crystal or glass cone}
{Instantaneous}
A blast of cold air erupts from your hands. Each creature in a 60-foot cone must make a Constitution saving throw. A creature takes 8d8 cold damage on a failed save, or half as much damage on a successful one.

A creature killed by this spell becomes a frozen statue until it thaws.

\addcontentsline{toc}{section}{Confusion}


\DndSpellHeader%
{Confusion}
{4th level Enchantment}
{1 action}
{90 feet}
{V, S, three nut shells}
{Concentration, up to 1 minute}
This spell assaults and twists creatures' minds, spawning delusions and provoking uncontrolled action. Each creature in a 10-foot-radius sphere centered on a point you choose within range must succeed on a Wisdom saving throw when you cast this spell or be affected by it.

An affected target can't take reactions and must roll a d10 at the start of each of its turns to determine its behavior for that turn.

\begin{DndTable}[header=Confusion Behavior]{cX}

d10 & Behavior\\
1 & The creature uses all its movement to move in a random direction. To determine the direction, roll a d8 and assign a direction to each die face. The creature doesn't take an action this turn.\\
2—6 & The creature doesn't move or take actions this turn.\\
7—8 & The creature uses its action to make a melee attack against a randomly determined creature within its reach. If there is no creature within its reach, the creature does nothing this turn.\\
9—10 & The creature can act and move normally.\\
\end{DndTable}

At the end of each of its turns, an affected target can make a Wisdom saving throw. If it succeeds, this effect ends for that target.

\addcontentsline{toc}{section}{Conjure Elemental}


\DndSpellHeader%
{Conjure Elemental}
{5th level Conjuration}
{1 minute}
{90 feet}
{V, S, burning incense for air, soft clay for earth, sulfur and phosphorus for fire, or water and sand for water}
{Concentration, up to 1 hour}
You call forth an elemental servant. Choose an area of air, earth, fire, or water that fills a 10-foot cube within range. An elemental of challenge rating 5 or lower appropriate to the area you chose appears in an unoccupied space within 10 feet of it. For example, a \textbf{fire elemental} emerges from a bonfire, and an \textbf{earth elemental} rises up from the ground. The elemental disappears when it drops to 0 hit points or when the spell ends.

The elemental is friendly to you and your companions for the duration. Roll initiative for the elemental, which has its own turns. It obeys any verbal commands that you issue to it (no action required by you). If you don't issue any commands to the elemental, it defends itself from hostile creatures but otherwise takes no actions.

If your concentration is broken, the elemental doesn't disappear. Instead, you lose control of the elemental, it becomes hostile toward you and your companions, and it might attack. An uncontrolled elemental can't be dismissed by you, and it disappears 1 hour after you summoned it.

The DM has the elemental's statistics.

\addcontentsline{toc}{section}{Conjure Minor Elementals}


\DndSpellHeader%
{Conjure Minor Elementals}
{4th level Conjuration}
{1 minute}
{90 feet}
{V, S}
{Concentration, up to 1 hour}
You summon elementals that appear in unoccupied spaces that you can see within range. You choose one the following options for what appears:


\begin{itemize}
\item One elemental of challenge rating 2 or lower
\item Two elementals of challenge rating 1 or lower
\item Four elementals of challenge rating 1/2 or lower
\item Eight elementals of challenge rating 1/4 or lower.
\end{itemize}

An elemental summoned by this spell disappears when it drops to 0 hit points or when the spell ends.

The summoned creatures are friendly to you and your companions. Roll initiative for the summoned creatures as a group, which has its own turns. They obey any verbal commands that you issue to them (no action required by you). If you don't issue any commands to them, they defend themselves from hostile creatures, but otherwise take no actions.

The DM has the creatures' statistics.

\addcontentsline{toc}{section}{Control Water}


\DndSpellHeader%
{Control Water}
{4th level Transmutation}
{1 action}
{300 feet}
{V, S, a drop of water and a pinch of dust}
{Concentration, up to 10 minutes}
Until the spell ends, you control any freestanding water inside an area you choose that is a cube up to 100 feet on a side. You can choose from any of the following effects when you cast this spell. As an action on your turn, you can repeat the same effect or choose a different one.

\subparagraph{Flood}
You cause the water level of all standing water in the area to rise by as much as 20 feet. If the area includes a shore, the flooding water spills over onto dry land.

If you choose an area in a large body of water, you instead create a 20-foot tall wave that travels from one side of the area to the other and then crashes down. Any Huge or smaller vehicles in the wave's path are carried with it to the other side. Any Huge or smaller vehicles struck by the wave have a No effect chance of capsizing.

The water level remains elevated until the spell ends or you choose a different effect. If this effect produced a wave, the wave repeats on the start of your next turn while the flood effect lasts.

\subparagraph{Part Water}
You cause water in the area to move apart and create a trench. The trench extends across the spell's area, and the separated water forms a wall to either side. The trench remains until the spell ends or you choose a different effect. The water then slowly fills in the trench over the course of the next round until the normal water level is restored.

\subparagraph{Redirect Flow}
You cause flowing water in the area to move in a direction you choose, even if the water has to flow over obstacles, up walls, or in other unlikely directions. The water in the area moves as you direct it, but once it moves beyond the spell's area, it resumes its flow based on the terrain conditions. The water continues to move in the direction you chose until the spell ends or you choose a different effect.

\subparagraph{Whirlpool}
This effect requires a body of water at least 50 feet square and 25 feet deep. You cause a whirlpool to form in the center of the area. The whirlpool forms a vortex that is 5 feet wide at the base, up to 50 feet wide at the top, and 25 feet tall. Any creature or object in the water and within 25 feet of the vortex is pulled 10 feet toward it. A creature can swim away from the vortex by making a Strength (Athletics) check against your spell save DC.

When a creature enters the vortex for the first time on a turn or starts its turn there, it must make a Strength saving throw. On a failed save, the creature takes 2d8 bludgeoning damage and is caught in the vortex until the spell ends. On a successful save, the creature takes half damage, and isn't caught in the vortex. A creature caught in the vortex can use its action to try to swim away from the vortex as described above, but has disadvantage on the Strength (Athletics) check to do so.

The first time each turn that an object enters the vortex, the object takes 2d8 bludgeoning damage; this damage occurs each round it remains in the vortex.

\addcontentsline{toc}{section}{Control Weather}


\DndSpellHeader%
{Control Weather}
{8th level Transmutation}
{10 minute}
{Self (5 miles radius)}
{V, S, burning incense and bits of earth and wood mixed in water}
{Concentration, up to 8 hours}
You take control of the weather within 5 miles of you for the duration. You must be outdoors to cast this spell. Moving to a place where you don't have a clear path to the sky ends the spell early.

When you cast the spell, you change the current weather conditions, which are determined by the DM based on the climate and season. You can change precipitation, temperature, and wind. It takes 1d4 × 10 minutes for the new conditions to take effect. Once they do so, you can change the conditions again. When the spell ends, the weather gradually returns to normal.

When you change the weather conditions, find a current condition on the following tables and change its stage by one, up or down. When changing the wind, you can change its direction.

\begin{DndTable}[header=Precipitation]{cX}

Stage & Condition\\
1 & Clear\\
2 & Light clouds\\
3 & Overcast or ground fog\\
4 & Rain, hail, or snow\\
5 & Torrential rain, driving hail, or blizzard\\
\end{DndTable}

\begin{DndTable}[header=Temperature]{cX}

Stage & Condition\\
1 & Unbearable heat\\
2 & Hot\\
3 & Warm\\
4 & Cool\\
5 & Cold\\
6 & Arctic cold\\
\end{DndTable}

\begin{DndTable}[header=Wind]{cX}

Stage & Condition\\
1 & Calm\\
2 & Moderate wind\\
3 & Strong wind\\
4 & Gale\\
5 & Storm\\
\end{DndTable}

\addcontentsline{toc}{section}{Counterspell}


\DndSpellHeader%
{Counterspell}
{3rd level Abjuration}
{1 reaction}
{60 feet}
{S}
{Instantaneous}
You attempt to interrupt a creature in the process of casting a spell. If the creature is casting a spell of 3rd level or lower, its spell fails and has no effect. If it is casting a spell of 4th level or higher, make an ability check using your spellcasting ability. The DC equals 10 + the spell's level. On a success, the creature's spell fails and has no effect.

\addcontentsline{toc}{section}{Dancing Lights}


\DndSpellHeader%
{Dancing Lights}
{Evocation cantrip}
{1 action}
{120 feet}
{V, S, a bit of phosphorus or wychwood, or a glowworm}
{Concentration, up to 1 minute}
You create up to four torch-sized lights within range, making them appear as torches, lanterns, or glowing orbs that hover in the air for the duration. You can also combine the four lights into one glowing vaguely humanoid form of Medium size. Whichever form you choose, each light sheds dim light in a 10-foot radius.

As a bonus action on your turn, you can move the lights up to 60 feet to a new spot within range. A light must be within 20 feet of another light created by this spell, and a light winks out if it exceeds the spell's range.

\addcontentsline{toc}{section}{Darkness}


\DndSpellHeader%
{Darkness}
{2nd level Evocation}
{1 action}
{60 feet}
{V, bat fur and a drop of pitch or piece of coal}
{Concentration, up to 10 minutes}
Magical darkness spreads from a point you choose within range to fill a 15-foot-radius sphere for the duration. The darkness spreads around corners. A creature with darkvision can't see through this darkness, and nonmagical light can't illuminate it.

If the point you choose is on an object you are holding or one that isn't being worn or carried, the darkness emanates from the object and moves with it. Completely covering the source of the darkness with an opaque object, such as a bowl or a helm, blocks the darkness.

If any of this spell's area overlaps with an area of light created by a spell of 2nd level or lower, the spell that created the light is dispelled.

\addcontentsline{toc}{section}{Detect Evil and Good}


\DndSpellHeader%
{Detect Evil and Good}
{1st level Divination}
{1 action}
{Self}
{V, S}
{Concentration, up to 10 minutes}
For the duration, you know if there is an aberration, celestial, elemental, fey, fiend, or undead within 30 feet of you, as well as where the creature is located. Similarly, you know if there is a place or object within 30 feet of you that has been magically consecrated or desecrated.

The spell can penetrate most barriers, but it is blocked by 1 foot of stone, 1 inch of common metal, a thin sheet of lead, or 3 feet of wood or dirt.

\addcontentsline{toc}{section}{Detect Magic}


\DndSpellHeader%
{Detect Magic}
{1st level Divination}
{1 action}
{Self}
{V, S}
{Concentration, up to 10 minutes}
For the duration, you sense the presence of magic within 30 feet of you. If you sense magic in this way, you can use your action to see a faint aura around any visible creature or object in the area that bears magic, and you learn its \textit{school of magic}, if any.

The spell can penetrate most barriers, but it is blocked by 1 foot of stone, 1 inch of common metal, a thin sheet of lead, or 3 feet of wood or dirt.

\addcontentsline{toc}{section}{Detect Thoughts}


\DndSpellHeader%
{Detect Thoughts}
{2nd level Divination}
{1 action}
{Self}
{V, S, a copper piece}
{Concentration, up to 1 minute}
For the duration, you can read the thoughts of certain creatures. When you cast the spell and as your action on each turn until the spell ends, you can focus your mind on any one creature that you can see within 30 feet of you. If the creature you choose has an Intelligence of 3 or lower or doesn't speak any language, the creature is unaffected.

You initially learn the surface thoughts of the creature—what is most on its mind in that moment. As an action, you can either shift your attention to another creature's thoughts or attempt to probe deeper into the same creature's mind. If you probe deeper, the target must make a Wisdom saving throw. If it fails, you gain insight into its reasoning (if any), its emotional state, and something that looms large in its mind (such as something it worries over, loves, or hates). If it succeeds, the spell ends. Either way, the target knows that you are probing into its mind, and unless you shift your attention to another creature's thoughts, the creature can use its action on its turn to make an Intelligence check contested by your Intelligence check; if it succeeds, the spell ends.

Questions verbally directed at the target creature naturally shape the course of its thoughts, so this spell is particularly effective as part of an interrogation.

You can also use this spell to detect the presence of thinking creatures you can't see. When you cast the spell or as your action during the duration, you can search for thoughts within 30 feet of you. The spell can penetrate barriers, but 2 feet of rock, 2 inches of any metal other than lead, or a thin sheet of lead blocks you. You can't detect a creature with an Intelligence of 3 or lower or one that doesn't speak any language.

Once you detect the presence of a creature in this way, you can read its thoughts for the rest of the duration as described above, even if you can't see it, but it must still be within range.

\addcontentsline{toc}{section}{Dimension Door}


\DndSpellHeader%
{Dimension Door}
{4th level Conjuration}
{1 action}
{500 feet}
{V}
{Instantaneous}
You teleport yourself from your current location to any other spot within range. You arrive at exactly the spot desired. It can be a place you can see, one you can visualize, or one you can describe by stating distance and direction, such as "200 feet straight downward" or "upward to the northwest at a 45-degree angle, 300 feet."

You can bring along objects as long as their weight doesn't exceed what you can carry. You can also bring one willing creature of your size or smaller who is carrying gear up to its carrying capacity. The creature must be within 5 feet of you when you cast this spell.

If you would arrive in a place already occupied by an object or a creature, you and any creature traveling with you each take 4d6 force damage, and the spell fails to teleport you.

\addcontentsline{toc}{section}{Disguise Self}


\DndSpellHeader%
{Disguise Self}
{1st level Illusion}
{1 action}
{Self}
{V, S}
{1 hour}
You make yourself—including your clothing, armor, weapons, and other belongings on your person—look different until the spell ends or until you use your action to dismiss it. You can seem 1 foot shorter or taller and can appear thin, fat, or in between. You can't change your body type, so you must adopt a form that has the same basic arrangement of limbs. Otherwise, the extent of the illusion is up to you.

The changes wrought by this spell fail to hold up to physical inspection. For example, if you use this spell to add a hat to your outfit, objects pass through the hat, and anyone who touches it would feel nothing or would feel your head and hair. If you use this spell to appear thinner than you are, the hand of someone who reaches out to touch you would bump into you while it was seemingly still in midair.

To discern that you are disguised, a creature can use its action to inspect your appearance and must succeed on an Intelligence (Investigation) check against your spell save DC.

\addcontentsline{toc}{section}{Disintegrate}


\DndSpellHeader%
{Disintegrate}
{6th level Transmutation}
{1 action}
{60 feet}
{V, S, a lodestone and a pinch of dust}
{Instantaneous}
A thin green ray springs from your pointing finger to a target that you can see within range. The target can be a creature, an object, or a creation of magical force, such as the wall created by \textit{wall of force}.

A creature targeted by this spell must make a Dexterity saving throw. On a failed save, the target takes 10d6 + 40 force damage. The target is disintegrated if this damage leaves it with 0 hit points.

A disintegrated creature and everything it is wearing and carrying, except magic items, are reduced to a pile of fine gray dust. The creature can be restored to life only by means of a \textit{true resurrection} or a \textit{wish} spell.

This spell automatically disintegrates a Large or smaller nonmagical object or a creation of magical force. If the target is a Huge or larger object or creation of force, this spell disintegrates a 10-foot-cube portion of it. A magic item is unaffected by this spell.

\addcontentsline{toc}{section}{Dispel Magic}


\DndSpellHeader%
{Dispel Magic}
{3rd level Abjuration}
{1 action}
{120 feet}
{V, S}
{Instantaneous}
Choose one creature, object, or magical effect within range. Any spell of 3rd level or lower on the target ends. For each spell of 4th level or higher on the target, make an ability check using your spellcasting ability. The DC equals 10 + the spell's level. On a successful check, the spell ends.

\addcontentsline{toc}{section}{Dominate Beast}


\DndSpellHeader%
{Dominate Beast}
{4th level Enchantment}
{1 action}
{60 feet}
{V, S}
{Concentration, up to 1 minute}
You attempt to beguile a beast that you can see within range. It must succeed on a Wisdom saving throw or be charmed by you for the duration. If you or creatures that are friendly to you are fighting it, it has advantage on the saving throw.

While the beast is charmed, you have a telepathic link with it as long as the two of you are on the same plane of existence. You can use this telepathic link to issue commands to the creature while you are conscious (no action required), which it does its best to obey. You can specify a simple and general course of action, such as "Attack that creature," "Run over there," or "Fetch that object." If the creature completes the order and doesn't receive further direction from you, it defends and preserves itself to the best of its ability.

You can use your action to take total and precise control of the target. Until the end of your next turn, the creature takes only the actions you choose, and doesn't do anything that you don't allow it to do. During this time, you can also cause the creature to use a reaction, but this requires you to use your own reaction as well.

Each time the target takes damage, it makes a new Wisdom saving throw against the spell. If the saving throw succeeds, the spell ends.

\addcontentsline{toc}{section}{Enlarge/Reduce}


\DndSpellHeader%
{Enlarge/Reduce}
{2nd level Transmutation}
{1 action}
{30 feet}
{V, S, a pinch of powdered iron}
{Concentration, up to 1 minute}
You cause a creature or an object you can see within range to grow larger or smaller for the duration. Choose either a creature or an object that is neither worn nor carried. If the target is unwilling, it can make a Constitution saving throw. On a success, the spell has no effect.

If the target is a creature, everything it is wearing and carrying changes size with it. Any item dropped by an affected creature returns to normal size at once.

\subparagraph{Enlarge}
The target's size doubles in all dimensions, and its weight is multiplied by eight. This growth increases its size by one category—from Medium to Large, for example. If there isn't enough room for the target to double its size, the creature or object attains the maximum possible size in the space available. Until the spell ends, the target also has advantage on Strength checks and Strength saving throws. The target's weapons also grow to match its new size. While these weapons are enlarged, the target's attacks with them deal 1d4 extra damage.

\subparagraph{Reduce}
The target's size is halved in all dimensions, and its weight is reduced to one-eighth of normal. This reduction decreases its size by one category—from Medium to Small, for example. Until the spell ends, the target also has disadvantage on Strength checks and Strength saving throws. The target's weapons also shrink to match its new size. While these weapons are reduced, the target's attacks with them deal 1d4 less damage (this can't reduce the damage below 1).

\addcontentsline{toc}{section}{Entangle}


\DndSpellHeader%
{Entangle}
{1st level Conjuration}
{1 action}
{90 feet}
{V, S}
{Concentration, up to 1 minute}
Grasping weeds and vines sprout from the ground in a 20-foot square starting from a point within range. For the duration, these plants turn the ground in the area into difficult terrain.

A creature in the area when you cast the spell must succeed on a Strength saving throw or be restrained by the entangling plants until the spell ends. A creature restrained by the plants can use its action to make a Strength check against your spell save DC. On a success, it frees itself.

When the spell ends, the conjured plants wilt away.

\addcontentsline{toc}{section}{Evard's Black Tentacles}


\DndSpellHeader%
{Evard's Black Tentacles}
{4th level Conjuration}
{1 action}
{90 feet}
{V, S, a piece of tentacle from a giant octopus or a giant squid}
{Concentration, up to 1 minute}
Squirming, ebony tentacles fill a 20-foot square on ground that you can see within range. For the duration, these tentacles turn the ground in the area into difficult terrain.

When a creature enters the affected area for the first time on a turn or starts its turn there, the creature must succeed on a Dexterity saving throw or take 3d6 bludgeoning damage and be restrained by the tentacles until the spell ends. A creature that starts its turn in the area and is already restrained by the tentacles takes 3d6 bludgeoning damage.

A creature restrained by the tentacles can use its action to make a Strength or Dexterity check (its choice) against your spell save DC. On a success, it frees itself.

\addcontentsline{toc}{section}{Expeditious Retreat}


\DndSpellHeader%
{Expeditious Retreat}
{1st level Transmutation}
{1 bonus}
{Self}
{V, S}
{Concentration, up to 10 minutes}
This spell allows you to move at an incredible pace. When you cast this spell, and then as a bonus action on each of your turns until the spell ends, you can take the Dash action.

\addcontentsline{toc}{section}{False Life}


\DndSpellHeader%
{False Life}
{1st level Necromancy}
{1 action}
{Self}
{V, S, a small amount of alcohol or distilled spirits}
{1 hour}
Bolstering yourself with a necromantic facsimile of life, you gain 1d4 + 4 temporary hit points for the duration.

\addcontentsline{toc}{section}{Fear}


\DndSpellHeader%
{Fear}
{3rd level Illusion}
{1 action}
{Self (30 feet cone)}
{V, S, a white feather or the heart of a hen}
{Concentration, up to 1 minute}
You project a phantasmal image of a creature's worst fears. Each creature in a 30-foot cone must succeed on a Wisdom saving throw or drop whatever it is holding and become frightened for the duration.

While frightened by this spell, a creature must take the Dash action and move away from you by the safest available route on each of its turns, unless there is nowhere to move. If the creature ends its turn in a location where it doesn't have line of sight to you, the creature can make a Wisdom saving throw. On a successful save, the spell ends for that creature.

\addcontentsline{toc}{section}{Feather Fall}


\DndSpellHeader%
{Feather Fall}
{1st level Transmutation}
{1 reaction}
{60 feet}
{V, a small feather or a piece of down}
{1 minute}
Choose up to five falling creatures within range. A falling creature's rate of descent slows to 60 feet per round until the spell ends. If the creature lands before the spell ends, it takes no falling damage and can land on its feet, and the spell ends for that creature.

\addcontentsline{toc}{section}{Fire Bolt}


\DndSpellHeader%
{Fire Bolt}
{Evocation cantrip}
{1 action}
{120 feet}
{V, S}
{Instantaneous}
You hurl a mote of fire at a creature or object within range. Make a ranged spell attack against the target. On a hit, the target takes 1d10 fire damage. A flammable object hit by this spell ignites if it isn't being worn or carried.

This spell's damage increases by 1d10 when you reach 5th level (2d10), 11th level (3d10), and 17th level (4d10).

\addcontentsline{toc}{section}{Fireball}


\DndSpellHeader%
{Fireball}
{3rd level Evocation}
{1 action}
{150 feet}
{V, S, a tiny ball of bat guano and sulfur}
{Instantaneous}
A bright streak flashes from your pointing finger to a point you choose within range and then blossoms with a low roar into an explosion of flame. Each creature in a 20-foot-radius sphere centered on that point must make a Dexterity saving throw. A target takes 8d6 fire damage on a failed save, or half as much damage on a successful one.

The fire spreads around corners. It ignites flammable objects in the area that aren't being worn or carried.

\addcontentsline{toc}{section}{Flaming Sphere}


\DndSpellHeader%
{Flaming Sphere}
{2nd level Conjuration}
{1 action}
{60 feet}
{V, S, a bit of tallow, a pinch of brimstone, and a dusting of powdered iron}
{Concentration, up to 1 minute}
A 5-foot-diameter sphere of fire appears in an unoccupied space of your choice within range and lasts for the duration. Any creature that ends its turn within 5 feet of the sphere must make a Dexterity saving throw. The creature takes 2d6 fire damage on a failed save, or half as much damage on a successful one.

As a bonus action, you can move the sphere up to 30 feet. If you ram the sphere into a creature, that creature must make the saving throw against the sphere's damage, and the sphere stops moving this turn.

When you move the sphere, you can direct it over barriers up to 5 feet tall and jump it across pits up to 10 feet wide. The sphere ignites flammable objects not being worn or carried, and it sheds bright light in a 20-foot radius and dim light for an additional 20 feet.

\addcontentsline{toc}{section}{Flesh to Stone}


\DndSpellHeader%
{Flesh to Stone}
{6th level Transmutation}
{1 action}
{60 feet}
{V, S, a pinch of lime, water, and earth}
{Concentration, up to 1 minute}
You attempt to turn one creature that you can see within range into stone. If the target's body is made of flesh, the creature must make a Constitution saving throw. On a failed save, it is restrained as its flesh begins to harden. On a successful save, the creature isn't affected.

A creature restrained by this spell must make another Constitution saving throw at the end of each of its turns. If it successfully saves against this spell three times, the spell ends. If it fails its saves three times, it is turned to stone and subjected to the petrified condition for the duration. The successes and failures don't need to be consecutive; keep track of both until the target collects three of a kind.

If the creature is physically broken while petrified, it suffers from similar deformities if it reverts to its original state.

If you maintain your concentration on this spell for the entire possible duration, the creature is turned to stone until the effect is removed.

\addcontentsline{toc}{section}{Fly}


\DndSpellHeader%
{Fly}
{3rd level Transmutation}
{1 action}
{Touch}
{V, S, a wing feather from any bird}
{Concentration, up to 10 minutes}
You touch a willing creature. The target gains a flying speed of 60 feet for the duration. When the spell ends, the target falls if it is still aloft, unless it can stop the fall.

\addcontentsline{toc}{section}{Fog Cloud}


\DndSpellHeader%
{Fog Cloud}
{1st level Conjuration}
{1 action}
{120 feet}
{V, S}
{Concentration, up to 1 hour}
You create a 20-foot-radius sphere of fog centered on a point within range. The sphere spreads around corners, and its area is Vision and Light. It lasts for the duration or until a wind of moderate or greater speed (at least 10 miles per hour) disperses it.

\addcontentsline{toc}{section}{Friends}


\DndSpellHeader%
{Friends}
{Enchantment cantrip}
{1 action}
{Self}
{S, a small amount of makeup applied to the face as this spell is cast}
{Concentration, up to 1 minute}
For the duration, you have advantage on all Charisma checks directed at one creature of your choice that isn't hostile toward you. When the spell ends, the creature realizes that you used magic to influence its mood and becomes hostile toward you. A creature prone to violence might attack you. Another creature might seek retribution in other ways (at the DM's discretion), depending on the nature of your interaction with it.

\addcontentsline{toc}{section}{Gaseous Form}


\DndSpellHeader%
{Gaseous Form}
{3rd level Transmutation}
{1 action}
{Touch}
{V, S, a bit of gauze and a wisp of smoke}
{Concentration, up to 1 hour}
You transform a willing creature you touch, along with everything it's wearing and carrying, into a misty cloud for the duration. The spell ends if the creature drops to 0 hit points. An incorporeal creature isn't affected.

While in this form, the target's only method of movement is a flying speed of 10 feet. The target can enter and occupy the space of another creature. The target has resistance to nonmagical damage, and it has advantage on Strength, Dexterity, and Constitution saving throws. The target can pass through small holes, narrow openings, and even mere cracks, though it treats liquids as though they were solid surfaces. The target can't fall and remains hovering in the air even when stunned or otherwise incapacitated.

While in the form of a misty cloud, the target can't talk or manipulate objects, and any objects it was carrying or holding can't be dropped, used, or otherwise interacted with. The target can't attack or cast spells.

\addcontentsline{toc}{section}{Glyph of Warding}


\DndSpellHeader%
{Glyph of Warding}
{3rd level Abjuration}
{1 hour}
{Touch}
{V, S, incense and powdered diamond worth at least 200 gp, which the spell consumes}
{Until dispelled or triggered}
When you cast this spell, you inscribe a glyph that later unleashes a magical effect. You inscribe it either on a surface (such as a table or a section of floor or wall) or within an object that can be closed (such as a book, a scroll, or a treasure chest) to conceal the glyph. The glyph can cover an area no larger than 10 feet in diameter. If the surface or object is moved more than 10 feet from where you cast this spell, the glyph is broken, and the spell ends without being triggered.

The glyph is nearly invisible and requires a successful Intelligence (Investigation) check against your spell save DC to be found.

You decide what triggers the glyph when you cast the spell. For glyphs inscribed on a surface, the most typical triggers include touching or standing on the glyph, removing another object covering the glyph, approaching within a certain distance of the glyph, or manipulating the object on which the glyph is inscribed. For glyphs inscribed within an object, the most common triggers include opening that object, approaching within a certain distance of the object, or seeing or reading the glyph. Once a glyph is triggered, this spell ends.

You can further refine the trigger so the spell activates only under certain circumstances or according to physical characteristics (such as height or weight), creature kind (for example, the ward could be set to affect aberrations or drow), or alignment. You can also set conditions for creatures that don't trigger the glyph, such as those who say a certain password.

When you inscribe the glyph, choose explosive runes or a spell glyph.

\subparagraph{Explosive Runes}
When triggered, the glyph erupts with magical energy in a 20-foot-radius sphere centered on the glyph. The sphere spreads around corners. Each creature in the area must make a Dexterity saving throw. A creature takes 5d8 acid, cold, fire, lightning, or thunder damage on a failed saving throw (your choice when you create the glyph), or half as much damage on a successful one.

\subparagraph{Spell Glyph}
You can store a prepared spell of 3rd level or lower in the glyph by casting it as part of creating the glyph. The spell must target a single creature or an area. The spell being stored has no immediate effect when cast in this way. When the glyph is triggered, the stored spell is cast. If the spell has a target, it targets the creature that triggered the glyph. If the spell affects an area, the area is centered on that creature. If the spell summons hostile creatures or creates harmful objects or traps, they appear as close as possible to the intruder and attack it. If the spell requires concentration, it lasts until the end of its full duration.

\addcontentsline{toc}{section}{Greater Invisibility}


\DndSpellHeader%
{Greater Invisibility}
{4th level Illusion}
{1 action}
{Touch}
{V, S}
{Concentration, up to 1 minute}
You or a creature you touch becomes invisible until the spell ends. Anything the target is wearing or carrying is invisible as long as it is on the target's person.

\addcontentsline{toc}{section}{Haste}


\DndSpellHeader%
{Haste}
{3rd level Transmutation}
{1 action}
{30 feet}
{V, S, a shaving of licorice root}
{Concentration, up to 1 minute}
Choose a willing creature that you can see within range. Until the spell ends, the target's speed is doubled, it gains a +2 bonus to AC, it has advantage on Dexterity saving throws, and it gains an additional action on each of its turns. That action can be used only to take the Attack (one weapon attack only), Dash, Disengage, Hide, or Use an Object action.

When the spell ends, the target can't move or take actions until after its next turn, as a wave of lethargy sweeps over it.

\addcontentsline{toc}{section}{Hold Monster}


\DndSpellHeader%
{Hold Monster}
{5th level Enchantment}
{1 action}
{90 feet}
{V, S, a small, straight piece of iron}
{Concentration, up to 1 minute}
Choose a creature that you can see within range. The target must succeed on a Wisdom saving throw or be paralyzed for the duration. This spell has no effect on undead. At the end of each of its turns, the target can make another Wisdom saving throw. On a success, the spell ends on the target.

\addcontentsline{toc}{section}{Hold Person}


\DndSpellHeader%
{Hold Person}
{2nd level Enchantment}
{1 action}
{60 feet}
{V, S, a small, straight piece of iron}
{Concentration, up to 1 minute}
Choose a humanoid that you can see within range. The target must succeed on a Wisdom saving throw or be paralyzed for the duration. At the end of each of its turns, the target can make another Wisdom saving throw. On a success, the spell ends on the target.

\addcontentsline{toc}{section}{Ice Storm}


\DndSpellHeader%
{Ice Storm}
{4th level Evocation}
{1 action}
{300 feet}
{V, S, a pinch of dust and a few drops of water}
{Instantaneous}
A hail of rock-hard ice pounds to the ground in a 20-foot-radius, 40-foot-high cylinder centered on a point within range. Each creature in the cylinder must make a Dexterity saving throw. A creature takes 2d8 bludgeoning damage and 4d6 cold damage on a failed save, or half as much damage on a successful one.

Hailstones turn the storm's area of effect into difficult terrain until the end of your next turn.

\addcontentsline{toc}{section}{Identify}


\DndSpellHeader%
{Identify}
{1st level Divination}
{1 minute}
{Touch}
{V, S, a pearl worth at least 100 gp and an owl feather}
{Instantaneous}
You choose one object that you must touch throughout the casting of the spell. If it is a magic item or some other magic-imbued object, you learn its properties and how to use them, whether it requires attunement to use, and how many charges it has, if any. You learn whether any spells are affecting the item and what they are. If the item was created by a spell, you learn which spell created it.

If you instead touch a creature throughout the casting, you learn what spells, if any, are currently affecting it.

\addcontentsline{toc}{section}{Invisibility}


\DndSpellHeader%
{Invisibility}
{2nd level Illusion}
{1 action}
{Touch}
{V, S, an eyelash encased in gum arabic}
{Concentration, up to 1 hour}
A creature you touch becomes invisible until the spell ends. Anything the target is wearing or carrying is invisible as long as it is on the target's person. The spell ends for a target that attacks or casts a spell.

\addcontentsline{toc}{section}{Jump}


\DndSpellHeader%
{Jump}
{1st level Transmutation}
{1 action}
{Touch}
{V, S, a grasshopper's hind leg}
{1 minute}
You touch a creature. The creature's \textit{jump distance} is tripled until the spell ends.

\addcontentsline{toc}{section}{Knock}


\DndSpellHeader%
{Knock}
{2nd level Transmutation}
{1 action}
{60 feet}
{V}
{Instantaneous}
Choose an object that you can see within range. The object can be a door, a box, a \textit{chest}, a set of \textit{manacles}, a padlock, or another object that contains a mundane or magical means that prevents access.

A target that is held shut by a mundane lock or that is stuck or barred becomes unlocked, unstuck, or unbarred. If the object has multiple locks, only one of them is unlocked.

If you choose a target that is held shut with \textit{arcane lock}, that spell is suppressed for 10 minutes, during which time the target can be opened and shut normally.

When you cast the spell, a loud knock, audible from as far away as 300 feet, emanates from the target object.

\addcontentsline{toc}{section}{Legend Lore}


\DndSpellHeader%
{Legend Lore}
{5th level Divination}
{10 minute}
{Self}
{V, S, incense worth at least 250 gp, which the spell consumes, and four ivory strips worth at least 50 gp each}
{Instantaneous}
Name or describe a person, place, or object. The spell brings to your mind a brief summary of the significant lore about the thing you named. The lore might consist of current tales, forgotten stories, or even secret lore that has never been widely known. If the thing you named isn't of legendary importance, you gain no information. The more information you already have about the thing, the more precise and detailed the information you receive is.

The information you learn is accurate but might be couched in figurative language. For example, if you have a mysterious magic axe on hand, the spell might yield this information: "Woe to the evildoer whose hand touches the axe, for even the haft slices the hand of the evil ones. Only a true Child of Stone, lover and beloved of Moradin, may awaken the true powers of the axe, and only with the sacred word Rudnogg on the lips."

\addcontentsline{toc}{section}{Levitate}


\DndSpellHeader%
{Levitate}
{2nd level Transmutation}
{1 action}
{60 feet}
{V, S, either a small leather loop or a piece of golden wire bent into a cup shape with a long shank on one end}
{Concentration, up to 10 minutes}
One creature or loose object of your choice that you can see within range rises vertically, up to 20 feet, and remains suspended there for the duration. The spell can levitate a target that weighs up to 500 pounds. An unwilling creature that succeeds on a Constitution saving throw is unaffected.

The target can move only by pushing or pulling against a fixed object or surface within reach (such as a wall or a ceiling), which allows it to move as if it were climbing. You can change the target's altitude by up to 20 feet in either direction on your turn. If you are the target, you can move up or down as part of your move. Otherwise, you can use your action to move the target, which must remain within the spell's range.

When the spell ends, the target floats gently to the ground if it is still aloft.

\addcontentsline{toc}{section}{Light}


\DndSpellHeader%
{Light}
{Evocation cantrip}
{1 action}
{Touch}
{V, a firefly or phosphorescent moss}
{1 hour}
You touch one object that is no larger than 10 feet in any dimension. Until the spell ends, the object sheds bright light in a 20-foot radius and dim light for an additional 20 feet. The light can be colored as you like. Completely covering the object with something opaque blocks the light. The spell ends if you cast it again or dismiss it as an action.

If you target an object held or worn by a hostile creature, that creature must succeed on a Dexterity saving throw to avoid the spell.

\addcontentsline{toc}{section}{Lightning Bolt}


\DndSpellHeader%
{Lightning Bolt}
{3rd level Evocation}
{1 action}
{Self (100 feet line)}
{V, S, a bit of fur and a rod of amber, crystal, or glass}
{Instantaneous}
A stroke of lightning forming a line 100 feet long and 5 feet wide blasts out from you in a direction you choose. Each creature in the line must make a Dexterity saving throw. A creature takes 8d6 lightning damage on a failed save, or half as much damage on a successful one.

The lightning ignites flammable objects in the area that aren't being worn or carried.

\addcontentsline{toc}{section}{Locate Object}


\DndSpellHeader%
{Locate Object}
{2nd level Divination}
{1 action}
{Self}
{V, S, a forked twig}
{Concentration, up to 10 minutes}
Describe or name an object that is familiar to you. You sense the direction to the object's location, as long as that object is within 1,000 feet of you. If the object is in motion, you know the direction of its movement.

The spell can locate a specific object known to you, as long as you have seen it up close—within 30 feet—at least once. Alternatively, the spell can locate the nearest object of a particular kind, such as a certain kind of apparel, jewelry, furniture, tool, or weapon.

This spell can't locate an object if any thickness of lead, even a thin sheet, blocks a direct path between you and the object.

\addcontentsline{toc}{section}{Mage Armor}


\DndSpellHeader%
{Mage Armor}
{1st level Abjuration}
{1 action}
{Touch}
{V, S, a piece of cured leather}
{8 hours}
You touch a willing creature who isn't wearing armor, and a protective magical force surrounds it until the spell ends. The target's base AC becomes 13 + its Dexterity modifier. The spell ends if the target dons armor or if you dismiss the spell as an action.

\addcontentsline{toc}{section}{Mage Hand}


\DndSpellHeader%
{Mage Hand}
{Conjuration cantrip}
{1 action}
{30 feet}
{V, S}
{1 minute}
A spectral, floating hand appears at a point you choose within range. The hand lasts for the duration or until you dismiss it as an action. The hand vanishes if it is ever more than 30 feet away from you or if you cast this spell again.

You can use your action to control the hand. You can use the hand to manipulate an object, open an unlocked door or container, stow or retrieve an item from an open container, or pour the contents out of a vial. You can move the hand up to 30 feet each time you use it.

The hand can't attack, activate magic items, or carry more than 10 pounds.

\addcontentsline{toc}{section}{Magic Missile}


\DndSpellHeader%
{Magic Missile}
{1st level Evocation}
{1 action}
{120 feet}
{V, S}
{Instantaneous}
You create three glowing darts of magical force. Each dart hits a creature of your choice that you can see within range. A dart deals 1d4 + 1 force damage to its target. The darts all strike simultaneously, and you can direct them to hit one creature or several.

\addcontentsline{toc}{section}{Magic Mouth}


\DndSpellHeader%
{Magic Mouth}
{2nd level Illusion}
{1 minute}
{30 feet}
{V, S, a small bit of honeycomb and jade dust worth at least 10 gp, which the spell consumes}
{Until dispelled}
You implant a message within an object in range, a message that is uttered when a trigger condition is met. Choose an object that you can see and that isn't being worn or carried by another creature. Then speak the message, which must be 25 words or less, though it can be delivered over as long as 10 minutes. Finally, determine the circumstance that will trigger the spell to deliver your message.

When that circumstance occurs, a magical mouth appears on the object and recites the message in your voice and at the same volume you spoke. If the object you chose has a mouth or something that looks like a mouth (for example, the mouth of a statue), the magical mouth appears there so that the words appear to come from the object's mouth. When you cast this spell, you can have the spell end after it delivers its message, or it can remain and repeat its message whenever the trigger occurs.

The triggering circumstance can be as general or as detailed as you like, though it must be based on visual or audible conditions that occur within 30 feet of the object. For example, you could instruct the mouth to speak when any creature moves within 30 feet of the object or when a silver bell rings within 30 feet of it.

\addcontentsline{toc}{section}{Major Image}


\DndSpellHeader%
{Major Image}
{3rd level Illusion}
{1 action}
{120 feet}
{V, S, a bit of fleece}
{Concentration, up to 10 minutes}
You create the image of an object, a creature, or some other visible phenomenon that is no larger than a 20-foot cube. The image appears at a spot that you can see within range and lasts for the duration. It seems completely real, including sounds, smells, and temperature appropriate to the thing depicted. You can't create sufficient heat or cold to cause damage, a sound loud enough to deal thunder damage or deafen a creature, or a smell that might sicken a creature (like a troglodyte's stench).

As long as you are within range of the illusion, you can use your action to cause the image to move to any other spot within range. As the image changes location, you can alter its appearance so that its movements appear natural for the image. For example, if you create an image of a creature and move it, you can alter the image so that it appears to be walking. Similarly, you can cause the illusion to make different sounds at different times, even making it carry on a conversation, for example.

Physical interaction with the image reveals it to be an illusion, because things can pass through it. A creature that uses its action to examine the image can determine that it is an illusion with a successful Intelligence (Investigation) check against your spell save DC. If a creature discerns the illusion for what it is, the creature can see through the image, and its other sensory qualities become faint to the creature.

\addcontentsline{toc}{section}{Melf's Acid Arrow}


\DndSpellHeader%
{Melf's Acid Arrow}
{2nd level Evocation}
{1 action}
{90 feet}
{V, S, powdered rhubarb leaf and an adder's stomach}
{Instantaneous}
A shimmering green arrow streaks toward a target within range and bursts in a spray of acid. Make a ranged spell attack against the target. On a hit, the target takes 4d4 acid damage immediately and 2d4 acid damage at the end of its next turn. On a miss, the arrow splashes the target with acid for half as much of the initial damage and no damage at the end of its next turn.

\addcontentsline{toc}{section}{Mending}


\DndSpellHeader%
{Mending}
{Transmutation cantrip}
{1 minute}
{Touch}
{V, S, two lodestones}
{Instantaneous}
This spell repairs a single break or tear in an object you touch, such as broken chain link, two halves of a broken key, a torn cloak, or a leaking wineskin. As long as the break or tear is no larger than 1 foot in any dimension, you mend it, leaving no trace of the former damage.

This spell can physically repair a magic item or construct, but the spell can't restore magic to such an object.

\addcontentsline{toc}{section}{Message}


\DndSpellHeader%
{Message}
{Transmutation cantrip}
{1 action}
{120 feet}
{V, S, a short piece of copper wire}
{1 round}
You point your finger toward a creature within range and whisper a message. The target (and only the target) hears the message and can reply in a whisper that only you can hear.

You can cast this spell through solid objects if you are familiar with the target and know it is beyond the barrier. Magical silence, 1 foot of stone, 1 inch of common metal, a thin sheet of lead, or 3 feet of wood blocks the spell. The spell doesn't have to follow a straight line and can travel freely around corners or through openings.

\addcontentsline{toc}{section}{Minor Illusion}


\DndSpellHeader%
{Minor Illusion}
{Illusion cantrip}
{1 action}
{30 feet}
{S, a bit of fleece}
{1 minute}
You create a sound or an image of an object within range that lasts for the duration. The illusion also ends if you dismiss it as an action or cast this spell again.

If you create a sound, its volume can range from a whisper to a scream. It can be your voice, someone else's voice, a lion's roar, a beating of drums, or any other sound you choose. The sound continues unabated throughout the duration, or you can make discrete sounds at different times before the spell ends.

If you create an image of an object—such as a chair, muddy footprints, or a small chest—it must be no larger than a 5-foot cube. The image can't create sound, light, smell, or any other sensory effect. Physical interaction with the image reveals it to be an illusion, because things can pass through it.

If a creature uses its action to examine the sound or image, the creature can determine that it is an illusion with a successful Intelligence (Investigation) check against your spell save DC. If a creature discerns the illusion for what it is, the illusion becomes faint to the creature.

\addcontentsline{toc}{section}{Mirror Image}


\DndSpellHeader%
{Mirror Image}
{2nd level Illusion}
{1 action}
{Self}
{V, S}
{1 minute}
Three illusory duplicates of yourself appear in your space. Until the spell ends, the duplicates move with you and mimic your actions, shifting position so it's impossible to track which image is real. You can use your action to dismiss the illusory duplicates.

Each time a creature targets you with an attack during the spell's duration, roll a d20 to determine whether the attack instead targets one of your duplicates.

If you have three duplicates, you must roll a 6 or higher to change the attack's target to a duplicate. With two duplicates, you must roll an 8 or higher. With one duplicate, you must roll an 11 or higher.

A duplicate's AC equals 10 + your Dexterity modifier. If an attack hits a duplicate, the duplicate is destroyed. A duplicate can be destroyed only by an attack that hits it. It ignores all other damage and effects. The spell ends when all three duplicates are destroyed.

A creature is unaffected by this spell if it can't see, if it relies on senses other than sight, such as blindsight, or if it can perceive illusions as false, as with truesight.

\addcontentsline{toc}{section}{Misty Step}


\DndSpellHeader%
{Misty Step}
{2nd level Conjuration}
{1 bonus}
{Self}
{V}
{Instantaneous}
Briefly surrounded by silvery mist, you teleport up to 30 feet to an unoccupied space that you can see.

\addcontentsline{toc}{section}{Pass without Trace}


\DndSpellHeader%
{Pass without Trace}
{2nd level Abjuration}
{1 action}
{Self}
{V, S, ashes from a burned leaf of mistletoe and a sprig of spruce}
{Concentration, up to 1 hour}
A veil of shadows and silence radiates from you, masking you and your companions from detection. For the duration, each creature you choose within 30 feet of you (including you) has a +10 bonus to Dexterity (Stealth) checks and can't be tracked except by magical means. A creature that receives this bonus leaves behind no tracks or other traces of its passage.

\addcontentsline{toc}{section}{Passwall}


\DndSpellHeader%
{Passwall}
{5th level Transmutation}
{1 action}
{30 feet}
{V, S, a pinch of sesame seeds}
{1 hour}
A passage appears at a point of your choice that you can see on a wooden, plaster, or stone surface (such as a wall, a ceiling, or a floor) within range, and lasts for the duration. You choose the opening's dimensions: up to 5 feet wide, 8 feet tall, and 20 feet deep. The passage creates no instability in a structure surrounding it.

When the opening disappears, any creatures or objects still in the passage created by the spell are safely ejected to an unoccupied space nearest to the surface on which you cast the spell.

\addcontentsline{toc}{section}{Phantasmal Force}


\DndSpellHeader%
{Phantasmal Force}
{2nd level Illusion}
{1 action}
{60 feet}
{V, S, a bit of fleece}
{Concentration, up to 1 minute}
You craft an illusion that takes root in the mind of a creature that you can see within range. The target must make an Intelligence saving throw. On a failed save, you create a phantasmal object, creature, or other visible phenomenon of your choice that is no larger than a 10-foot cube and that is perceivable only to the target for the duration. This spell has no effect on undead or constructs.

The phantasm includes sound, temperature, and other stimuli, also evident only to the creature.

The target can use its action to examine the phantasm with an Intelligence (Investigation) check against your spell save DC. If the check succeeds, the target realizes that the phantasm is an illusion, and the spell ends.

While a target is affected by the spell, the target treats the phantasm as if it were real. The target rationalizes any illogical outcomes from interacting with the phantasm. For example, a target attempting to walk across a phantasmal bridge that spans a chasm falls once it steps onto the bridge. If the target survives the fall, it still believes that the bridge exists and comes up with some other explanation for its fall—it was pushed, it slipped, or a strong wind might have knocked it off.

An affected target is so convinced of the phantasm's reality that it can even take damage from the illusion. A phantasm created to appear as a creature can attack the target. Similarly, a phantasm created to appear as fire, a pool of acid, or lava can burn the target. Each round on your turn, the phantasm can deal 1d6 psychic damage to the target if it is in the phantasm's area or within 5 feet of the phantasm, provided that the illusion is of a creature or hazard that could logically deal damage, such as by attacking. The target perceives the damage as a type appropriate to the illusion.

\addcontentsline{toc}{section}{Phantasmal Killer}


\DndSpellHeader%
{Phantasmal Killer}
{4th level Illusion}
{1 action}
{120 feet}
{V, S}
{Concentration, up to 1 minute}
You tap into the nightmares of a creature you can see within range and create an illusory manifestation of its deepest fears, visible only to that creature. The target must make a Wisdom saving throw. On a failed save, the target becomes frightened for the duration. At the end of each of the target's turns before the spell ends, the target must succeed on a Wisdom saving throw or take 4d10 psychic damage. On a successful save, the spell ends.

\addcontentsline{toc}{section}{Phantom Steed}


\DndSpellHeader%
{Phantom Steed}
{3rd level Illusion}
{1 minute}
{30 feet}
{V, S}
{1 hour}
A Large quasi-real, horselike creature appears on the ground in an unoccupied space of your choice within range. You decide the creature's appearance, but it is equipped with a saddle, bit, and bridle. Any of the equipment created by the spell vanishes in a puff of smoke if it is carried more than 10 feet away from the steed.

For the duration, you or a creature you choose can ride the steed. The creature uses the statistics for a \textbf{riding horse}, except it has a speed of 100 feet and can travel 10 miles in an hour, or 13 miles at a fast pace. When the spell ends, the steed gradually fades, giving the rider 1 minute to dismount. The spell ends if you use an action to dismiss it or if the steed takes any damage.

\addcontentsline{toc}{section}{Plane Shift}


\DndSpellHeader%
{Plane Shift}
{7th level Conjuration}
{1 action}
{Touch}
{V, S, a forked, metal rod worth at least 250 gp, attuned to a particular plane of existence}
{Instantaneous}
You and up to eight willing creatures who link hands in a circle are transported to a different plane of existence. You can specify a target destination in general terms, such as the City of Brass on the Elemental Plane of Fire or the palace of Dispater on the second level of the Nine Hells, and you appear in or near that destination. If you are trying to reach the City of Brass, for example, you might arrive in its Street of Steel, before its Gate of Ashes, or looking at the city from across the Sea of Fire, at the DM's discretion.

Alternatively, if you know the sigil sequence of a teleportation circle on another plane of existence, this spell can take you to that circle. If the teleportation circle is too small to hold all the creatures you transported, they appear in the closest unoccupied spaces next to the circle.

You can use this spell to banish an unwilling creature to another plane. Choose a creature within your reach and make a melee spell attack against it. On a hit, the creature must make a Charisma saving throw. If the creature fails this save, it is transported to a random location on the plane of existence you specify. A creature so transported must find its own way back to your current plane of existence.

\addcontentsline{toc}{section}{Poison Spray}


\DndSpellHeader%
{Poison Spray}
{Conjuration cantrip}
{1 action}
{10 feet}
{V, S}
{Instantaneous}
You extend your hand toward a creature you can see within range and project a puff of noxious gas from your palm. The creature must succeed on a Constitution saving throw or take 1d12 poison damage.

This spell's damage increases by 1d12 when you reach 5th level (2d12), 11th level (3d12), and 17th level (4d12).

\addcontentsline{toc}{section}{Polymorph}


\DndSpellHeader%
{Polymorph}
{4th level Transmutation}
{1 action}
{60 feet}
{V, S, a caterpillar cocoon}
{Concentration, up to 1 hour}
This spell transforms a creature that you can see within range into a new form. An unwilling creature must make a Wisdom saving throw to avoid the effect. The spell has no effect on a shapechanger or a creature with 0 hit points.

The transformation lasts for the duration, or until the target drops to 0 hit points or dies. The new form can be any beast whose challenge rating is equal to or less than the target's (or the target's level, if it doesn't have a challenge rating). The target's game statistics, including mental ability scores, are replaced by the statistics of the chosen beast. It retains its alignment and personality.

The target assumes the hit points of its new form. When it reverts to its normal form, the creature returns to the number of hit points it had before it transformed. If it reverts as a result of dropping to 0 hit points, any excess damage carries over to its normal form. As long as the excess damage doesn't reduce the creature's normal form to 0 hit points, it isn't knocked unconscious.

The creature is limited in the actions it can perform by the nature of its new form, and it can't speak, cast spells, or take any other action that requires hands or speech.

The target's gear melds into the new form. The creature can't activate, use, wield, or otherwise benefit from any of its equipment.

\addcontentsline{toc}{section}{Prestidigitation}


\DndSpellHeader%
{Prestidigitation}
{Transmutation cantrip}
{1 action}
{10 feet}
{V, S}
{Up to 1 hour}
This spell is a minor magical trick that novice spellcasters use for practice. You create one of the following magical effects within range:


\begin{itemize}
\item You create an instantaneous, harmless sensory effect, such as a shower of sparks, a puff of wind, faint musical notes, or an odd odor.
\item You instantaneously light or snuff out a candle, a torch, or a small campfire.
\item You instantaneously clean or soil an object no larger than 1 cubic foot.
\item You chill, warm, or flavor up to 1 cubic foot of nonliving material for 1 hour.
\item You make a color, a small mark, or a symbol appear on an object or a surface for 1 hour.
\item You create a nonmagical trinket or an illusory image that can fit in your hand and that lasts until the end of your next turn.
\end{itemize}

If you cast this spell multiple times, you can have up to three of its non-instantaneous effects active at a time, and you can dismiss such an effect as an action.

\addcontentsline{toc}{section}{Prismatic Spray}


\DndSpellHeader%
{Prismatic Spray}
{7th level Evocation}
{1 action}
{Self (60 feet cone)}
{V, S}
{Instantaneous}
Eight multicolored rays of light flash from your hand. Each ray is a different color and has a different power and purpose. Each creature in a 60-foot cone must make a Dexterity saving throw. For each target, roll a d8 to determine which color ray affects it.

\subparagraph{1-Red}
The target takes 10d6 fire damage on a failed save, or half as much damage on a successful one.

\subparagraph{2-Orange}
The target takes 10d6 acid damage on a failed save, or half as much damage on a successful one.

\subparagraph{3-Yellow}
The target takes 10d6 lightning damage on a failed save, or half as much damage on a successful one.

\subparagraph{4-Green}
The target takes 10d6 poison damage on a failed save, or half as much damage on a successful one.

\subparagraph{5-Blue}
The target takes 10d6 cold damage on a failed save, or half as much damage on a successful one.

\subparagraph{6-Indigo}
On a failed save, the target is restrained. It must then make a Constitution saving throw at the end of each of its turns. If it successfully saves three times, the spell ends. If it fails its save three times, it permanently turns to stone and is subjected to the petrified condition. The successes and failures don't need to be consecutive, keep track of both until the target collects three of a kind.

\subparagraph{7-Violet}
On a failed save, the target is blinded. It must then make a Wisdom saving throw at the start of your next turn. A successful save ends the blindness. If it fails that save, the creature is transported to another plane of existence of the DM's choosing and is no longer blinded. (Typically, a creature that is on a plane that isn't its home plane is banished home, while other creatures are usually cast into the Astral or Ethereal planes.)

\subparagraph{8-Special}
The target is struck by two rays. Roll twice more, rerolling any 8.

\addcontentsline{toc}{section}{Ray of Enfeeblement}


\DndSpellHeader%
{Ray of Enfeeblement}
{2nd level Necromancy}
{1 action}
{60 feet}
{V, S}
{Concentration, up to 1 minute}
A black beam of enervating energy springs from your finger toward a creature within range. Make a ranged spell attack against the target. On a hit, the target deals only half damage with weapon attacks that use Strength until the spell ends.

At the end of each of the target's turns, it can make a Constitution saving throw against the spell. On a success, the spell ends.

\addcontentsline{toc}{section}{Ray of Frost}


\DndSpellHeader%
{Ray of Frost}
{Evocation cantrip}
{1 action}
{60 feet}
{V, S}
{Instantaneous}
A frigid beam of blue-white light streaks toward a creature within range. Make a ranged spell attack against the target. On a hit, it takes 1d8 cold damage, and its speed is reduced by 10 feet until the start of your next turn.

The spell's damage increases by 1d8 when you reach 5th level (2d8), 11th level (3d8), and 17th level (4d8).

\addcontentsline{toc}{section}{Ray of Sickness}


\DndSpellHeader%
{Ray of Sickness}
{1st level Necromancy}
{1 action}
{60 feet}
{V, S}
{Instantaneous}
A ray of sickening greenish energy lashes out toward a creature within range. Make a ranged spell attack against the target. On a hit, the target takes 2d8 poison damage and must make a Constitution saving throw. On a failed save, it is also poisoned until the end of your next turn.

\addcontentsline{toc}{section}{Remove Curse}


\DndSpellHeader%
{Remove Curse}
{3rd level Abjuration}
{1 action}
{Touch}
{V, S}
{Instantaneous}
At your touch, all curses affecting one creature or object end. If the object is a cursed magic item, its curse remains, but the spell breaks its owner's attunement to the object so it can be removed or discarded.

\addcontentsline{toc}{section}{Scorching Ray}


\DndSpellHeader%
{Scorching Ray}
{2nd level Evocation}
{1 action}
{120 feet}
{V, S}
{Instantaneous}
You create three rays of fire and hurl them at targets within range. You can hurl them at one target or several.

Make a ranged spell attack for each ray. On a hit, the target takes 2d6 fire damage.

\addcontentsline{toc}{section}{Scrying}


\DndSpellHeader%
{Scrying}
{5th level Divination}
{10 minute}
{Self}
{V, S, a focus worth at least 1,000 gp, such as a crystal ball, a silver mirror, or a font filled with holy water}
{Concentration, up to 10 minutes}
You can see and hear a particular creature you choose that is on the same plane of existence as you. The target must make a Wisdom saving throw, which is modified by how well you know the target and the sort of physical connection you have to it. If a target knows you're casting this spell, it can fail the saving throw voluntarily if it wants to be observed.

\begin{DndTable}[header=Knowledge of Target]{Xc}

Knowledge & Save Modifier\\
Secondhand (you have heard of the target) & +5\\
Firsthand (you have met the target) & +0\\
Familiar (you know the target well) & -5\\
\end{DndTable}

\begin{DndTable}[header=Connection to Target]{Xc}

Connection & Save Modifier\\
Likeness or picture & -2\\
Possession or garment & -4\\
Body part, lock of hair, bit of nail, or the like & -10\\
\end{DndTable}

On a successful save, the target isn't affected, and you can't use this spell against it again for 24 hours.

On a failed save, the spell creates an invisible sensor within 10 feet of the target. You can see and hear through the sensor as if you were there. The sensor moves with the target, remaining within 10 feet of it for the duration. A creature that can see invisible objects sees the sensor as a luminous orb about the size of your fist.

Instead of targeting a creature, you can choose a location you have seen before as the target of this spell. When you do, the sensor appears at that location and doesn't move.

\addcontentsline{toc}{section}{Shatter}


\DndSpellHeader%
{Shatter}
{2nd level Evocation}
{1 action}
{60 feet}
{V, S, a chip of mica}
{Instantaneous}
A sudden loud ringing noise, painfully intense, erupts from a point of your choice within range. Each creature in a 10-foot-radius sphere centered on that point must make a Constitution saving throw. A creature takes 3d8 thunder damage on a failed save, or half as much damage on a successful one. A creature made of inorganic material such as stone, crystal, or metal has disadvantage on this saving throw.

A nonmagical object that isn't being worn or carried also takes the damage if it's in the spell's area.

\addcontentsline{toc}{section}{Shield}


\DndSpellHeader%
{Shield}
{1st level Abjuration}
{1 reaction}
{Self}
{V, S}
{1 round}
An invisible barrier of magical force appears and protects you. Until the start of your next turn, you have a +5 bonus to AC, including against the triggering attack, and you take no damage from \textit{magic missile}.

\addcontentsline{toc}{section}{Shocking Grasp}


\DndSpellHeader%
{Shocking Grasp}
{Evocation cantrip}
{1 action}
{Touch}
{V, S}
{Instantaneous}
Lightning springs from your hand to deliver a shock to a creature you try to touch. Make a melee spell attack against the target. You have advantage on the attack roll if the target is wearing armor made of metal. On a hit, the target takes 1d8 lightning damage, and it can't take reactions until the start of its next turn.

The spell's damage increases by 1d8 when you reach 5th level (2d8), 11th level (3d8), and 17th level (4d8).

\addcontentsline{toc}{section}{Silent Image}


\DndSpellHeader%
{Silent Image}
{1st level Illusion}
{1 action}
{60 feet}
{V, S, a bit of fleece}
{Concentration, up to 10 minutes}
You create the image of an object, a creature, or some other visible phenomenon that is no larger than a 15-foot cube. The image appears at a spot within range and lasts for the duration. The image is purely visual; it isn't accompanied by sound, smell, or other sensory effects.

You can use your action to cause the image to move to any spot within range. As the image changes location, you can alter its appearance so that its movements appear natural for the image. For example, if you create an image of a creature and move it, you can alter the image so that it appears to be walking.

Physical interaction with the image reveals it to be an illusion, because things can pass through it. A creature that uses its action to examine the image can determine that it is an illusion with a successful Intelligence (Investigation) check against your spell save DC. If a creature discerns the illusion for what it is, the creature can see through the image.

\addcontentsline{toc}{section}{Sleep}


\DndSpellHeader%
{Sleep}
{1st level Enchantment}
{1 action}
{90 feet}
{V, S, a pinch of fine sand, rose petals, or a cricket}
{1 minute}
This spell sends creatures into a magical slumber. Roll 5d8; the total is how many hit points of creatures this spell can affect. Creatures within 20 feet of a point you choose within range are affected in ascending order of their current hit points (ignoring unconscious creatures).

Starting with the creature that has the lowest current hit points, each creature affected by this spell falls unconscious until the spell ends, the sleeper takes damage, or someone uses an action to shake or slap the sleeper awake. Subtract each creature's hit points from the total before moving on to the creature with the next lowest hit points. A creature's hit points must be equal to or less than the remaining total for that creature to be affected.

Undead and creatures immune to being charmed aren't affected by this spell.

\addcontentsline{toc}{section}{Slow}


\DndSpellHeader%
{Slow}
{3rd level Transmutation}
{1 action}
{120 feet}
{V, S, a drop of molasses}
{Concentration, up to 1 minute}
You alter time around up to six creatures of your choice in a 40-foot cube within range. Each target must succeed on a Wisdom saving throw or be affected by this spell for the duration.

An affected target's speed is halved, it takes a −2 penalty to AC and Dexterity saving throws, and it can't use reactions. On its turn, it can use either an action or a bonus action, not both. Regardless of the creature's abilities or magic items, it can't make more than one melee or ranged attack during its turn.

If the creature attempts to cast a spell with a casting time of 1 action, roll a d20. On an 11 or higher, the spell doesn't take effect until the creature's next turn, and the creature must use its action on that turn to complete the spell. If it can't, the spell is wasted.

A creature affected by this spell makes another Wisdom saving throw at the end of each of its turns. On a successful save, the effect ends for it.

\addcontentsline{toc}{section}{Speak with Animals}


\DndSpellHeader%
{Speak with Animals}
{1st level Divination}
{1 action}
{Self}
{V, S}
{10 minutes}
You gain the ability to comprehend and verbally communicate with beasts for the duration. The knowledge and awareness of many beasts is limited by their intelligence, but at minimum, beasts can give you information about nearby locations and monsters, including whatever they can perceive or have perceived within the past day. You might be able to persuade a beast to perform a small favor for you, at the DM's discretion.

\addcontentsline{toc}{section}{Stinking Cloud}


\DndSpellHeader%
{Stinking Cloud}
{3rd level Conjuration}
{1 action}
{90 feet}
{V, S, a rotten egg or several skunk cabbage leaves}
{Concentration, up to 1 minute}
You create a 20-foot-radius sphere of yellow, nauseating gas centered on a point within range. The cloud spreads around corners, and its area is Vision and Light. The cloud lingers in the air for the duration.

Each creature that is completely within the cloud at the start of its turn must make a Constitution saving throw against poison. On a failed save, the creature spends its action that turn retching and reeling. Creatures that don't need to breathe or are immune to poison automatically succeed on this saving throw.

A moderate wind (at least 10 miles per hour) disperses the cloud after 4 rounds. A strong wind (at least 20 miles per hour) disperses it after 1 round.

\addcontentsline{toc}{section}{Stone Shape}


\DndSpellHeader%
{Stone Shape}
{4th level Transmutation}
{1 action}
{Touch}
{V, S, soft clay, which must be worked into roughly the desired shape of the stone object}
{Instantaneous}
You touch a stone object of Medium size or smaller or a section of stone no more than 5 feet in any dimension and form it into any shape that suits your purpose. So, for example, you could shape a large rock into a weapon, idol, or coffer, or make a small passage through a wall, as long as the wall is less than 5 feet thick. You could also shape a stone door or its frame to seal the door shut. The object you create can have up to two hinges and a latch, but finer mechanical detail isn't possible.

\addcontentsline{toc}{section}{Stoneskin}


\DndSpellHeader%
{Stoneskin}
{4th level Abjuration}
{1 action}
{Touch}
{V, S, diamond dust worth 100 gp, which the spell consumes}
{Concentration, up to 1 hour}
This spell turns the flesh of a willing creature you touch as hard as stone. Until the spell ends, the target has resistance to nonmagical bludgeoning, piercing, and slashing damage.

\addcontentsline{toc}{section}{Suggestion}


\DndSpellHeader%
{Suggestion}
{2nd level Enchantment}
{1 action}
{30 feet}
{V, a snake's tongue and either a bit of honeycomb or a drop of sweet oil}
{Concentration, up to 8 hours}
You suggest a course of activity (limited to a sentence or two) and magically influence a creature you can see within range that can hear and understand you. Creatures that can't be charmed are immune to this effect. The suggestion must be worded in such a manner as to make the course of action sound reasonable. Asking the creature to stab itself, throw itself onto a spear, immolate itself, or do some other obviously harmful act ends the spell.

The target must make a Wisdom saving throw. On a failed save, it pursues the course of action you described to the best of its ability. The suggested course of action can continue for the entire duration. If the suggested activity can be completed in a shorter time, the spell ends when the subject finishes what it was asked to do.

You can also specify conditions that will trigger a special activity during the duration. For example, you might suggest that a knight give her warhorse to the first beggar she meets. If the condition isn't met before the spell expires, the activity isn't performed.

If you or any of your companions damage the target, the spell ends.

\addcontentsline{toc}{section}{Telekinesis}


\DndSpellHeader%
{Telekinesis}
{5th level Transmutation}
{1 action}
{60 feet}
{V, S}
{Concentration, up to 10 minutes}
You gain the ability to move or manipulate creatures or objects by thought. When you cast the spell, and as your action each round for the duration, you can exert your will on one creature or object that you can see within range, causing the appropriate effect below. You can affect the same target round after round, or choose a new one at any time. If you switch targets, the prior target is no longer affected by the spell.

\subparagraph{Creature}
You can try to move a Huge or smaller creature. Make an ability check with your spellcasting ability contested by the creature's Strength check. If you win the contest, you move the creature up to 30 feet in any direction, including upward but not beyond the range of this spell. Until the end of your next turn, the creature is restrained in your telekinetic grip. A creature lifted upward is suspended in mid-air.

On subsequent rounds, you can use your action to attempt to maintain your telekinetic grip on the creature by repeating the contest.

\subparagraph{Object}
You can try to move an object that weighs up to 1,000 pounds. If the object isn't being worn or carried, you automatically move it up to 30 feet in any direction, but not beyond the range of this spell.

If the object is worn or carried by a creature, you must make an ability check with your spellcasting ability contested by that creature's Strength check. If you succeed, you pull the object away from that creature and can move it up to 30 feet in any direction but not beyond the range of this spell.

You can exert fine control on objects with your telekinetic grip, such as manipulating a simple tool, opening a door or a container, stowing or retrieving an item from an open container, or pouring the contents from a vial.

\addcontentsline{toc}{section}{Thunderwave}


\DndSpellHeader%
{Thunderwave}
{1st level Evocation}
{1 action}
{Self (15 feet cube)}
{V, S}
{Instantaneous}
A wave of thunderous force sweeps out from you. Each creature in a 15-foot cube originating from you must make a Constitution saving throw. On a failed save, a creature takes 2d8 thunder damage and is pushed 10 feet away from you. On a successful save, the creature takes half as much damage and isn't pushed.

In addition, unsecured objects that are completely within the area of effect are automatically pushed 10 feet away from you by the spell's effect, and the spell emits a thunderous boom audible out to 300 feet.

\addcontentsline{toc}{section}{Tongues}


\DndSpellHeader%
{Tongues}
{3rd level Divination}
{1 action}
{Touch}
{V, a small clay model of a ziggurat}
{1 hour}
This spell grants the creature you touch the ability to understand any spoken language it hears. Moreover, when the target speaks, any creature that knows at least one language and can hear the target understands what it says.

\addcontentsline{toc}{section}{True Seeing}


\DndSpellHeader%
{True Seeing}
{6th level Divination}
{1 action}
{Touch}
{V, S, an ointment for the eyes that costs 25 gp; is made from mushroom powder, saffron, and fat; and is consumed by the spell}
{1 hour}
This spell gives the willing creature you touch the ability to see things as they actually are. For the duration, the creature has truesight, notices secret doors hidden by magic, and can see into the Ethereal Plane, all out to a range of 120 feet.

\addcontentsline{toc}{section}{Unseen Servant}


\DndSpellHeader%
{Unseen Servant}
{1st level Conjuration}
{1 action}
{60 feet}
{V, S, a piece of string and a bit of wood}
{1 hour}
This spell creates an invisible, mindless, shapeless, Medium force that performs simple tasks at your command until the spell ends. The servant springs into existence in an unoccupied space on the ground within range. It has AC 10, 1 hit point, and a Strength of 2, and it can't attack. If it drops to 0 hit points, the spell ends.

Once on each of your turns as a bonus action, you can mentally command the servant to move up to 15 feet and interact with an object. The servant can perform simple tasks that a human servant could do, such as fetching things, cleaning, mending, folding clothes, lighting fires, serving food, and pouring wine. Once you give the command, the servant performs the task to the best of its ability until it completes the task, then waits for your next command.

If you command the servant to perform a task that would move it more than 60 feet away from you, the spell ends.

\addcontentsline{toc}{section}{Vampiric Touch}


\DndSpellHeader%
{Vampiric Touch}
{3rd level Necromancy}
{1 action}
{Self}
{V, S}
{Concentration, up to 1 minute}
The touch of your shadow-wreathed hand can siphon life force from others to heal your wounds. Make a melee spell attack against a creature within your reach. On a hit, the target takes 3d6 necrotic damage, and you regain hit points equal to half the amount of necrotic damage dealt. Until the spell ends, you can make the attack again on each of your turns as an action.

\addcontentsline{toc}{section}{Wall of Fire}


\DndSpellHeader%
{Wall of Fire}
{4th level Evocation}
{1 action}
{120 feet}
{V, S, a small piece of phosphorus}
{Concentration, up to 1 minute}
You create a wall of fire on a solid surface within range. You can make the wall up to 60 feet long, 20 feet high, and 1 foot thick, or a ringed wall up to 20 feet in diameter, 20 feet high, and 1 foot thick. The wall is opaque and lasts for the duration.

When the wall appears, each creature within its area must make a Dexterity saving throw. On a failed save, a creature takes 5d8 fire damage, or half as much damage on a successful save.

One side of the wall, selected by you when you cast this spell, deals 5d8 fire damage to each creature that ends its turn within 10 feet of that side or inside the wall. A creature takes the same damage when it enters the wall for the first time on a turn or ends its turn there. The other side of the wall deals no damage.

\addcontentsline{toc}{section}{Wall of Force}


\DndSpellHeader%
{Wall of Force}
{5th level Evocation}
{1 action}
{120 feet}
{V, S, a pinch of powder made by crushing a clear gemstone}
{Concentration, up to 10 minutes}
An invisible wall of force springs into existence at a point you choose within range. The wall appears in any orientation you choose, as a horizontal or vertical barrier or at an angle. It can be free floating or resting on a solid surface. You can form it into a hemispherical dome or a sphere with a radius of up to 10 feet, or you can shape a flat surface made up of ten 10-foot-by-10-foot panels. Each panel must be contiguous with another panel. In any form, the wall is 1/4 inch thick. It lasts for the duration. If the wall cuts through a creature's space when it appears, the creature is pushed to one side of the wall (your choice which side).

Nothing can physically pass through the wall. It is immune to all damage and can't be dispelled by \textit{dispel magic}. A \textit{disintegrate} spell destroys the wall instantly, however. The wall also extends into the Ethereal Plane, blocking ethereal travel through the wall.

\addcontentsline{toc}{section}{Wall of Ice}


\DndSpellHeader%
{Wall of Ice}
{6th level Evocation}
{1 action}
{120 feet}
{V, S, a small piece of quartz}
{Concentration, up to 10 minutes}
You create a wall of ice on a solid surface within range. You can form it into a hemispherical dome or a sphere with a radius of up to 10 feet, or you can shape a flat surface made up of ten 10-foot-square panels. Each panel must be contiguous with another panel. In any form, the wall is 1 foot thick and lasts for the duration.

If the wall cuts through a creature's space when it appears, the creature within its area is pushed to one side of the wall and must make a Dexterity saving throw. On a failed save, the creature takes 10d6 cold damage, or half as much damage on a successful save.

The wall is an object that can be damaged and thus breached. It has AC 12 and 30 hit points per 10-foot section, and it is vulnerable to fire damage. Reducing a 10-foot section of wall to 0 hit points destroys it and leaves behind a sheet of frigid air in the space the wall occupied. A creature moving through the sheet of frigid air for the first time on a turn must make a Constitution saving throw. That creature takes 5d6 cold damage on a failed save, or half as much damage on a successful one.

\addcontentsline{toc}{section}{Web}


\DndSpellHeader%
{Web}
{2nd level Conjuration}
{1 action}
{60 feet}
{V, S, a bit of spiderweb}
{Concentration, up to 1 hour}
You conjure a mass of thick, sticky webbing at a point of your choice within range. The webs fill a 20-foot cube from that point for the duration. The webs are difficult terrain and lightly obscure their area.

If the webs aren't anchored between two solid masses (such as walls or trees) or layered across a floor, wall, or ceiling, the conjured web collapses on itself, and the spell ends at the start of your next turn. Webs layered over a flat surface have a depth of 5 feet.

Each creature that starts its turn in the webs or that enters them during its turn must make a Dexterity saving throw. On a failed save, the creature is restrained as long as it remains in the webs or until it breaks free.

A creature restrained by the webs can use its action to make a Strength check against your spell save DC. If it succeeds, it is no longer restrained.

The webs are flammable. Any 5-foot cube of webs exposed to fire burns away in 1 round, dealing 2d4 fire damage to any creature that starts its turn in the fire.

\addcontentsline{toc}{section}{Wish}


\DndSpellHeader%
{Wish}
{9th level Conjuration}
{1 action}
{Self}
{V}
{Instantaneous}
Wish is the mightiest spell a mortal creature can cast. By simply speaking aloud, you can alter the very foundations of reality in accord with your desires.

The basic use of this spell is to duplicate any other spell of 8th level or lower. You don't need to meet any requirements in that spell, including costly components. The spell simply takes effect. Alternatively, you can create one of the following effects of your choice:


\begin{itemize}
\item You create one object of up to 25,000 gp in value that isn't a magic item. The object can be no more than 300 feet in any dimension, and it appears in an unoccupied space you can see on the ground.
\item You allow up to twenty creatures that you can see to regain all hit points, and you end all effects on them described in the \textit{greater restoration} spell.
\item You grant up to ten creatures that you can see resistance to a damage type you choose.
\item You grant up to ten creatures you can see immunity to a single spell or other magical effect for 8 hours. For instance, you could make yourself and all your companions immune to a \textbf{lich}'s life drain attack.
\item You undo a single recent event by forcing a reroll of any roll made within the last round (including your last turn). Reality reshapes itself to accommodate the new result. For example, a \textit{wish} spell could undo an opponent's successful save, a foe's critical hit, or a friend's failed save. You can force the reroll to be made with advantage or disadvantage, and you can choose whether to use the reroll or the original roll.
\end{itemize}

You might be able to achieve something beyond the scope of the above examples. State your wish to the DM as precisely as possible. The DM has great latitude in ruling what occurs in such an instance; the greater the wish, the greater the likelihood that something goes wrong. This spell might simply fail, the effect you desire might only be partly achieved, or you might suffer some unforeseen consequence as a result of how you worded the wish. For example, wishing that a villain were dead might propel you forward in time to a period when that villain is no longer alive, effectively removing you from the game. Similarly, wishing for a legendary magic item or artifact might instantly transport you to the presence of the item's current owner.

The stress of casting this spell to produce any effect other than duplicating another spell weakens you. After enduring that stress, each time you cast a spell until you finish a long rest, you take 1d10 necrotic damage per level of that spell. This damage can't be reduced or prevented in any way. In addition, your Strength drops to 3, if it isn't 3 or lower already, for 2d4 days. For each of those days that you spend resting and doing nothing more than light activity, your remaining recovery time decreases by 2 days. Finally, there is a Able to cast again chance that you are unable to cast wish ever again if you suffer this stress.

\end{document}
